\newcommand{\DoNotLoadEpstopdf}{}
\documentclass[11pt,reqno]{amsart}
\usepackage[T1]{fontenc}
\usepackage{lmodern,amsmath,amssymb,amsthm,mathtools,microtype,array,mathrsfs}
\usepackage[margin=1.1in]{geometry}
\usepackage[hidelinks,hypertexnames=false]{hyperref}
\hypersetup{pdftitle={Collision records, arithmetic, and raw local inversion in a triangular Lorentz gas},pdfauthor={Qian Qi}}
\setlength{\emergencystretch}{2em}
\allowdisplaybreaks[2]
\numberwithin{equation}{section}
\newtheorem{theorem}{Theorem}[section]
\newtheorem{leadtheorem}{Theorem}
\newtheorem{maintheorem}{Theorem}
\newtheorem{revisiontheorem}{Theorem}
\renewcommand{\therevisiontheorem}{\Alph{revisiontheorem}}
\renewcommand{\themaintheorem}{\Alph{maintheorem}}
\newtheorem{lemma}[theorem]{Lemma}
\newtheorem{proposition}[theorem]{Proposition}
\newtheorem{corollary}[theorem]{Corollary}
\theoremstyle{definition}\newtheorem{definition}[theorem]{Definition}
\theoremstyle{remark}\newtheorem{remark}[theorem]{Remark}
\newcommand{\R}{\mathbb R}\newcommand{\Z}{\mathbb Z}\newcommand{\T}{\mathbb T}
\newcommand{\dd}{\,\mathrm d}\newcommand{\one}{\mathbf1}
\newcommand{\supp}{\operatorname{supp}}\newcommand{\TV}{\operatorname{TV}}
\title[Technical supplement]{Technical supplement to Collision records, arithmetic, and raw local inversion in a triangular Lorentz gas}
\author{Qian Qi}
\date{October 12, 2026. A2-DYN, revision 105; complete retained source through revision 104}
\subjclass[2020]{37D50, 37A50, 60F05, 42A38}
\keywords{Sinai billiards, collision records, local limits, inducing, Fourier inversion, geometric response}

\theoremstyle{definition}\newtheorem{inputstatement}[theorem]{Input statement}
\newcommand{\esssup}{\operatorname*{ess\,sup}}
\providecommand{\diver}{\operatorname{div}}
\usepackage{xstring}
\begin{document}
\raggedbottom
\begin{abstract}
We study exact displacement, collision count and flight time for actual
returns in a triangular periodic Lorentz gas. We prove that the complete
physical boundary source has a finite distributional derivative at each
collision count, uniformly in the radius, the original exact labels and
the guard width at that count. The proof writes the unchanged smooth
guard as a signed threshold integral and combines definable finite-word
geometry with the complete density bound. A two-sided averaging inequality
then bounds the unrecovered height by paired incoming and outgoing flux,
without charging an isolated jump. Finally the ordered scalar pointwise
error equals the bounded-Lipschitz path-numerator error on the same roofs.
The collision-uniform decay of the paired flux remains the quantitative
estimate needed for the unrestricted arithmetic pointwise law.
\end{abstract}
\maketitle

\paragraph{The physical record and the pointwise problem.}
Let $R\in I=[9/20,47/100]$ and let
$J_{n,R}=(K_{n,R},N_{n,R},T_{n,R})$ be the actual return record
under $\nu_R^*=c_*^{-1}\nu|_{Y_R^*}$, where $c_*=91/(10000\pi)$.
The occupation convention counts collisions $0,\ldots,m-1$ and tests
terminal membership at collision $m$. Write
\[
 Z_{m,R}(k,t,n)=\frac{(k_1,k_2,t-m\bar\tau_R,n-mc_*)}{\sqrt m}.
\]
The original pointwise objective is, for every $M<\infty$,
\begin{equation}\label{eq:intro-v73-target}
 \sup_{R,\,1\le n\le m,\,k\in\Z^2}
 \operatorname*{ess\,sup}_{t:\,|Z_{m,R}(k,t,n)|\le M}
 |m^2p_{n,R}(k,m,t)-\mathcal L_{m,R}(k_1,k_2,t,n)|\longrightarrow0.
\end{equation}
Here $\mathcal L_{m,R}$ retains the full finite arithmetic transition
kernel, including zero classes. The new results concern this exact
record; they do not replace it by a window or a pooled discrete label.
Let $\mathcal E_M$ be the collision limsup of the left side of
\eqref{eq:intro-v73-target}. For a family of densities set
\begin{equation}\label{eq:intro-v73-ordered-height}
 \mathcal H(f)=\limsup_{m\to\infty}\sup_{R,n,k}m^2
                         \|f_{n,k,m,R}\|_{L^\infty(\R)}.
\end{equation}
Every density norm is an essential norm. Widths and reserves are fixed
before this collision limit.

\paragraph{The complete source.}
Let $b=b^\varepsilon_{n,k,m,R}$ be the density of
$(L_{m,R})_\#((1-H_{m,R}^\varepsilon)\one_{\Pi_{n,k,m,R}}\nu_R^*)$.
Its original positive decomposition is
\[
 b=s_{71,1}+d_{71,1}+b^{\varepsilon,\mathrm{inc}}+e_{70}.
\]
The term $e_{70}$ contains first clearance outside the selected disks;
first incidence is separate. We keep all these states. The canonical
current is $Db$, constructed directly from this full source. Our first
result strengthens its topology without changing the source.

\begin{revisiontheorem}[Finite variation of the complete current]
\label{thm:intro-v73-bv}
At each fixed collision count $m\ge2$, the full density $b$ belongs to
$BV(\R)$ and its derivative has total mass zero. There is a finite
constant $C_m$, uniform in $R,\varepsilon,n,k$ at that count, such
that $\|Db\|_{\TV}\le C_m$. A possible bound is
\[
 C_m=4(J_m+1)(CA^m/c_*)\|d\Theta\|_{\TV}^{2m+1},
\]
where $J_m$ is the uniform monotonicity number of the finite threshold
family. No collision-uniform bound on $m^2C_m$ is part of this assertion.
\end{revisiontheorem}
\begin{proof}
Lemma~\ref{lem:v73-threshold-mixture} represents the original smooth
guard exactly. Lemma~\ref{lem:v73-threshold-monotonicity} combines
finite physical geometry, parameterized integration and definable
monotonicity. Theorem~\ref{thm:v73-full-source-bv} integrates the
resulting variation bound to obtain the original full current.
\end{proof}

\begin{revisiontheorem}[Paired flux and raw height]
\label{thm:intro-v73-paired}
For a nonnegative compactly supported $BV$ density,
\[
 [b-2A_\delta b]_+(t)
 \le\min\{(Db)^+([t-\delta,t]),(Db)^-([t,t+\delta])\}
 \quad\hbox{for a.e. }t.
\]
Let $\mathcal P_\varepsilon(\delta)$ be the ordered normalized
essential supremum of the right side, defined in
\eqref{eq:v73-paired-height}. For every sufficiently large fixed $B$,
with $\varepsilon(B)=A_0B^{-1/12}$ and any fixed $\delta(B)>0$,
\begin{equation}\label{eq:intro-v73-budget}
 \mathcal E_M\le C_MB^{-1/192}
                  +\mathcal P_{\varepsilon(B)}(\delta(B)).
\end{equation}
The choice of $B$ and its width precedes $m\to\infty$; only then
may $B\to\infty$. The reserve is the constant two.
\end{revisiontheorem}
\begin{proof}
Lemma~\ref{lem:v73-paired-flux} compares both one-sided averages.
Theorem~\ref{thm:v73-paired-budget} applies the resulting bound to
the original source and the fixed-width local-variation estimate.
The source-preserving allocation with reserve two still retains
$s_{71,1}$; no original residual is discarded.
\end{proof}

\begin{revisiontheorem}[Equality of scalar and path errors]
\label{thm:intro-v73-path}
On the inherited common positive path-remainder theorem, the ordered
bounded-Lipschitz path-numerator error equals $\mathcal E_M$ on every
original central target. Thus the scalar pointwise endpoint and the
same-roof collision-path numerator have the same vanishing condition.
Positive-reference normalization gives the conditional collision bridge
and forward essential arithmetic likelihood when that endpoint holds.
\end{revisiontheorem}
\begin{proof}
Theorem~\ref{thm:v73-scalar-path-equality} proves the equality by
positivity of a common path remainder. Corollary~\ref{cor:v73-central-excess}
gives the normalization and the exact central excess comparison.
Zero arithmetic classes are never normalized.
\end{proof}

The missing quantitative assertion is decay of
$\mathcal P_{\varepsilon(B)}(\delta(B))$ for a legal choice of widths,
or directly of the equivalent central excess. Finite variation at each
count does not imply this assertion. In particular the present results
do not yet establish \eqref{eq:intro-v73-target}. They prove the full
finite-measure property, remove the charge for isolated physical jumps,
and identify the scalar estimate that also controls the path numerator.
The complete inherited proof chain follows the new argument. Its
preceding opening is retained in Appendix~\ref{app:v73-retained-front}.

\input{appendices/v85_dependency_guide}
\input{appendices/v86_dependency_guide}
\input{appendices/v87_dependency_guide}
\input{appendices/v88_dependency_guide}
\input{appendices/v89_dependency_guide}
\part{Positive angular intervals and paid critical arms}
\begingroup
\let\ArchiveLabel\label
\let\ArchiveRef\ref
\renewcommand{\label}[1]{\ArchiveLabel{v104archive:#1}}
\renewcommand{\ref}[1]{%
  \IfBeginWith{#1}{supp-}{%
    \StrBehind{#1}{supp-}[\ArchiveTarget]%
    \ArchiveRef{\ArchiveTarget}%
  }{\ArchiveRef{v104archive:#1}}%
}
\makeatletter
\let\ltx@label\label
\makeatother
\paragraph{Retained revision-104 abstract.}

We study the exact-label pointwise raw local limit problem for actual
returns in a triangular Lorentz gas. The remaining upper source is
the positive minimum of two weighted endpoint traces on short regular
segments. We prove an angular-interval comparison for this same
positive measure. The comparison uses endpoint order, integrated
trace balance and physical overlap, but no reciprocal endpoint
velocity. On regular backward-critical arms it gives an essential-height
bound through stationary outgoing roofs, uniformly in the relative
position inside the segment. An integral comparison first isolates
any shortage of outgoing trace. On these critical arms its entire
shortage is then paid by the actual uncut optical volume source,
without an outgoing trace-fraction lower bound. The recovered source has height $O(B^{-1/192})$
for uniformly balanced arms, or $O(B^{-1/384})$ under a growing balance
budget. An exact positive decomposition retains uncovered physical pairs;
no complete coverage theorem is inferred. The original pointwise assertion remains
unchanged; decay of that complete deficit is not asserted.


\input{journal-problem-v93}

\subsection{Comparison at a stationary endpoint}
The previous endpoint-distance flows compare nearby trajectories of
the same exact event. Their paired coarea densities are
$c^\varepsilon/|a'|$ and $c^\varepsilon/|b'|$, where
$c^\varepsilon$ is the positive minimum of the weighted traces at
the two ends of a maximal regular segment. These formulas remove
the relative-position degeneracy but still require nonzero endpoint
speed along a useful comparison path.

The argument below uses intervals of actual angular trajectories
instead of a single comparison path. If all trajectories in a
comparison interval contain a one-sided roof translate of every
source trajectory, the integral of their positive trace controls the
source at its original roof. The ratio of the two integrated traces
and the physical overlap are the only numerical costs. Stationary
endpoint velocity and rapid local variation of the minimum do not
enter that ratio. Theorem~\ref{thm:v104-stationary} verifies it on
critical arms with an explicit level-width condition and the original uncut optical variation. Theorem~
ef{thm:v104-paid-arms}
also pays the entire comparison-trace shortage. Its whole-family
height bound includes all
relative positions, without a lower original residence cutoff.

The pointwise problem is still \eqref{eq:v76-target}. The new
theorems estimate specified positive source families at that topology,
not a replacement finite-norm problem. Theorem~\ref{thm:v104-recovery}
retains a nonnegative remainder with a precise comparison-trace
deficit. Its vanishing on the complete physical family remains the
last requirement in the resulting scalar reduction.

\section{The inherited physical source in the required form}
Fix the radius, collision count and exact displacement and occupation
labels. On the disjoint original space of regular pairing sheets,
write $a(\theta)<u<b(\theta)$ for a maximal roof segment. On the regular set away from $g=0$, put
\[
 q=(\alpha_0,\alpha_m),\quad F=L,\quad
 g=\nabla_qF=R(-p_0,p_m),\quad z=|g|,\quad
 I_m=1+\sum_{i=1}^{m-1}q_0^{d_i}/c_i.
\]
Here $\omega=|\partial_{p_0}\alpha_m|^{-1}$ and
$\chi=|\det D_q(F,\arg g)|$. On the regular set away from $g=0$, put
\[
 q=(\alpha_0,\alpha_m),\quad F=L,\quad
 g=\nabla_qF=R(-p_0,p_m),\quad z=|g|,\quad
 I_m=1+\sum_{i=1}^{m-1}q_0^{d_i}/c_i.
\]
Here $\omega=|\partial_{p_0}\alpha_m|^{-1}$ and
$\chi=|\det D_q(F,\arg g)|$. The physical angular density is
\[
 A(u,\theta)=\frac{\omega(u,\theta)}{4\pi c_*\chi(u,\theta)},
 \qquad f^\varepsilon=A W^\varepsilon,
 \qquad c^\varepsilon(\theta)
     =\min\{f^\varepsilon(a+,\theta),f^\varepsilon(b-,\theta)\}.
\]
The initial section visit is counted and the terminal section visit
is tested separately, as in \eqref{eq:v76-event}. Every angular
sheet carries this same exact event; its incident physical component
is part of its identity.

\begin{inputstatement}[Positive pairing and local concentration]
\label{input:v104-pair}
The original paired measure is
$c^\varepsilon(\theta)\,d\theta\,du$ on $a<u<b$.
For $0<\Delta$ let $\mathsf T_\Delta^\varepsilon$ be its roof
density restricted to $0<\ell=b-a<\Delta$, without an additional
selector. At every previously fixed $r,\Delta,\varepsilon>0$,
\begin{equation}\label{eq:v104-paired-average}
 \limsup_m\sup_{R,n,k}m^2
       \|A_r\mathsf T_\Delta^\varepsilon\|_\infty
 \le C\{\varepsilon^{1/16}+(1+\varepsilon^{-1})\Delta\},
 \qquad A_rf(t)=\frac1{2r}\int_{t-r}^{t+r}f(u)\,du.
\end{equation}
The leading constant is independent of the previously fixed widths.
For fixed $m,R,n,k$, all nonnegative angular restrictions use the
same coarea kernels outside a roof-null set. Countable selections
use common versions. The uniform order is $\sup_R\esssup_t$;
one universal null set for all radii is not asserted.
\end{inputstatement}
This is the positive barycentric source and complete paired-average
estimate of Supplement Theorem~\ref{supp-v103archive:thm:v102-pair-coarea}
and equation~\eqref{supp-v103archive:eq:v102-paired-average}.
The paired-average estimate follows there from the original weighted
source and the actual uncut inverse-gradient volume source. It is
not an assumed essential-height estimate for the paired source.
The same inherited analytic input gives the actual uncut source
\begin{equation}\label{eq:v104-uncut-input}
 v^0(t)\,dt=F_\#(I_m z^{-1}\one_\Pi\nu_R^*),\qquad
 \limsup_m\sup_{R,n,k}m^2\|A_\rho v^0\|_\infty\le C
\end{equation}
for every previously fixed $\rho>0$, with a width-independent leading
constant. It is finite at fixed count. On each regular segment,
\begin{equation}\label{eq:v104-weighted-variation-input}
 |\partial_u f^\varepsilon(u,\theta)|
 \le C(1+\varepsilon^{-1})A(u,\theta)I_m(u,\theta)/z(u,\theta).
\end{equation}
These are the actual uncut volume and weighted optical variation
statements in Supplement Input statement~\ref{supp-v103archive:input:v102-analytic}
and its physical input. No unweighted pointwise height bound for
$v^0$ is assumed. The bound is used only on a previously fixed
positive window after the physical sum.

The Borel physical sheet summation retains multiplicity one; its
historical proof and both endpoint trace conventions are unchanged.

\begin{inputstatement}[Regular critical entrance and scalar reduction]
\label{input:v104-scalar}
At a regular nongrazing backward-critical extension $q_*$, the
incoming roof is constant and
\begin{equation}\label{eq:v104-critical-input}
 f^\varepsilon(a+,\theta)
 =\frac{N_*}{e_\perp(\theta)^{\mathsf T}K_*e_\perp(\theta)},
 \quad N_*=\frac{\omega(q_*)W^\varepsilon(q_*)}{4\pi c_*},
 \quad \lambda_\circ I\le K_*\le\Lambda_\circ I,
\end{equation}
where $\lambda_\circ=9\sqrt{19}/200$ and $\Lambda_\circ=47/6$.
The incoming trace can be positive. A singular stationary extension
is not included in this formula merely because its incoming roof
is critical.

Under the original outer schedule
\begin{equation}\label{eq:v104-schedule}
 \varepsilon(B)=A_0B^{-1/12},\qquad q_0^*=1/8,\qquad
 r_0=c_0\varepsilon(B)^4,\qquad\delta=6q_0^*r_0,
\end{equation}
the preceding positive remainder $\mathsf U_B^{\rm rem}$ is the
original short paired source on failure of the relative clock and
both endpoint-distance clocks. With
\[
 \mathcal H_M(T)=\limsup_m\sup_{R,n,k}m^2
       \esssup_{t\in\mathcal C_{m,R,n,k}(M)}T(t),
\]
one has
\begin{equation}\label{eq:v104-old-remainder}
 \mathcal E_M\le\mathcal H_M(\mathsf U_B^{\rm rem})
                                       +C_MB^{-1/384}.
\end{equation}
Along this same schedule, the original pointwise assertion is
equivalent to $\mathcal H_M(\mathsf U_B^{\rm rem})\to0$.
\end{inputstatement}
The trace formula is Supplement
Proposition~\ref{supp-v103archive:prop:v103-critical-trace}; the
scalar statement is Supplement Theorem~\ref{supp-v103archive:thm:v103-recovery}.
The earlier relative and endpoint-distance comparisons remain
available and unchanged. The new comparison is applied specifically
to their positive failure source below.

\input{core/250_positive_angular_interval_comparison}
\input{core/251_stationary_critical_arm_recovery}
\input{core/252_positive_deficit_and_original_endpoint}

\appendix
\input{appendices/v104_dependency_map}

\endgroup
\part{Endpoint-distance coarea and positive recovery}
\begingroup
\let\ArchiveLabel\label
\let\ArchiveRef\ref
\renewcommand{\label}[1]{\ArchiveLabel{v103archive:#1}}
\renewcommand{\ref}[1]{%
  \IfBeginWith{#1}{supp-}{%
    \StrBehind{#1}{supp-}[\ArchiveTarget]%
    \ArchiveRef{\ArchiveTarget}%
  }{\ArchiveRef{v103archive:#1}}%
}
\makeatletter
\let\ltx@label\label
\makeatother
\paragraph{Retained revision-103 abstract.}

We study the exact-label pointwise raw local limit problem for actual
returns in a triangular Lorentz gas. The remaining upper source is
the positive minimum of the two weighted endpoint traces on each
short regular segment. For this same measure we construct two
comparisons preserving distance in roof time from one endpoint,
rather than relative position between the endpoints. Their coarea
densities contain neither a residence factor nor a relative-position
factor. An expanding-endpoint criterion therefore admits all relative
positions, including approach to a positive backward-critical trace
and to a fold of the relative-position map. After physical summation,
the certified source has essential height $O(B^{-1/192})$ under the
original outer schedule. The two clocks also recover part of the
previous failure source with height $O(B^{-1/384})$. The residual
which fails all three comparisons remains positive and is not proved
to vanish. The full pointwise target, exact labels, original guard,
arithmetic classes and preceding proofs are retained.


\input{journal-problem-v93}

\subsection{Two additional directions on the positive paired measure}
Write $a(\theta)<u<b(\theta)$ for an original maximal regular
segment, $\ell=b-a$, and
$c^\varepsilon=\min\{f^\varepsilon(a+),f^\varepsilon(b-)\}$.
The earlier transverse comparison fixes $s=(u-a)/\ell$ and has
speed $(1-s)a'+sb'$. At a backward critical entrance with $a'=0$
this speed degenerates as $s\downarrow0$; it may also vanish when
the two endpoint derivatives have opposite signs.

We instead fix either $u-a$ or $b-u$. The exact paired densities
in the two resulting roof coordinates are $c^\varepsilon/|a'|$
and $c^\varepsilon/|b'|$. The absence of $\ell$ and $s$ is an
identity of measures, not an asymptotic cancellation. A one-sided
expansion of $\ell$ keeps every original interior position inside
its segment, even when the starting residence is arbitrarily short.
The comparison still moves between actual trajectories of the same
exact event; it does not continue the radial orbit beyond its first
exit or change any section decision.

Theorem~\ref{thm:v103-gap-height} gives a finite-count pointwise
comparison after physical sheets have been summed. Theorem~\ref{thm:v103-certificate}
and Corollary~\ref{cor:v103-corners} prove its uniformity in relative
position, with the precise reach and trace-variation hypotheses.
Theorem~\ref{thm:v103-recovery} pays the recovered positive source
and retains the exact remainder. In particular a conditional endpoint
certificate is not identified with admission of every physical pair.

The original barycentric comparison is stated first with its proof.
It remains available unchanged and is combined with, not replaced by,
the two endpoint-distance comparisons. The additional critical-trace
calculation distinguishes a positive critical entrance from a genuinely
grazing end. All clocks below use pullback Jacobians for the paired
measure; their integral covariance is printed explicitly.

\input{journal-inputs-v102}
\input{core/244_barycentric_positive_trace_coarea}
\input{core/245_transverse_clock_and_pointwise_trace_height}
\input{core/246_transverse_certificates_and_scalar_schedule}
\input{core/247_endpoint_distance_coarea_and_critical_traces}
\input{core/248_gap_clock_pointwise_recovery}
\input{core/249_expanding_endpoint_certificate}

\appendix
\input{appendices/v103_dependency_map}

\endgroup
\part{Positive paired coarea and the relative transverse clock}
\begingroup
\let\ArchiveLabel\label
\let\ArchiveRef\ref
\renewcommand{\label}[1]{\ArchiveLabel{v102archive:#1}}
\renewcommand{\ref}[1]{%
  \IfBeginWith{#1}{supp-}{%
    \StrBehind{#1}{supp-}[\ArchiveTarget]%
    \ArchiveRef{\ArchiveTarget}%
  }{\ArchiveRef{v102archive:#1}}%
}
\makeatletter
\let\ltx@label\label
\makeatother
\paragraph{Retained revision-102 abstract.}

We study the exact-label pointwise raw local limit problem for actual
returns in a triangular Lorentz gas. The remaining upper source is
the positive minimum of two physical endpoint traces on each short
regular segment. We represent this minimum by a barycentric coarea
measure and construct a transverse comparison which preserves relative
position between the two ends. Its Jacobian is computed before
physical sheets are summed. A capped transverse clock then gives an
essential-height estimate for the admitted positive paired source
itself, rather than for its outgoing-minus-incoming difference.
A relative endpoint certificate applies without a lower residence
cutoff. Under the original reconstruction schedule the admitted
contribution is $O(B^{-1/384})$. The unadmitted transverse source is
retained explicitly; no assertion that it vanishes on the full Lorentz
family is made. The original pointwise objective, exact labels,
physical guard, arithmetic classes and all historical proofs are
unchanged.


\input{journal-problem-v93}

\subsection{Positive paired mass and a transverse comparison}
The uniform momentum-capacity theorem has removed the singular
volume-source premise from the preceding scalar reduction. Its
remaining boundary quantity is positive: two large endpoint traces
can have a small difference. We therefore put the positive minimum
itself into a coarea measure. The new comparison moves in angular
direction, keeping fixed the relative position between the same two
physical ends. It does not require a radial trajectory to survive
beyond its original short residence.

Theorem~\ref{thm:v102-pair-coarea} constructs this measure and its
Jacobian. Theorem~\ref{thm:v102-admitted-height} estimates its whole
admitted part in essential height after physical summation, with no
word-count or chart-overlap factor. Theorem~\ref{thm:v102-certificate}
gives a test in terms of the two endpoint roof maps and their positive
trace density. Its constants do not deteriorate merely because the
original residence tends to zero. Theorem~\ref{thm:v102-scalar}
inserts the estimate into the unchanged scalar target.

Admission on the complete nongrazing paired family is not established.
The remaining positive source includes short transverse comparison
lifetimes, folds of the endpoint roof maps, and insufficient control
of relative trace variation. A backward critical trace is not set to
zero. These distinctions are part of the theorem statements, not
changes of the governing objective.

\input{journal-inputs-v102}
\input{core/244_barycentric_positive_trace_coarea}
\input{core/245_transverse_clock_and_pointwise_trace_height}
\input{core/246_transverse_certificates_and_scalar_schedule}

\appendix
\input{appendices/v102_dependency_map}

\endgroup
\part{Uniform geometric capacity and the positive paired trace}
\begingroup
\let\ArchiveLabel\label
\let\ArchiveRef\ref
\renewcommand{\label}[1]{\ArchiveLabel{v100archive:#1}}
\renewcommand{\ref}[1]{%
  \IfBeginWith{#1}{supp-}{%
    \StrBehind{#1}{supp-}[\ArchiveTarget]%
    \ArchiveRef{\ArchiveTarget}%
  }{\ArchiveRef{v100archive:#1}}%
}
\makeatletter
\let\ltx@label\label
\makeatother
\paragraph{Retained revision-100 abstract.}

We study the exact-label pointwise raw local limit problem for actual
returns in a triangular Lorentz gas. A finite transverse testing family,
together with uniform growth and whole-cell distortion estimates,
bounds the geometric stable-Jacobian pressure uniformly over the
prescribed interval of radii and every collision count. The construction
uses no parameter-uniform mixing time. The previously established
marked comparison then bounds the actual weighted endpoint-momentum
capacity uniformly, discharging its spectral-capacity condition at
every fixed auxiliary band. A positive two-cap layer cake gives uncut
first inverse-gradient concentration and nearby negative-gradient
moments on the original exact-label source. The long critical residual
and the full short-trace comparison error consequently have uniform
essential-height bounds without an additional capacity premise.
The original pointwise objective is retained. The positive paired
minimum on nongrazing short segments remains a distinct unproved
height estimate; no unrestricted scalar or point-conditioned law is
inferred from capacity alone. All preceding mathematics and both
historical source variants are preserved.


\input{journal-problem-v93}

\subsection{Uniform pressure on the original radius family}
The principal target remains \eqref{eq:v76-target}. The complete
finite-count gradient cap, critical target certificate and paired
physical trace comparison are retained below with their proofs.
The global momentum uniqueness and marked pressure comparison are
also retained: they are the precise bridge from dynamical pressure
to the original optical fiber, not an analogy with a symbolic count.

The new geometric argument is Theorem~\ref{thm:v100-uniform-pressure}.
It compares every long backward continuity cell to one of finitely
many short stable test curves and pays their homogeneous inverse
branches by a uniform growth estimate. Dropping only the final stable
factor leaves a uniformly finite physical split. This removes the
need to infer uniform constants from a fixed-table mixing argument.
The resulting bound is simultaneous in radius and collision count.

Theorem~\ref{thm:v100-capacity} discharges the full-family spectral
capacity condition. Theorem~\ref{thm:v100-uncut} gives the actual
uncut source with weight $I_mz^{-p}$ for $1\le p<129/128$; this is
not a higher power of $I_m/z$. Theorem~\ref{thm:v100-recovery} then
bounds the complete critical residual and every selected short-source
comparison error in essential height, uniformly in the radius.
These are new family-wide conclusions, not an interchange of a
fixed-radius limit with a supremum.

The remaining term in Theorem~\ref{thm:v100-endpoint} is the positive
minimum of the two original weighted traces on each nongrazing short
segment. Its decay is not proved by the pressure theorem. The original
uniform pointwise objective is not replaced by the fixed-radius or
finite-norm consequences already in the archive.

\input{journal-inputs-v94}
\input{journal-residual-v98}
\input{core/235_quadratic_gradient_cap_and_uncut_insertion}
\input{core/236_disjoint_residuals_and_critical_height_certificate}
\input{core/237_paired_physical_traces_and_scalar_reduction}

\input{core/238_global_momentum_uniqueness_and_stable_weight}
\input{core/239_marked_pressure_comparison_and_growth}
\input{core/240_fixed_radius_recovery_and_uniform_scope}

\input{core/241_uniform_geometric_pressure_by_finite_tests}
\input{core/242_uniform_capacity_and_uncut_gradient_moments}
\input{core/243_uniform_source_recovery_and_positive_trace}

\appendix
\input{appendices/v100_dependency_map}

\endgroup
\part{Global momentum uniqueness and marked physical pressure}
\begingroup
\let\ArchiveLabel\label
\let\ArchiveRef\ref
\renewcommand{\label}[1]{\ArchiveLabel{v99archive:#1}}
\renewcommand{\ref}[1]{%
  \IfBeginWith{#1}{supp-}{%
    \StrBehind{#1}{supp-}[\ArchiveTarget]%
    \ArchiveRef{\ArchiveTarget}%
  }{\ArchiveRef{v99archive:#1}}%
}
\makeatletter
\let\ltx@label\label
\makeatother
\paragraph{Retained revision-99 abstract.}

We study the exact-label pointwise raw local limit problem for actual
returns in a triangular Lorentz gas. Convexity of the Cartesian flight
action proves global uniqueness for a prescribed lifted itinerary and
two endpoint momenta. A comparison of the inverse optical twist with
the invariant stable Jacobian then controls the actual weighted
momentum-fiber multiplicity by a marked pressure sum. The mark removes
one stable factor; a finite physical step alphabet pays the resulting
multiplicity without an infinite homogeneity-strip count. Established
fixed-table pressure bounds imply bounded optical capacity at every
fixed radius. Thus the first inverse-gradient insertion, long critical
residual height and grazing-ended short-source estimates no longer
require a capacity hypothesis at fixed radius. A submultiplicative
finite-block theorem separately gives a radius-uniform exponential
growth bound and a sufficient spectral margin. The original uniform
pointwise target is unchanged: uniform certification of that margin
and decay of the positive nongrazing paired short trace remain to be
proved. All earlier proofs and both historical source variants are
retained.


\input{journal-problem-v93}

\subsection{The weighted multiplicity and the pressure comparison}
The principal target is \eqref{eq:v76-target}, with the original source,
normalization and complete arithmetic reference. The finite-count
quadratic cap and both physical residual mechanisms are retained below
with their proofs. Their capacity condition concerns a weighted sum over
\emph{endpoint-momentum fibers}, not an ordinary symbolic word count.
Theorem~\ref{thm:v99-unique} removes the possible multiplicity inside
one physical itinerary by a global Cartesian convexity argument.
Lemma~\ref{lem:v99-stable-weight} then compares its precise optical
weight with the stable Jacobian and removes each singular incidence
at its marked time.

The resulting theorem compares the complete momentum-fiber capacity
with geometric partition sums. For each fixed radius it is bounded
uniformly in collision count by the established fixed-table pressure
theory; see Corollary~\ref{cor:v99-fixed-radius}. For the radius-uniform
problem, Theorem~\ref{thm:v99-block} gives a finite-block exponential
growth bound and its explicit strict spectral margin. These are
different quantifiers. The former is not silently promoted to the
latter by compactness of the parameter interval.

At fixed radius the new estimate makes the critical recovery and the
complete short-trace comparison unconditional relative to the retained
analytic and geometric inputs. The positive same-segment minimum on
short nongrazing segments remains an additional source-height question,
even at one radius. The original uniform theorem is not replaced by
the fixed-radius consequence, nor is the paired minimum identified
with a projected signed cancellation.

\input{journal-inputs-v94}
\input{journal-residual-v98}
\input{core/235_quadratic_gradient_cap_and_uncut_insertion}
\input{core/236_disjoint_residuals_and_critical_height_certificate}
\input{core/237_paired_physical_traces_and_scalar_reduction}

\input{core/238_global_momentum_uniqueness_and_stable_weight}
\input{core/239_marked_pressure_comparison_and_growth}
\input{core/240_fixed_radius_recovery_and_uniform_scope}

\appendix
\input{appendices/v99_dependency_map}

\endgroup
\part{Quadratic momentum caps and paired physical traces}
\begingroup
\let\ArchiveLabel\label
\let\ArchiveRef\ref
\renewcommand{\label}[1]{\ArchiveLabel{v98archive:#1}}
\renewcommand{\ref}[1]{%
  \IfBeginWith{#1}{supp-}{%
    \StrBehind{#1}{supp-}[\ArchiveTarget]%
    \ArchiveRef{\ArchiveTarget}%
  }{\ArchiveRef{v98archive:#1}}%
}
\makeatletter
\let\ltx@label\label
\makeatother
\paragraph{Retained revision-98 abstract.}

We retain the exact-label pointwise raw local limit problem for actual
returns in a triangular Lorentz gas. A physical area formula in the two
endpoint momenta gives a quadratic small-gradient cap with an explicit
weighted optical multiplicity. Combining this cap with the inherited
spectral estimate makes the first inverse-gradient layer cake integrable
to zero at every finite collision count, with its entire error retained.
We separate the preceding residual into long near-critical targets and
short physical segments. The first is paid by the new untruncated
insertion; the second is compared to the minimum of its two actual
weighted endpoint traces. The error of that comparison is a complete
physical volume term. An explicit spectral-capacity condition would
make these errors negligible and give a critical-residual height bound
of order $(qr)^{257/256}$. That condition is not verified, and the remaining paired short-trace
height is not proved to vanish for the full family. All
original sources, labels, arithmetic classes and earlier proofs remain
unchanged; the pointwise objective is not replaced by a finite-norm law.


\input{journal-problem-v93}

\subsection{The finite-count estimate and the two physical mechanisms}
The new estimate acts on the actual insertion $I_m/|g|$, without a
positive lower gradient cutoff. Its leading local term is
$C b^{1/128}(|J|+B_*^{-1})$. Its complete finite-count error is
bounded by an explicit function of the weighted two-dimensional
momentum-fiber multiplicity $\mathcal K_m$ and the fixed-band decay
factor $\theta_{B_*}^m$; see Theorem~\ref{thm:v98-uncut-insertion}.
The weighted multiplicity is finite at each count, but its growth has
not been bounded at the rate needed for scalar closure.

The residual is split disjointly by total physical residence, with the
overlap assigned to the short component. On the longer component a
low capped clock forces a definite decrease of the original angular
density on a full forward interval. On the short component we compare
the target density to its two physical endpoint traces using the same
Borel segment, with no trace on an artificial target-selector graph.
Both comparisons hold at finite count before any asymptotic condition.
The conditional height consequences and their exact additional
hypothesis are stated separately, in Theorems~\ref{thm:v98-critical-certificate}
and \ref{thm:v98-paired-trace}.

\input{journal-inputs-v94}
\input{journal-residual-v98}
\input{core/235_quadratic_gradient_cap_and_uncut_insertion}
\input{core/236_disjoint_residuals_and_critical_height_certificate}
\input{core/237_paired_physical_traces_and_scalar_reduction}

\appendix
\input{appendices/v98_dependency_map}

\endgroup
\part{Capped coarea occupation and failure-source recovery}
\begingroup
\let\ArchiveLabel\label
\let\ArchiveRef\ref
\renewcommand{\label}[1]{\ArchiveLabel{v97archive:#1}}
\renewcommand{\ref}[1]{%
  \IfBeginWith{#1}{supp-}{%
    \StrBehind{#1}{supp-}[\ArchiveTarget]%
    \ArchiveRef{\ArchiveTarget}%
  }{\ArchiveRef{v97archive:#1}}%
}
\makeatletter
\let\ltx@label\label
\makeatother
\paragraph{Retained revision-97 abstract.}

We study the upper pointwise local limit problem for the exact record
of actual returns in a triangular Lorentz gas. Capping the physical
coarea ratio before time integration gives an occupation anchor using
both available directions and only returns to the identical selected
event. This recovers targets excluded by a whole-interval or worst-time
test. Independently, we bound the running compression of the complete
physical angular density at each target roof on prescribed same-event
intervals. The positive-cutoff inverse-gradient input pays this maximum
without counting words or visits. Together these estimates quantitatively
reduce the preceding three-way failure source to a subsource near an
original entrance with small gradient, or on a short original segment.
The artificial endpoint collar has no contribution at a uniform small
radius. Every recovered contribution has an explicit vanishing height
bound along the reconstruction schedule. The residual essential height
is not proved zero; the original pointwise theorem and all exact physical
and arithmetic data remain unchanged.


\input{journal-problem-v93}

\subsection{Quantitative recovery of failed targets}
The new results concern the original failed-test source, not only a
newly selected family of regular targets. With $Z_{r,D,\beta}^\varepsilon$
denoting the preceding disjoint interruption, compression and
occupation failures, Theorem~\ref{thm:v97-reduction} constructs an
original subsource $Y^\varepsilon\le Z^\varepsilon$ and proves
\[
 \mathcal H_M(Z^\varepsilon-Y^\varepsilon)
 \le Cq^{-1}\{\varepsilon^{1/16}+\varepsilon^{-1}\sqrt r\}+Cqr,
 \qquad q=\beta/D\le1/2.
\]
This is essential height at the target roof. It is not an integrated
translation or exceptional-set estimate. The residual support is
specified by the original lifetimes in \eqref{eq:v97-support}.
The running maximum responsible for the $Cqr$ term is estimated in
Theorem~\ref{thm:v97-running} before taking the collision limit.

All new theorems use the following explicitly retained inputs. The
comparison-domain anchor uses no source estimate outside the exact
event. The singular-volume estimate for running compression is applied
only while the trajectory remains in that event. This distinction
prevents label changes from entering a fixed-label spectral estimate.

\input{journal-inputs-v94}
\input{journal-comparison-v97}
\input{core/232_capped_coarea_occupation}
\input{core/233_running_compression_height}
\input{core/234_quantitative_failure_source_reduction}

\appendix
\input{appendices/v97_dependency_map}
\enlargethispage{3\baselineskip}

\endgroup
\part{Exact-label physical continuation and occupation}
\begingroup
\let\ArchiveLabel\label
\let\ArchiveRef\ref
\renewcommand{\label}[1]{\ArchiveLabel{v96archive:#1}}
\renewcommand{\ref}[1]{%
  \IfBeginWith{#1}{supp-}{%
    \StrBehind{#1}{supp-}[\ArchiveTarget]%
    \ArchiveRef{\ArchiveTarget}%
  }{\ArchiveRef{v96archive:#1}}%
}
\makeatletter
\let\ltx@label\label
\makeatother
\paragraph{Retained revision-96 abstract.}

We study the remaining upper pointwise local limit problem for the exact
record of actual returns in a triangular Lorentz gas. A section-decision
end need not terminate the underlying regular physical word. We construct
a comparison flow across such decisions and prove an exact-label
occupation inequality at every finite collision count. Only sampled
states returning to the original selected event supply anchor mass;
no occupation labels are pooled. This estimates the positive two-nongrazing
short-segment packet itself on the bounded-compression recurrent class,
and estimates the entire recurrent target-critical core without a lower
gradient cutoff or a finite-annulus truncation. A prescribed-radius
pointwise inequality localizes the remaining source to physical
interruption, compression, and insufficient same-label occupation.
The original source, guard, arithmetic classes and pointwise objective
are unchanged. Exhaustion of all possible spikes by the controlled
class is not proved, and the unrestricted pointwise theorem is not
claimed complete.


\input{journal-problem-v93}

\subsection{The positive packet and its returning anchors}
The principal target remains \eqref{eq:v76-target}. The new question is
whether the positive two-ended packet can be paid by the original
source at other roofs, even when its own exact-event segment is short.
Theorem~\ref{thm:v96-anchor} answers this on a class defined by physical
continuation, a bound on the inverse physical Jacobian, and the time
spent in the same exact selected event. Its bound has no word-count or
return-count factor. Theorem~\ref{thm:v96-packet-decay} applies it to
the previously residual two-nongrazing minimum, not only to the error
between that minimum and the target source.

Theorem~\ref{thm:v96-source-height} also gives a target-critical estimate
with all gradients below the upper threshold present at finite count.
This closes the disappearing-annulus issue on the admitted class;
it does not claim that every critical target is admitted. The full
original source is partitioned exactly in
Theorem~\ref{thm:v96-pointwise}. The three failed tests there, rather
than a silently removed physical subsource, remain in the endpoint
budget. All prior results and their hypotheses are retained.

\input{journal-inputs-v94}
\input{core/229_same_label_continuation_occupation}
\input{core/230_recurrent_packets_and_untruncated_radial_core}
\input{core/231_continuation_failure_and_joint_forcing}

\appendix
\input{appendices/v96_dependency_map}

\endgroup
\part{Two-ended packets and intrinsic radial residence}
\begingroup
\let\ArchiveLabel\label
\let\ArchiveRef\ref
\renewcommand{\label}[1]{\ArchiveLabel{v95archive:#1}}
\renewcommand{\ref}[1]{%
  \IfBeginWith{#1}{supp-}{%
    \StrBehind{#1}{supp-}[\ArchiveTarget]%
    \ArchiveRef{\ArchiveTarget}%
  }{\ArchiveRef{v95archive:#1}}%
}
\makeatletter
\let\ltx@label\label
\makeatother
\paragraph{Retained revision-95 abstract.}

We continue the upper pointwise local limit problem for the exact record
of actual returns in a triangular Lorentz gas. On short physical segments,
the original weighted source is approximated in essential height by the
smaller of its two endpoint traces. At every fixed positive entrance
gradient this eliminates the zero-residence contribution of all segments
with a grazing endpoint. A two-sided anchor on the intrinsic radial scale
also estimates complete critical annuli without requiring a prescribed
shell crossing. The residual thin segments cannot originate at a critical
point. Finally we retain signed projection in both roof directions; the
minimum of the two resulting forcings is no larger than the previous
backward forcing and gives an equivalent pointwise criterion. These
results preserve the original source and every exact label. The remaining
two-nongrazing bridge and zero-gradient concentrations are not asserted
to vanish; the unrestricted pointwise theorem remains the objective.


\input{journal-problem-v93}

\subsection{Two physical ends rather than one positive inlet}
The target remains \eqref{eq:v76-target}. The new proof addresses the
short-segment and critical-core sources directly. It does not interpolate
a finite norm into an essential supremum or normalize the upper excess.
Theorem~\ref{thm:v95-packets} proves a complete, finite-count comparison
with a two-ended packet. One of its consequences is a vanishing
zero-residence height for every grazing-ended subfamily at fixed positive
entrance gradient. No positive lower residence cutoff is inserted.

Theorem~\ref{thm:v95-anchor} treats targets which remain inside a
previously prescribed critical shell but have enough residence on the
scale of their own gradient. It chooses an anchor on either physical
side and sums the physical energies with multiplicity one. The leading
bound in one annulus is $Cq^{129/128}+C(1+\varepsilon^{-1})q^{257/128}$.
Its finite-count remainder is displayed, so the difference between a
finite annulus sum and a zero-gradient limit is explicit.

Theorem~\ref{thm:v95-bidirectional} retains entrance and terminal signs
through projection in both directions. A critical entrance which was
positive in the backward balance is negative in the forward balance.
Their minimum is no larger than either individual forcing; it is
not a sum of two positive estimates. Its vanishing is equivalent to the
same original scalar law (Theorem~\ref{thm:v95-endpoint}); verification
of that condition still requires control of the remaining concentrations.

\input{journal-inputs-v95}
\input{core/226_two_ended_short_segment_packets}
\input{core/227_intrinsic_radial_residence_anchors}
\input{core/228_bidirectional_signed_pointwise_forcing}

\section{The remaining upper-source estimate}
The short noncritical source is now represented, up to a vanishing
height error, by a packet requiring two nongrazing physical traces.
At fixed positive entrance gradient its grazing-ended part has zero
ordered residence charge. Critical
annuli with sufficient relative residence have an additional
pointwise estimate even when they do not cross the previously
prescribed shell. The complementary relatively thin segments have
strictly positive entrance gradient.

The pointwise law is not concluded from these statements. A family
of positive annular estimates with a common leading constant does
not authorize removal of a lower gradient cutoff inside a collision
limsup. Nor does the minimum of two unknown signed forcings have to
vanish. The complete two-nongrazing short packet and a possible
concentration in disappearing critical annuli remain positive
source-level obligations. All finite-norm, one-sided lower and
conditioning theorems retain their historical scope; unrestricted
point-conditioned laws and arithmetic-zero normalization are not
promoted here. The full preceding article and every inherited proof
remain in the technical archive.

\appendix
\input{appendices/v95_dependency_map}

\endgroup
\part{Radial entrance packets and signed projected forcing}
\begingroup
\let\ArchiveLabel\label
\let\ArchiveRef\ref
\renewcommand{\label}[1]{\ArchiveLabel{v94archive:#1}}
\renewcommand{\ref}[1]{%
  \IfBeginWith{#1}{supp-}{%
    \StrBehind{#1}{supp-}[\ArchiveTarget]%
    \ArchiveRef{\ArchiveTarget}%
  }{\ArchiveRef{v94archive:#1}}%
}
\makeatletter
\let\ltx@label\label
\makeatother
\paragraph{Retained revision-94 abstract.}

We study the upper pointwise local limit problem for the exact collision
record of actual returns in a triangular Lorentz gas. The complete
noncritical entrance source has an essential-height estimate at every
fixed positive residence cutoff, with a leading constant independent
of both cutoffs. Its zero-residence contribution is identified with
the original weighted short-segment source. We also estimate all
critical histories which traverse a prescribed radial shell before
the target roof. Their entrance and growth are retained as a signed
packet; its actual target height is bounded without a residence or
incidence cutoff. Finally, projecting the complete signed physical
balance before taking positive parts gives a forcing term smaller
than the separate positive entrance and growth budgets. Its vanishing
is equivalent to the original pointwise assertion, in the prescribed
order of limits. The quantitative entrance and radial-transit estimates
do not establish vanishing of the remaining critical-target and
zero-residence concentration. The original pointwise theorem and all
exact physical data remain unchanged.


\input{journal-problem-v93}

\subsection{Entrance geometry and signed upper forcing}
The positive scalar remainder makes the upper sign the outstanding
part of \eqref{eq:v76-target}. The previous signed prefix balance
identified two sufficient positive mechanisms, entrance traces and
critical positive growth. Those positive pieces may overestimate the
projected source, since physical and prefix losses can cancel them.
The present argument retains that cancellation while estimating two
complete geometric families in essential height.

Theorem~\ref{thm:v94-entrance-height} controls every noncritical inlet
at each fixed positive residence cutoff by $C(r^{1/32}+r)$, uniformly
in that cutoff and the original guard. Proposition~\ref{prop:v94-inlet-charge}
identifies the precise uncut remainder as an original collision-source
restriction. Theorem~\ref{thm:v94-transit} controls the actual target
source of every history crossing the shell $\sqrt r\le z\le2\sqrt r$
within the backward window by
\[
 C r^{129/256}+C(1+\varepsilon^{-1})r^{257/256}.
\]
This includes genuine critical directional births; zero area at a
critical point is never used to erase its trace.

Theorems~\ref{thm:v94-signed-comparison} and \ref{thm:v94-endpoint}
then give a smaller, signed pointwise forcing and an equivalent upper
criterion. These are not finite-norm embeddings. The forcing remains
inside the essential supremum and collision limsup before its gradient
cutoff is removed. Its zero value is not proved by the two partial
height estimates.

\input{journal-inputs-v94}
\input{core/223_noncritical_entrance_residence}
\input{core/224_critical_radial_transit_packets}
\input{core/225_signed_projected_forcing_and_endpoint}

\section{The remaining upper-source estimate}
The entry theorem uses a complete physical-volume anchor for each
selected segment and only a finite residence-band sum before the
collision limit. The radial theorem obtains its anchor length from
actual passage through a gradient shell. Neither argument sums a
single-sheet modulus by an uncontrolled transverse multiplicity.
The signed theorem removes completed prefixes and preserves all
negative physical and artificial-cut contributions through projection.

The pointwise target nevertheless still requires vanishing of the
signed forcing in \eqref{eq:v94-equivalence}. Equation
\eqref{eq:v94-localization} records what is not estimated: a target
which remains in the small-gradient core, and the original weighted
zero-residence inlet concentration. Small fixed-window mass of these
sources does not bound their essential height. The equality of the
inlet and original-source charges is not a proof that either is zero.
Nor does a signed cancellation identity alone supply a bound for its
remaining positive part.

All original labels, section decisions, first-hit tests, weights and
arithmetic classes remain in \eqref{eq:v76-target}. The earlier lower
pointwise half, finite-norm laws and their conditioned consequences
retain their exact hypotheses. No point-conditioned likelihood, kernel,
unrestricted ratio or normalization of arithmetic zero classes is
promoted. The complete preceding article and every inherited proof
remain in the separate technical archive.

\appendix
\input{appendices/v94_dependency_map}

\endgroup
\part{Absolute optical traces and signed ancestral upper forcing}
\begingroup
\let\ArchiveLabel\label
\let\ArchiveRef\ref
\renewcommand{\label}[1]{\ArchiveLabel{v93archive:#1}}
\renewcommand{\ref}[1]{%
  \IfBeginWith{#1}{supp-}{%
    \StrBehind{#1}{supp-}[\ArchiveTarget]%
    \ArchiveRef{\ArchiveTarget}%
  }{\ArchiveRef{v93archive:#1}}%
}
\makeatletter
\let\ltx@label\label
\makeatother
\paragraph{Retained revision-93 abstract.}

We return to the upper pointwise local limit problem for the exact
collision record of actual returns in a triangular Lorentz gas.
The positive optical denominator cancels each individual incidence
singularity in the absolute angular density, giving a collision-uniform
one-half modulus on every physical sheet. Its complete fixed-count
signed balance shows that the original critical-or-short-residence
selector creates only negative interior cuts. Positive increments come
from actual entrance traces and positive physical density growth.
We prove a prescribed-scale pointwise comparison retaining only the
ancestors of the selected target fiber, together with an upper-spike
persistence dichotomy on that same source. The complete noncritical
growth term has ordered height of order the roof width, with a leading
constant independent of every fixed positive gradient cutoff. Actual
ancestral entrances and critical positive growth remain explicit in
the scalar upper bound. No vanishing of their concentration is inferred,
and the original unrestricted pointwise theorem remains to be completed.
All earlier finite-norm and conditioning results are retained.


\input{journal-problem-v93}

\subsection{The upper sign and its physical origins}
The organizing assertion is still \eqref{eq:v76-target}. The historical
one-sided lower law and the finite-norm conclusions remain available,
but do not answer the upper-spike question. The new argument studies
the original positive weighted density before taking its absolute
roof increments. It uses the monotonicity of the actual exceptional
selector along each maximal physical segment, not an embedding of a
finite norm into an essential supremum.

Theorems~\ref{thm:v93-balance} and \ref{thm:v93-pointwise} give its
signed balance and a pointwise inequality at a prescribed roof radius.
Theorem~\ref{thm:v93-gain} estimates the full noncritical positive
variation after the complete physical summation. Thus a high value
either persists on the prescribed backward interval or is paid by a
quantified entrance or critical-growth contribution of the same target's
ancestors. The final upper bound in Theorem~\ref{thm:v93-upper-endpoint}
keeps those contributions rather than declaring them small.

\input{journal-inputs-v93}
\input{core/220_absolute_optical_traces}
\input{core/221_signed_prefix_birth_balance}
\input{core/222_upper_birth_and_critical_gain}

\section{The original endpoint and the retained consequences}
The comparison just proved acts at almost every original roof, not
only after integrating starting roofs. It also retains all true
entrances and losses. An interior cutoff contributes a negative trace,
so the original exceptional union does not acquire an uncontrolled
positive cutoff-boundary term. These conclusions do not use the
positive-reference denominator or suppress upper amplitudes by taking
a reciprocal.

What is not proved is that the two positive forcing terms in
\eqref{eq:v93-endpoint} vanish. A uniform one-sheet modulus cannot be
summed with an unestimated angular multiplicity. A critical point's
zero area does not remove its possible directional entrance trace.
The remaining task is therefore a complete-source estimate of these
specific positive mechanisms or another source-specific upper bound.
No new unrestricted charge-zero or pointwise likelihood assertion is
made here.

All earlier weak, strong subcritical, logarithmic, diffusive-window,
one-sided lower, reciprocal and conditioning theorems retain exactly
their previous hypotheses. Both revision-88 variants remain distinct.
The technical archive includes the full preceding article and every
inherited proof; no alternate model, averaged target or pooled label
has replaced \eqref{eq:v76-target}.

\appendix
\input{appendices/v93_dependency_map}

\endgroup
\part{One-sided lower inversion and positive-reference conditioning}
\begingroup
\let\ArchiveLabel\label
\let\ArchiveRef\ref
\renewcommand{\label}[1]{\ArchiveLabel{v92archive:#1}}
\renewcommand{\ref}[1]{%
  \IfBeginWith{#1}{supp-}{%
    \StrBehind{#1}{supp-}[\ArchiveTarget]%
    \ArchiveRef{\ArchiveTarget}%
  }{\ArchiveRef{v92archive:#1}}%
}
\makeatletter
\let\ltx@label\label
\makeatother
\paragraph{Retained revision-92 abstract.}

We retain the pointwise local limit problem for the complete displacement,
collision-count, occupation and roof record of actual returns in a
triangular Lorentz gas. The retained positive-remainder lower law holds
pointwise and uniformly in radius and exact labels.
Consequently the actual roof density has a uniform lower bound wherever
the finite arithmetic reference has a prescribed positive floor; this
conclusion does not require an upper-height estimate. Combining it with
the unweakened critical weak limit yields strong convergence of the
reciprocal density in every finite local norm. Under the retained positive
path-remainder theorem, the same argument gives finite-norm convergence
of the exact same-roof collision kernels and quantitative posterior
bounds for the original exceptional union. Conditioning on arbitrary
moving measurable roof sets with fixed positive reference mass is stable
in the joint roof--path bounded-Lipschitz metric. The roof marginals also
converge in both directions of relative entropy. All source normalizations,
exact labels and arithmetic classes are retained. Uniform exclusion of
upper roof spikes, and hence the full pointwise theorem, remains a separate
unproved requirement.



\input{journal-problem-v89}

\subsection{The pointwise conclusion and the conditioning conclusions}
The historical lower law, restated in Theorem~\ref{thm:v92-lower}, gives
\[
 \sup_{R,n,k}\esssup_{\mathcal C(M)}[\mathcal L_{m,R}-m^2p_{n,R}]_+
                                                     \longrightarrow0.
\]
This is a pointwise statement, deduced from the sign of the physical
remainder rather than from a finite-norm embedding. It proves a uniform
physical denominator floor on each quantified positive-reference class.
It does not control the opposite sign of the raw error.

The ensuing normalization results retain that distinction.
Theorem~\ref{thm:v92-reciprocal} proves every finite-norm limit for the
bounded reciprocal error. Theorem~\ref{thm:v92-kernels} concerns the
bounded-Lipschitz distance between the original same-roof collision path
law and its retained bridge reference, not path-space total variation.
Theorem~\ref{thm:v92-entropy} concerns relative entropy of the roof
marginals, not entropy of path laws. Their full central-window versions
use exactly $dt/\sqrt m$.

These results use neither a new source nor an unproved pointwise upper
bound. In particular they do not replace \eqref{eq:v76-target} by a
window theorem. They give additional consequences of the original
positive source and its now-established critical weak convergence.

\input{journal-inputs-v92}
\input{core/217_one_sided_pointwise_inversion}
\input{core/218_positive_disintegration_and_posterior_stability}
\input{core/219_moving_roof_conditioning_and_entropy}

\section{The remaining pointwise requirement}
The retained estimate controls the absolute negative part in the
original uniform essential-supremum norm. Thus the unresolved part of
\eqref{eq:v76-target} is entirely an upper-excess problem. Its original
weighted-source charge and its order of limits are unchanged. Neither
Theorem~\ref{thm:v92-lower} nor the conditional finite-norm and entropy
results proves that this charge is zero. The upper-pulse example in
Remark~\ref{rem:v92-scope} makes the remaining distinction explicit.

Every reconstruction band is fixed before the collision limit and is
sent to infinity last. Reference floors and reference-mass floors are
also fixed, not chosen as shrinking functions of collision count.
The conditional assertions are uniform on those stated classes, not on
all sets of arbitrarily small positive mass. The actual-return clock
and any uniform forward-likelihood assertion retain their separate
hypotheses. No arithmetic zero class is normalized.

\appendix
\input{appendices/v92_dependency_map}

\endgroup
\part{The unweakened critical limit and all-radius source control}
\begingroup
\let\ArchiveLabel\label
\let\ArchiveRef\ref
\renewcommand{\label}[1]{\ArchiveLabel{v91archive:#1}}
\renewcommand{\ref}[1]{%
  \IfBeginWith{#1}{supp-}{%
    \StrBehind{#1}{supp-}[\ArchiveTarget]%
    \ArchiveRef{\ArchiveTarget}%
  }{\ArchiveRef{v91archive:#1}}%
}
\makeatletter
\let\ltx@label\label
\makeatother
\paragraph{Retained revision-91 abstract.}

We retain the pointwise local limit problem for the exact displacement,
collision-count, occupation and roof record of actual returns in a
triangular Lorentz gas. We prove convergence of the original raw error
in the unweakened critical weak norm of exponent $32/31$, removing the
previous logarithmic loss. The decisive step glues the complete surface
estimate on roof windows matched to the exit width before taking the
collision limit; the resulting surface norm decays as $r^{1/32}$ on
every fixed window. A positive domination argument then gives the
original critical-or-short-residence source a polynomial critical-norm
bound, uniformly over its smaller cutoffs and its original weights.
A maximal version controls every local averaging radius outside one
quantified exceptional roof set. The full central-window limit is
stated in diffusive roof measure. All exact labels, arithmetic classes
and prior proofs are retained. These finite-critical conclusions do
not assert that the remaining essential-height charge is zero.



\input{journal-problem-v89}

\subsection{The unweakened critical limit}
The new critical surface decay is proved in
Theorem~\ref{thm:v91-surface}. Its input is a scale-matched bound
already valid for the complete one-sided first-exit surface. A finite
partition into windows of that same length supplies the missing power
on a fixed roof interval. Theorem~\ref{thm:v91-raw} then proves
\[
 \sup_{R,n,k}\sup_{|J|=L}
 \mathcal N_{32/31}\bigl((m^2p-\mathcal L)\one_{\mathcal C(M)};J\bigr)
                                      \longrightarrow0
\]
for every fixed $M,L>0$, without a logarithmic denominator. This is
a theorem at a finite distribution exponent, not a redefinition of
\eqref{eq:v76-target}. The same result holds on the complete central
window with the explicit measure $dt/\sqrt m$.

For the original weighted exceptional source, the cofinal cutoffs
$z<h^{1/8}$ or $\ell<h^{1/4}$ give a critical norm of order
$h^{31/4096}$ in the ordered collision limit. The bound concerns the
union itself, not just common survivors. The all-radius consequence
estimates a single exceptional roof set, while leaving the absence
of essential-height concentration as a separate pointwise obligation.

\input{journal-inputs-v91}
\input{core/214_scale_gluing_and_unweakened_critical_limit}
\input{core/215_original_source_critical_power}
\input{core/216_all_radius_maximal_source_control}

\section{Proof dependencies and preserved conclusions}
The three new steps are scale-matched gluing, bounded-reference
domination, and an interval covering argument. The first uses the
complete physical surface theorem, including its capped residence
sum. The second uses the resulting raw critical convergence and the
actual critical/residence mass bounds. The third is applied to that
same original union before the collision limit. None uses the
pointwise theorem as an input.

The complete previous article and all inherited physical, spectral,
arithmetic and source-oscillation proofs are retained in the technical
supplement. The earlier logarithmic critical theorem remains valid
but is strengthened by Theorem~\ref{thm:v91-raw}; its historical
status is not edited. The original pointwise theorem still requires
a uniform source-specific exclusion of collapsing high roofs.
Independent verification of the inherited billiard and spectral
inputs remains distinct from source preservation and finite checks.


\endgroup
\part{All-scale original-source oscillation}
\begingroup
\let\ArchiveLabel\label
\let\ArchiveRef\ref
\renewcommand{\label}[1]{\ArchiveLabel{v90archive:#1}}
\renewcommand{\ref}[1]{%
  \IfBeginWith{#1}{supp-}{%
    \StrBehind{#1}{supp-}[\ArchiveTarget]%
    \ArchiveRef{\ArchiveTarget}%
  }{\ArchiveRef{v90archive:#1}}%
}
\makeatletter
\let\ltx@label\label
\makeatother
\paragraph{Retained revision-90 abstract.}

We retain the pointwise local limit problem for the complete displacement,
collision-count, occupation and roof record of actual returns in a
triangular Lorentz gas. We prove finite-count oscillation estimates on
the original collision source, including both true physical ends and
backward critical entrances. A capped end-layer estimate and radial
motion yield a uniform $L^1$ translation modulus for the full weighted
roof density. A one-endpoint-per-segment argument then gives a modulus
of order $h^{1/33}$ for every measurable initial segment restriction,
uniformly in the endpoint location. This includes the original union
of initial criticality and short total residence, with both cutoffs
arbitrarily small. Its complete microscopic fractional-variation tail
is bounded before the collision limit, and its fractional local norm
has a quantitative vanishing bound. All exact labels, physical weights,
arithmetic classes and previous proofs are retained. The new estimates
control integrated roof oscillation; the unchanged essential-height
charge still requires a uniform starting-roof estimate.



\input{journal-problem-v89}

\subsection{An all-scale estimate on the original exceptional source}
Let $\tau^-,\tau^+$ be the stopped roof-flow lifetimes, $\ell=\tau^-+\tau^+$
and $z=|\nabla_qF|$. The original exceptional source is
\[
 R_{a,\delta}^\varepsilon(t)
   =\int W^\varepsilon(x)\one_{\{z(x)<a\ \mathrm{or}\ \ell(x)<\delta\}}
                                                       \,d\mu_t(x).
\]
Its union is counted once. Neither endpoint section membership nor a
true boundary is removed. Our principal new assertion is the following.
\begin{theorem}[Microscopic oscillation on the complete original union]
\label{thm:v90-main}
Under the precise physical and spectral inputs stated below, for every
fixed $L$ and $0<\varepsilon\le1$ there is $K_L<\infty$ such that
\[
 \sup_{m\ge2,R,n,k}\sup_{a,\delta>0}\sup_{|J|\le L}m^2
 \int_J|R_{a,\delta}^\varepsilon(t+h)-R_{a,\delta}^\varepsilon(t)|\,dt
       \le K_L\bigl(h^{1/33}+\varepsilon^{-1}\sqrt h\bigr)
\]
for every sufficiently small $h>0$. The constant is obtained at one
fixed auxiliary band. The supremum over finite collision counts and
both cutoffs precedes $h\downarrow0$.
For $0<s<1/33$, the full integral of these differences against
$h^{-1-s}dh$ down to zero has a uniform vanishing microscopic tail.
For the cofinal choice $a=r^{1/8}$, $\delta=r^{1/4}$,
\[
 \limsup_m\sup_{R,n,k}\sup_{|J|=L}
 \|m^2R_{a,\delta}^\varepsilon\|_{s,1;J}
             \le C_{L,s}(1+\varepsilon^{-1})^{33s}r^{(1-33s)/128},
\]
where the fractional norm is defined in \eqref{eq:v90-gagliardo}.
\end{theorem}
The proof supplies a genuine pre-limit envelope, rather than a sum of
fixed-scale limiting estimates. The complete density $b^\varepsilon$
has the stronger translation exponent $1/32$. The loss to $1/33$ for
the original union comes from the single new interior cutoff on a
physical segment. Its location needs no regularity in angle.

These are oscillation statements in the roof variable, not another
claim that a local mass estimate is a pointwise bound. In particular the
integral over starting roofs is visible in the theorem. The exact
pointwise objective remains \eqref{eq:v76-target}; no unconditional
point-conditioned law is deduced without its essential-height closure.

\input{journal-inputs-v90}
\input{core/211_two_sided_end_layers_and_source_oscillation}
\input{core/212_cutoff_uniform_original_source_prefixes}
\input{core/213_microscopic_fractional_variation_and_endpoint_scope}

\section{Dependency and scalar-first order}
The new argument has three steps. The two-sided end-layer theorem pays
all true boundary changes, including critical backward ends. The
weighted segment calculation controls surviving variation without an
initial inverse-gradient cutoff. The prefix theorem pays the extra
cut introduced by the original exceptional union without a derivative
of its angular graph. Every estimate is summed over actual physical
interiors using the same Borel disintegration.

The preceding critical-distribution and diffusive-window theorems remain
in the complete preceding article in the supplement. They are not
replaced, and the two revision-88 variants remain separately preserved.
The new physical oscillation envelope does not establish zero
essential-height charge. A further uniform starting-roof or comparable
pointwise persistence bound is still required for that conclusion.
Only after scalar closure may the retained positive-arithmetic
normalization, same-roof bridge, clock transfer and likelihood arguments
be used for unconditional point-conditioned laws.


\endgroup
\part{Critical distribution and the complete diffusive window}
\begingroup
\let\ArchiveLabel\label
\let\ArchiveRef\ref
\renewcommand{\label}[1]{\ArchiveLabel{v89archive:#1}}
\renewcommand{\ref}[1]{%
  \IfBeginWith{#1}{supp-}{%
    \StrBehind{#1}{supp-}[\ArchiveTarget]%
    \ArchiveRef{\ArchiveTarget}%
  }{\ArchiveRef{v89archive:#1}}%
}
\makeatletter
\let\ltx@label\label
\makeatother
\paragraph{Retained revision-89 abstract.}

We study the pointwise local limit problem for the complete displacement,
collision-count, occupation and roof record of actual returns in a
triangular Lorentz gas. Retaining the measure of a tested roof set
through the complete residence sum yields a critical weak estimate
at exponent $32/31$ for the untruncated first-exit surface. Its
finite-count remainder remains summable down to residence zero.
The original raw density consequently has a uniform critical weak
bound. Together with the established local mass convergence, this
proves raw-error convergence at the critical exponent with every
positive logarithmic loss. The same logarithmic estimate acts directly
on the original critical-or-short-residence source. We also prove
convergence over the full central window in diffusive roof measure,
with the scale factor explicit. All exact labels, arithmetic classes,
physical weights and earlier proofs are retained. The organizing
pointwise assertion is unchanged; its residual essential-height
charge is not inferred to vanish from these distribution bounds.



\input{journal-problem-v89}

\subsection{Complete sources without a positive residence cutoff}
The pointwise error in \eqref{eq:v76-target} remains the organizing
objective. The new assertions estimate the same raw density in finite
local norms and the actual outgoing surface with all positive residence
lengths retained. These statements are not conditional on a microscopic
charge being zero.

For $1<q<\infty$ let $\mathcal N_q(f;J)$ be the measurable-set norm
in \eqref{eq:v88-weak-norm}, equivalent to weak $L^q$ on $J$.
For the original first-exit surface write $\overline{\mathcal F}_r$
for its average over exit widths between zero and $r$, as defined below.
\begin{theorem}[Untruncated surface and raw-density estimates]
\label{thm:v88-main}
Under the explicit inherited inputs of the next section, for every
$1<q<32/31$ and fixed $L>0$,
\[
 \limsup_m\sup_{R,n,k}\sup_{|J|=L}
 \mathcal N_q(m^2\overline{\mathcal F}_r;J)
                      \le C_{q,L}r^{1/q-31/32}.
\]
There is no lower gradient, residence or incidence cutoff in this
estimate. For every fixed $M,L>0$, the raw scalar error satisfies
\[
 \sup_{R,n,k}\sup_{|J|=L}
 \mathcal N_q\!\left(
 (m^2p_{n,R}(k,m,\cdot)-\mathcal L_{m,R}(k_1,k_2,\cdot,n))
 \one_{\mathcal C_{m,R,n,k}(M)};J\right)\longrightarrow0.
\]
The corresponding strong $L^p$ convergence holds for every
$1\le p<32/31$. All arithmetic classes remain unnormalized.
\end{theorem}
The proofs are Theorems~\ref{thm:v88-full-surface} and
\ref{thm:v88-raw-limit}. The new source theorem first establishes a
negative moment of total regular residence by a capped layer cake at
finite collision count. Its geometric trace application uses the
whole same-label segment; the source term is not replaced by a
boundary collar. A finite $q$ is essential to the conclusion. The
unchanged essential-supremum problem remains separate.

\subsection{Critical distribution and the original exceptional source}
Set $Q=32/31$ and $\gamma=1/32$. Theorem~\ref{thm:v89-endpoint-surface}
proves a finite-count measurable-set estimate for the complete surface
at the previously excluded exponent $Q$. Its proof keeps the factor
$\min\{|E|,\ell\}$ until after the residence sum; it does not
substitute an endpoint negative moment into a subcritical theorem.
Proposition~\ref{prop:v89-density-bound} then gives a critical weak
bound for the raw density itself.

The logarithmic norm in \eqref{eq:v89-log-norm} yields two further
statements. The raw error tends to zero in that norm for every
positive logarithmic exponent (Theorem~\ref{thm:v89-log-raw}).
The original weighted union of initial criticality and short total
residence has an explicit logarithmic modulus in the same critical
norm (Theorem~\ref{thm:v89-exceptional-log}). The latter estimate
acts on the exceptional source, not only its complementary surviving
fibers. Theorem~\ref{thm:v89-diffusive} gives the corresponding full
central-window conclusions with measure $dt/\sqrt m$.

These statements preserve \eqref{eq:v76-target} as the pointwise
objective. Neither a logarithmically weakened critical norm nor the
normalized central-window theorem is identified with essential height.
The exponent $32/31$ comes from the available residence exponent,
not a claim of physical optimality.

\input{journal-inputs-v89}

\input{core/205_capped_negative_residence_moments}
\input{core/206_uncut_surface_lorentz_control}
\input{core/207_exact_label_raw_density_local_norm}

\input{core/208_critical_surface_distribution}
\input{core/209_logarithmic_endpoint_and_exceptional_source}
\input{core/210_diffusive_window_and_exact_limit_order}

\section{The unchanged pointwise endpoint}
The new negative residence theorem is a finite-count assertion for an
actual singular insertion. Its full layer cake includes arbitrarily
short physical segments before the collision limit. The outgoing
trace theorem similarly sums all incident components before applying
a norm. Thus the complete surface weak-norm decay and raw local-norm
limit proved here require no positive residence cutoff.

These results do not identify a finite local norm with essential
height. The endpoint in \eqref{eq:v76-target} continues to require
\[
 \lim_{B\to\infty}\mathfrak R_{\varepsilon(B),M}=0,
\]
where $\mathfrak R$ is the original weighted-source concentration
charge of Supplement Theorem~\ref{supp-v87archive:thm:v87-endpoint}.
The critical distribution and logarithmic estimates constrain this same source but do not
exclude narrow peaks in its essential supremum. The older sufficient
pairs $(\mathfrak A,\mathfrak B)$ and $(\mathfrak A,\mathfrak Q)$
remain valid with their exact hypotheses in the complete archive.
No charge-zero assertion is inferred here.

In particular, a strong local $L^p$ theorem for $p<32/31$ is not a
uniform denominator theorem on positive arithmetic classes. The
scalar-first path remainder, collision-to-return clock transfer and
forward likelihood statements are preserved, not promoted to new
unconditional same-roof ratios. Zero arithmetic classes are never
normalized. All previous mathematical files and the complete preceding
article remain in the supplement and provenance.


\endgroup
\part{Untruncated residence moments and the raw local norm}
\begingroup
\let\ArchiveLabel\label
\let\ArchiveRef\ref
\renewcommand{\label}[1]{\ArchiveLabel{v88archive:#1}}
\renewcommand{\ref}[1]{%
  \IfBeginWith{#1}{supp-}{%
    \StrBehind{#1}{supp-}[\ArchiveTarget]%
    \ArchiveRef{\ArchiveTarget}%
  }{\ArchiveRef{v88archive:#1}}%
}
\makeatletter
\let\ltx@label\label
\makeatother
\paragraph{Retained revision-88 abstract.}

We study the pointwise local limit problem for the complete displacement,
collision-count, occupation and roof record of actual returns in a
triangular Lorentz gas. A count-independent physical cap for very short
regular residence can be combined with the finite-count spectral bound
before taking either limit. This yields actual negative residence moments
at every order below $1/32$. Integrating the exact angular-density
variation over a whole maximal segment converts those volume moments
into a weighted outgoing trace without an initial-gradient cutoff.
An interval-pulse inequality then proves decay in weak $L^q$ for the
complete, untruncated first-exit surface, for every $1<q<32/31$.
The same range gives uniform local-norm convergence of the original
raw exact-label density to its full arithmetic reference, with strong
$L^p$ convergence for $1\le p<32/31$. The original pointwise objective
is unchanged: these complete finite-norm theorems do not yet establish
vanishing of its critical and zero-residence essential-height charge.



\section{The unchanged physical problem}
Put $R_-=9/20$ and $R_+=47/100$, and let $R\in I=[R_-,R_+]$. Circular scatterers have centers in
$\Lambda=\Z(1,0)+\Z(1/2,\sqrt3/2)$. Write $T_R$ for the regular
collision map and $(\alpha,p)\in(\R/2\pi\Z)\times(-1,1)$ for its
coordinates; $p$ is the outgoing tangential velocity and
$\dd\nu=\dd\alpha\,\dd p/(4\pi)$. Physical flights satisfy
$\tau_-\le\tau_R\le\tau_+$, where $\tau_-=3/50$ and $\tau_+$
is the fixed finite-horizon bound.

The section is specified here to fix every event in the argument.
Put $h_0=1/200$. Let $Y_R$ be the union of the four open coordinate
squares of half-side $h_0$ centered at
\[
 (0,0),\quad(\pi/6,-1/2),\quad
 (-\theta_R,\sin\theta_R),\quad
 (\pi/3-\theta_R,\sin\theta_R).
\]
Here $\theta_R=2\arctan q_R$, and $q_R$ is the unique root in
$(49/100,53/100)$ of
\[
 8Rq(1+q^2)-4q(1-q^2)+\sqrt3(1-6q^2+q^4)=0.
\]
Set
\[
 U=\{|\alpha-\pi/6|<3/50,\ |p|<3/20\},\qquad
 Y_R^*=Y_R\cup U,\qquad c_*=\frac{91}{10000\pi}.
\]
The rectangles are disjoint and $\nu(Y_R^*)=c_*$. We always use
$\nu_R^*=c_*^{-1}\nu|_{Y_R^*}$, not a renormalization of a smaller
source.

At $m$ flights, let $L=L_{m,R}$ be their total length,
$K^{\rm c}_{m,R}$ their lattice displacement, and
$A_{m,R}=\sum_{j=0}^{m-1}\one_{Y_R^*}\circ T_R^j$. The event is
\begin{equation}\label{eq:v76-event}
 \Pi_{n,k,m,R}=Y_R^*\cap\{K^{\rm c}_{m,R}=k,\ A_{m,R}=n\}
                         \cap T_R^{-m}Y_R^*.
\end{equation}
The initial visit is counted; the terminal visit is tested separately.
On this event the actual $n$th return has collision count $m$ and
roof $L$. Denote its density by $p_{n,R}(k,m,t)$.

Write $c_j=\sqrt{1-p_j^2}$ and let $\Delta_j$ be the least clearance
of flight $j$ from any third disk, excluding its two incident disks.
Keep the original smooth function $0\le\Theta\le1$, zero on
$(-\infty,1]$ and one on $[2,\infty)$. With $q_0=47/53$,
$d_j=\min(j,m-j)$ and $d'_j=\min(j,m-1-j)$, set
\begin{align}
 H^\varepsilon&=
 \prod_{j=0}^{m}\Theta\!\left(\frac{c_j}{\varepsilon q_0^{d_j/4}}\right)
 \prod_{j=0}^{m-1}\Theta\!\left(\frac{\Delta_j}{\varepsilon q_0^{d'_j/4}}\right),
 \qquad W^\varepsilon=1-H^\varepsilon,\label{eq:v76-guard}\\
 b^\varepsilon(t)\dd t
   &=L_\#(W^\varepsilon\one_{\Pi_{n,k,m,R}}\nu_R^*).
 \label{eq:v76-source}
\end{align}
In particular first incidence and first clearance outside any selected
critical neighborhoods are still present. No section-distance factor
has been inserted into this guard.

\subsection{The arithmetic reference and the target}
For $y=(k_1,k_2,t,n)$ put
\[
 \mu_R=(0,0,\bar\tau_R,c_*),\quad
 Z=(y-m\mu_R)/\sqrt m,\quad
 \mathcal C_{m,R,n,k}(M)=\{t:|Z|\le M\}.
\]
We retain the finite spectral-peak reference, including peaks close to
but not on the unit circle. To specify its meaning, the collision
record observable is $(\kappa_1,\kappa_2,\tau,\one_{Y_R^*})$.
Let $\lambda_j(R,\vartheta)$ and $\mathsf P_j(R,\vartheta)$ be its
continued simple twisted eigenvalue and projection in each of the
fixed peak charts, and let $\chi_j$ be their frequency partition.
At the real maximum $a_j$ of $\log|\lambda_j|$ put
\[
 z_j=\lambda_j(R,a_j),\quad
 \mu_j=i^{-1}\nabla\log\lambda_j(R,a_j),\quad
 H_j=-\nabla^2\log\lambda_j(R,a_j),
\]
with $\operatorname{Re}H_j\ge aI_4$ for a common $a>0$. If
$\eta_R=\one_{Y_R^*}$ and $\ell$ denotes mass evaluation, put
$B_j=\chi_j(R,a_j)\ell(M_{\eta_R}\mathsf P_j(R,a_j)(\eta_R\nu))$.
Define
\begin{align}
 G_H(z)&=(2\pi)^{-4}\int_{\R^4}e^{-iv\cdot z-v^{\mathsf T}Hv/2}\dd v,
 \notag\\
 \mathcal L_{m,R}(y)&=\operatorname{Re}\left[
 c_*^{-1}\sum_j e^{-ia_j\cdot y}z_j^m B_j
      G_{H_j}\left((y-m\mu_j)/\sqrt m\right)\right].
 \label{eq:v76-reference}
\end{align}
Only the charts active at $R$ are summed. This is the same transition
reference as in the retained spectral theorem. The existence and
uniform estimates for these scalar data belong to that analytic
input, not to the transport argument below. No frequency branch or
arithmetic zero class is deleted.

The governing assertion remains
\begin{equation}\label{eq:v76-target}
 \sup_{R,\,1\le n\le m,\,k\in\Z^2}
 \esssup_{t\in\mathcal C_{m,R,n,k}(M)}
 |m^2p_{n,R}(k,m,t)-\mathcal L_{m,R}(k_1,k_2,t,n)|\longrightarrow0.
\end{equation}
Let $\mathcal E_M$ denote this error's collision limsup. An empty
central target contributes zero. The order of the uniformity is
$\sup_R\esssup_t$, so exceptional density sets may depend on $R$.


\subsection{Complete sources without a positive residence cutoff}
The pointwise error in \eqref{eq:v76-target} remains the organizing
objective. The new assertions estimate the same raw density in finite
local norms and the actual outgoing surface with all positive residence
lengths retained. These statements are not conditional on a microscopic
charge being zero.

For $1<q<\infty$ let $\mathcal N_q(f;J)$ be the measurable-set norm
in \eqref{eq:v88-weak-norm}, equivalent to weak $L^q$ on $J$.
For the original first-exit surface write $\overline{\mathcal F}_r$
for its average over exit widths between zero and $r$, as defined below.
\begin{theorem}[Untruncated surface and raw-density estimates]
\label{thm:v88-main}
Under the explicit inherited inputs of the next section, for every
$1<q<32/31$ and fixed $L>0$,
\[
 \limsup_m\sup_{R,n,k}\sup_{|J|=L}
 \mathcal N_q(m^2\overline{\mathcal F}_r;J)
                      \le C_{q,L}r^{1/q-31/32}.
\]
There is no lower gradient, residence or incidence cutoff in this
estimate. For every fixed $M,L>0$, the raw scalar error satisfies
\[
 \sup_{R,n,k}\sup_{|J|=L}
 \mathcal N_q\!\left(
 (m^2p_{n,R}(k,m,\cdot)-\mathcal L_{m,R}(k_1,k_2,\cdot,n))
 \one_{\mathcal C_{m,R,n,k}(M)};J\right)\longrightarrow0.
\]
The corresponding strong $L^p$ convergence holds for every
$1\le p<32/31$. All arithmetic classes remain unnormalized.
\end{theorem}
The proofs are Theorems~\ref{thm:v88-full-surface} and
\ref{thm:v88-raw-limit}. The new source theorem first establishes a
negative moment of total regular residence by a capped layer cake at
finite collision count. Its geometric trace application uses the
whole same-label segment; the source term is not replaced by a
boundary collar. A finite $q$ is essential to the conclusion. The
unchanged essential-supremum problem remains separate.

\section{Explicit analytic and geometric inputs}
We collect the indispensable inherited assertions in the exact form
used by the new proofs. Their source theorems and operator dependencies
are specified here and in Supplement Appendix~\ref{supp-v87archive:app:v83-kernels}. The new
potential and first-exit estimates do not assume a surface spectral
insertion theorem or an unweighted pointwise local bound.

\begin{inputstatement}[Physical flow and source comparison]
\label{input:v83-geometric}
On the regular part $\Omega$ of each exact event, endpoint angles
$q=(\alpha_0,\alpha_m)$ give $F=L$, $g=\nabla_qF=R(-p_0,p_m)$,
$K=D_q^2F$ and $\vartheta=\arg g$. Put
\[
 \lambda_\circ=9\sqrt{19}/200,\qquad \Lambda_\circ=47/6,\qquad
 \omega=|\partial_{p_0}\alpha_m|^{-1},\qquad
 \chi=|\det D_q(F,\vartheta)|.
\]
The bounds $\lambda_\circ I\le K\le\Lambda_\circ I$ and
$\lambda_\circ\le\chi\le\Lambda_\circ$ hold uniformly.
The explicit field
\[
 X_q=\frac{K^{-1}g}{g^{\mathsf T}K^{-1}g}
\]
is stopped at true physical and section-decision boundaries, not at
chart seams. Its flow satisfies $XF=1$, $X\vartheta=0$ and
\begin{equation}\label{eq:v83-radial-input}
 2\lambda_\circ s\le |g(\Phi_sx)|^2-|g(x)|^2
                                      \le2\Lambda_\circ s.
\end{equation}
Its endpoint-curve length is at most $C\sqrt s$. For
$I_m=1+\sum_{i=1}^{m-1}q_0^{d_i}/c_i$ one has
$|X\omega|\le C\omega I_m/|g|$ on regular survivors. At a genuine
grazing exit some internal incidence tends to zero and $\omega$ tends
to zero. Forward interior critical termination cannot occur.

Every starting point with forward lifetime at most $r$ belongs to a
layer of one of the following forms, at an actual collision or flight:
\begin{equation}\label{eq:v83-layer-input}
 c_j\le C_Eq_0^{d_j}\sqrt r,\quad
 \Delta_j\le C_Eq_0^{d'_j}\sqrt r,\quad
 \operatorname{dist}(T_R^jx,\partial Y_R^*)\le C_Eq_0^{d_j}\sqrt r.
\end{equation}
The physical coarea covariance, proved in the appendix in the convention
used here, has Jacobian $J_h=(\omega(\Phi_hx)/\omega(x))\chi(x)/\chi(\Phi_hx)$.
Writing $E_h^\varepsilon$ for all weighted starting exit mass and
\[
 O_h^\varepsilon(t)=\int_{S_h(t)} W^\varepsilon(x)
            [1-\omega(\Phi_hx)/\omega(x)]_+\,d\mu_t(x),
\]
the complete source satisfies, for each fixed $r$ and almost every $t$,
\begin{equation}\label{eq:v83-shift-input}
 b^\varepsilon(t)\le2A_rb^\varepsilon(t)+\bar E_r^\varepsilon(t)
       +\bar O_r^\varepsilon(t)
       +C_\Theta(1+\varepsilon^{-1})\sqrt r\,A_rp(t),
 \qquad A_rf(t)=\frac1{2r}\int_{t-r}^{t+r}f(s)\,ds.
\end{equation}
\end{inputstatement}
These are the endpoint-field, all-contact motion, uniform cubic-action
and separated-volume results in retained modules 174, 177, 180--182.
The guard term in \eqref{eq:v83-shift-input} is the complete product
guard, not a protective subword source. The derivative and continuation
assertions are confined to regular physical trajectories.

\begin{inputstatement}[Complete local source estimates]
\label{input:v83-local}
For every fixed auxiliary band $B_*$, constants $C_{B_*}<\infty$ and
$\rho_{B_*}<1$ give the following bounds uniformly in radius, exact
labels, interval $J$ and mark $0\le j\le m$. Let $S_{R,s}$ be the
union of strips of width $Cs$ about the finitely many coordinate lines
supporting $\partial Y_R^*$. For $s$ small,
\begin{align}
 m^2\nu\{K_m^{\rm c}=k,A_m=n,L_m\in J,
                       T_R^jx\in S_{R,s}\}
 &\le Cs^{1/16}(|J|+B_*^{-1})
   +C_{B_*}(1+B_*|J|)m^2\rho_{B_*}^{m/2},\label{eq:v83-strip-input}\\
 m^2\int_Jb^\delta(t)\,dt
 &\le C\delta^{1/16}(|J|+B_*^{-1})
   +C_{B_*}(1+B_*|J|)m^3\rho_{B_*}^{m/2}.\label{eq:v83-guard-input}
\end{align}
The local extension to a fixed remaining compact range of $\delta$
is included. At auxiliary parameter one $W^1=1$, so $b^1=p$.

Put $z=|g|$, $\alpha=129/128$ and, for $0<\eta<1/2$, define the
actual nonnegative source measures
\begin{align}
 v^{\eta,b}(t)\,dt&=L_\#(I_m z^{-\eta}\one_{\{0<z<b\}}
                                       \one_\Pi\nu_R^*),\notag\\
 v^\eta(t)\,dt&=L_\#(I_m z^{-\eta}\one_\Pi\nu_R^*).
                                                   \label{eq:v83-fractional-input}
\end{align}
They are finite and absolutely continuous at every fixed count, and
for $0<b<b_1$,
\begin{equation}\label{eq:v83-fractional-finite-input}
 m^2\int_Jv^{\eta,b}(t)\,dt
 \le C_\eta b^{\alpha-\eta}(|J|+B_*^{-1})
  +C_{B_*,\eta}(1+B_*|J|)m^2b^{1/2-\eta}\rho_{B_*}^{m/4}.
\end{equation}
Adding the bounded-gradient part gives, for every fixed $L>0$,
\begin{equation}\label{eq:v83-fixed-width-input}
 \limsup_m\sup_{R,n,k}\sup_{|J|=L}m^2\int_Jv^\eta\le C_\eta L,
 \qquad
 \limsup_m\sup_{R,n,k}\sup_{|J|=L}m^2\int_Jv^{\eta,b}
                                      \le C_\eta b^{\alpha-\eta}L.
\end{equation}
Every auxiliary band is fixed before the collision limit and enlarged
only afterwards; leading constants are independent of that band and
of the previously fixed interval length.
\end{inputstatement}
Equation~\eqref{eq:v83-strip-input} is Supplement
Theorem~\ref{supp-thm:v48-marked-strip-local}. The physical source
bound \eqref{eq:v83-guard-input} is Supplement
Theorem~\ref{supp-thm:v49-physical-remainder-local}, in particular its
finite-window estimate. The fractional bounds are the positive
layer-cake conclusions of retained module 187, stated in full above.
The joint spectral and common-space assumptions used there are recorded
in the appendix. No estimate for an untruncated first inverse gradient
or for a singular roof convolution is part of this input.

\begin{inputstatement}[The scalar positive remainder]
\label{input:v83-scalar}
For large fixed $B$, with $\varepsilon(B)=A_0B^{-1/12}$,
\begin{equation}\label{eq:v76-remainder}
 m^2p-\mathcal L=m^2b^{\varepsilon(B)}+e_B,\qquad
 \limsup_m\sup_{R,n,k}\esssup_{\mathcal C(M)}|e_B|
                                      \le C_MB^{-1/192}.
\end{equation}
The source $b^\varepsilon$ is nonnegative and has a compactly supported
$BV$ density at fixed count. No collision-uniform variation bound is
assumed. Supplement Theorem~\ref{supp-thm:v51-positive-error} gives
this remainder with the same complete finite-peak reference.
\end{inputstatement}

\subsection{Maximal segments and the two existing charges}
The retained Borel parametrization, Supplement
Theorem~\ref{supp-v85archive:thm:v85-segment-map}, is a bijection
\[
 \Psi:\coprod_j\{(\theta,u):\theta\in E_j,\ a_j(\theta)<u<b_j(\theta)\}
                      \longrightarrow\Omega,
\]
with Borel inverse. Its exact density is
$A_j(u,\theta)=\omega(\Psi_j(\theta,u))/(4\pi c_*\chi(\Psi_j(\theta,u)))$.
Every segment remains in one regular exact event until a true boundary;
Borel chart seams are not exits. Nongrazing outgoing faces retain their
incident physical component and carry
$d\zeta_+=A_j(b_j-,\theta)\,d\theta$. The full angular variation is
$|\partial_u A_j|\le C A_j I_m/z$.
These are measure identities with the null-set convention already
specified, not bounds with an unknown number of words.

The true first-exit assignment is
\[
 \mathcal F_h(t)=\int_{\Gamma_{\rm ng}^+}
  \one_{\{0<b-t\le h,\ b-t<\tau^-(y)\}}\,d\zeta_+(y),\qquad
 \overline{\mathcal F}_r=r^{-1}\int_0^r\mathcal F_h\,dh.
\]
No lower gradient cutoff is imposed. Put $l=b-a$ and define
\[
 \mathcal Z_\delta(t)=\int_{\Gamma_{\rm ng}^+}
       \one_{\{0<l\le\delta\}}\one_{\{a<t<b\}}\,d\zeta_+,
 \qquad q_\delta(t)\,dt
  =F_\#(\one_{\Omega_+^{\rm ng}}\one_{\{l\le\delta\}}\bar\nu).
\]
The second object is the original collision volume restricted to the
same actual short segments, not a volume collar for a boundary.
For source families set
\[
 \mathcal H_M(f)=\limsup_m\sup_{R,n,k}m^2
            \esssup_{t\in\mathcal C_{m,R,n,k}(M)}f(t),
 \quad \mathfrak B_M=\lim_{\delta\downarrow0}\mathcal H_M(\mathcal Z_\delta),
 \quad \mathfrak Q_M=\lim_{\delta\downarrow0}\mathcal H_M(q_\delta).
\]
These limits are monotone extended nonnegative numbers.

For $0<\eta<1/2$, with $\kappa=(1-\eta)/2$, retain
\[
 \mathsf P_{\eta,h}v^\eta(t)=\int_0^h s^{-\kappa}v^\eta(t+s)\,ds,
 \qquad\mathfrak A_{\eta,M}
       =\lim_{h\downarrow0}\mathcal H_M(\mathsf P_{\eta,h}v^\eta).
\]
Supplement Theorems~\ref{supp-v85archive:thm:v85-cutoff-surface}
and \ref{supp-v85archive:thm:v85-boundary-charge} give
\begin{equation}\label{eq:v86-surface-input}
 \mathfrak B_M\le\mathcal H_M(\overline{\mathcal F}_r)
                         \le\mathfrak B_M+C_\eta r^{1/32}.
\end{equation}
Their band sum is finite before the collision limsup at each positive
residence cutoff. Its leading constant is cutoff independent; the
complete zero-residence tail is still inside $\mathcal H_M$ before
its cutoff tends to zero. The active offset correction to historical
module 189 remains Supplement equation
\eqref{supp-v83archive:eq:v83-dyadic-tail}; the historical inclusion is
explicitly marked superseded. None of these inputs asserts charge zero.


\subsection{The precise finite exit and surface inputs}
For total residence $\ell=\tau^-+\tau^+$, let
$e_s(t)=\int\one_{\{\tau^+\le s\}}\,d\mu_t$.
The finite-count statement of Supplement
Theorem~\ref{supp-v83archive:thm:v83-exit-local-mass} is
\begin{equation}\label{eq:v88-exit-input}
 m^2\int_J e_s\le C s^{1/32}(|J|+B_*^{-1})
       +C_{B_*}(1+B_*|J|)m^3\rho_{B_*}^{m/2}.
\end{equation}
It holds for all sufficiently small $s$, every interval $J$ and all
finite counts, with one fixed auxiliary band. The proof uses the
actual layers in \eqref{eq:v83-layer-input}, the marked section
estimate and the complete original guard, not a count of words.
The one-flight envelope used in its geometric input is described in
Supplement Lemma~\ref{supp-lem:v49-physical-envelope}; its explicit
coordinates also give the physical area cap proved below.

The actual untruncated surface and original weighted defects satisfy
\begin{equation}\label{eq:v88-optical-input}
 E_h^\varepsilon+O_h^\varepsilon
     \le C\mathcal F_h+C_\eta\mathsf P_{\eta,h}v^\eta.
\end{equation}
This is Supplement Theorem~\ref{supp-v83archive:thm:v83-two-sided-flux},
with the same joint physical coarea kernels and initial weights.
For the part $\mathcal L_r$ of $\overline{\mathcal F}_r$ having
backward residence $\ell>2r$, Supplement
Theorem~\ref{supp-v84archive:thm:v84-long-height} proves
\begin{equation}\label{eq:v88-long-input}
 \limsup_m\sup_{R,n,k}m^2\|\mathcal L_r\|_\infty\le Cr^{1/32}.
\end{equation}
Its persistence proof precedes the mass-to-height deduction. We do
not extend \eqref{eq:v88-long-input} to zero residence. The short
part is estimated below in a different, explicitly finite norm.

All fiber identities have the following convention. For fixed
$m,R,n,k$, the Borel maximal-segment parametrization fixes common
kernels outside a null set in the roof, and common product kernels
outside a joint null set for roof and time. Countably many widths
share versions; the uniform order is $\sup_R\operatorname*{ess\,sup}_t$,
not a claim to one null set for all radii. All positive restrictions
and infinite moment integrals below use these kernels.

\input{core/205_capped_negative_residence_moments}
\input{core/206_uncut_surface_lorentz_control}
\input{core/207_exact_label_raw_density_local_norm}

\section{The unchanged pointwise endpoint}
The new negative residence theorem is a finite-count assertion for an
actual singular insertion. Its full layer cake includes arbitrarily
short physical segments before the collision limit. The outgoing
trace theorem similarly sums all incident components before applying
a norm. Thus the complete surface weak-norm decay and raw local-norm
limit proved here require no positive residence cutoff.

These results do not identify a finite local norm with essential
height. The endpoint in \eqref{eq:v76-target} continues to require
\[
 \lim_{B\to\infty}\mathfrak R_{\varepsilon(B),M}=0,
\]
where $\mathfrak R$ is the original weighted-source concentration
charge of Supplement Theorem~\ref{supp-v87archive:thm:v87-endpoint}.
The local-norm theorem constrains such concentration but does not
exclude narrow peaks in its essential supremum. The older sufficient
pairs $(\mathfrak A,\mathfrak B)$ and $(\mathfrak A,\mathfrak Q)$
remain valid with their exact hypotheses in the complete archive.
No charge-zero assertion is inferred here.

In particular, a strong local $L^p$ theorem for $p<32/31$ is not a
uniform denominator theorem on positive arithmetic classes. The
scalar-first path remainder, collision-to-return clock transfer and
forward likelihood statements are preserved, not promoted to new
unconditional same-roof ratios. Zero arithmetic classes are never
normalized. All previous mathematical files and the complete preceding
article remain in the supplement and provenance.


\endgroup
\part{Complete segment traces and the original concentration source}
\begingroup
\let\ArchiveLabel\label
\let\ArchiveRef\ref
\renewcommand{\label}[1]{\ArchiveLabel{v87archive:#1}}
\renewcommand{\ref}[1]{%
  \IfBeginWith{#1}{supp-}{%
    \StrBehind{#1}{supp-}[\ArchiveTarget]%
    \ArchiveRef{\ArchiveTarget}%
  }{\ArchiveRef{v87archive:#1}}%
}
\makeatletter
\let\ltx@label\label
\makeatother
\let\ArchivePageRef\pageref
\renewcommand{\pageref}[1]{%
  \IfBeginWith{#1}{supp-}{%
    \StrBehind{#1}{supp-}[\ArchivePageTarget]%
    \ArchivePageRef{\ArchivePageTarget}%
  }{\ArchivePageRef{v87archive:#1}}%
}
\paragraph{Retained revision-87 abstract.}

We study the unchanged exact-label pointwise local limit problem for
actual returns in a triangular Lorentz gas. A complete optical moment
strictly above one yields a trace estimate for the physical angular
density on maximal regular segments. Summation uses an exact Borel
parametrization and has multiplicity one. The trace estimate gives a
finite-count optical envelope on all initial states whose gradient and
total residence exceed explicit powers of the microscopic roof width.
The envelope includes every optical sensitivity and every positive
incidence; for one choice its normalized bound is $K h^{31/1408}$,
uniformly before the collision limit. Inserting this estimate into the
original first-exit argument localizes the full bad-source height to
initial criticality and zero total residence. The organizing pointwise
error is exactly the outer limit of this original weighted-source
concentration charge. Its vanishing remains to be established; no
surface cutoff or arithmetic class is removed from the theorem.



\section{The unchanged physical problem}
Put $R_-=9/20$ and $R_+=47/100$, and let $R\in I=[R_-,R_+]$. Circular scatterers have centers in
$\Lambda=\Z(1,0)+\Z(1/2,\sqrt3/2)$. Write $T_R$ for the regular
collision map and $(\alpha,p)\in(\R/2\pi\Z)\times(-1,1)$ for its
coordinates; $p$ is the outgoing tangential velocity and
$\dd\nu=\dd\alpha\,\dd p/(4\pi)$. Physical flights satisfy
$\tau_-\le\tau_R\le\tau_+$, where $\tau_-=3/50$ and $\tau_+$
is the fixed finite-horizon bound.

The section is specified here to fix every event in the argument.
Put $h_0=1/200$. Let $Y_R$ be the union of the four open coordinate
squares of half-side $h_0$ centered at
\[
 (0,0),\quad(\pi/6,-1/2),\quad
 (-\theta_R,\sin\theta_R),\quad
 (\pi/3-\theta_R,\sin\theta_R).
\]
Here $\theta_R=2\arctan q_R$, and $q_R$ is the unique root in
$(49/100,53/100)$ of
\[
 8Rq(1+q^2)-4q(1-q^2)+\sqrt3(1-6q^2+q^4)=0.
\]
Set
\[
 U=\{|\alpha-\pi/6|<3/50,\ |p|<3/20\},\qquad
 Y_R^*=Y_R\cup U,\qquad c_*=\frac{91}{10000\pi}.
\]
The rectangles are disjoint and $\nu(Y_R^*)=c_*$. We always use
$\nu_R^*=c_*^{-1}\nu|_{Y_R^*}$, not a renormalization of a smaller
source.

At $m$ flights, let $L=L_{m,R}$ be their total length,
$K^{\rm c}_{m,R}$ their lattice displacement, and
$A_{m,R}=\sum_{j=0}^{m-1}\one_{Y_R^*}\circ T_R^j$. The event is
\begin{equation}\label{eq:v76-event}
 \Pi_{n,k,m,R}=Y_R^*\cap\{K^{\rm c}_{m,R}=k,\ A_{m,R}=n\}
                         \cap T_R^{-m}Y_R^*.
\end{equation}
The initial visit is counted; the terminal visit is tested separately.
On this event the actual $n$th return has collision count $m$ and
roof $L$. Denote its density by $p_{n,R}(k,m,t)$.

Write $c_j=\sqrt{1-p_j^2}$ and let $\Delta_j$ be the least clearance
of flight $j$ from any third disk, excluding its two incident disks.
Keep the original smooth function $0\le\Theta\le1$, zero on
$(-\infty,1]$ and one on $[2,\infty)$. With $q_0=47/53$,
$d_j=\min(j,m-j)$ and $d'_j=\min(j,m-1-j)$, set
\begin{align}
 H^\varepsilon&=
 \prod_{j=0}^{m}\Theta\!\left(\frac{c_j}{\varepsilon q_0^{d_j/4}}\right)
 \prod_{j=0}^{m-1}\Theta\!\left(\frac{\Delta_j}{\varepsilon q_0^{d'_j/4}}\right),
 \qquad W^\varepsilon=1-H^\varepsilon,\label{eq:v76-guard}\\
 b^\varepsilon(t)\dd t
   &=L_\#(W^\varepsilon\one_{\Pi_{n,k,m,R}}\nu_R^*).
 \label{eq:v76-source}
\end{align}
In particular first incidence and first clearance outside any selected
critical neighborhoods are still present. No section-distance factor
has been inserted into this guard.

\subsection{The arithmetic reference and the target}
For $y=(k_1,k_2,t,n)$ put
\[
 \mu_R=(0,0,\bar\tau_R,c_*),\quad
 Z=(y-m\mu_R)/\sqrt m,\quad
 \mathcal C_{m,R,n,k}(M)=\{t:|Z|\le M\}.
\]
We retain the finite spectral-peak reference, including peaks close to
but not on the unit circle. To specify its meaning, the collision
record observable is $(\kappa_1,\kappa_2,\tau,\one_{Y_R^*})$.
Let $\lambda_j(R,\vartheta)$ and $\mathsf P_j(R,\vartheta)$ be its
continued simple twisted eigenvalue and projection in each of the
fixed peak charts, and let $\chi_j$ be their frequency partition.
At the real maximum $a_j$ of $\log|\lambda_j|$ put
\[
 z_j=\lambda_j(R,a_j),\quad
 \mu_j=i^{-1}\nabla\log\lambda_j(R,a_j),\quad
 H_j=-\nabla^2\log\lambda_j(R,a_j),
\]
with $\operatorname{Re}H_j\ge aI_4$ for a common $a>0$. If
$\eta_R=\one_{Y_R^*}$ and $\ell$ denotes mass evaluation, put
$B_j=\chi_j(R,a_j)\ell(M_{\eta_R}\mathsf P_j(R,a_j)(\eta_R\nu))$.
Define
\begin{align}
 G_H(z)&=(2\pi)^{-4}\int_{\R^4}e^{-iv\cdot z-v^{\mathsf T}Hv/2}\dd v,
 \notag\\
 \mathcal L_{m,R}(y)&=\operatorname{Re}\left[
 c_*^{-1}\sum_j e^{-ia_j\cdot y}z_j^m B_j
      G_{H_j}\left((y-m\mu_j)/\sqrt m\right)\right].
 \label{eq:v76-reference}
\end{align}
Only the charts active at $R$ are summed. This is the same transition
reference as in the retained spectral theorem. The existence and
uniform estimates for these scalar data belong to that analytic
input, not to the transport argument below. No frequency branch or
arithmetic zero class is deleted.

The governing assertion remains
\begin{equation}\label{eq:v76-target}
 \sup_{R,\,1\le n\le m,\,k\in\Z^2}
 \esssup_{t\in\mathcal C_{m,R,n,k}(M)}
 |m^2p_{n,R}(k,m,t)-\mathcal L_{m,R}(k_1,k_2,t,n)|\longrightarrow0.
\end{equation}
Let $\mathcal E_M$ denote this error's collision limsup. An empty
central target contributes zero. The order of the uniformity is
$\sup_R\esssup_t$, so exceptional density sets may depend on $R$.

\subsection{Physical traces and a uniform microscopic envelope}
The sensitivity is the original additive function
$I_m=1+\sum_{i=1}^{m-1}q_0^{d_i}/c_i$. Write $z=|g|$ for the
endpoint-gradient norm and $\ell$ for the total residence of its
maximal regular exact-event segment. We use the common Borel coarea
kernels $\mu_t$. At fixed $m,R,n,k$ the identities hold outside a
null set in $t$, or a joint null set for product kernels; the uniform
norm is $\sup_R\esssup_t$, not a claim to one null set for all radii.

\begin{theorem}[Microscopic envelope on a shrinking initial family]
\label{thm:v87-main}
Set $a(h)=h^{1/128}$ and $\delta(h)=h^{1/64}$. On the original
reachable target fibers define
\[
 V_h^{\rm reg}(t)=\int_0^h\int
  \one_{\{\tau^-(y)>s,\ z(\Phi_{-s}y)\ge a(h),\ \ell(y)\ge\delta(h)\}}
                      \frac{I_m(y)}{z(y)}\,d\mu_{t+s}(y)\,ds.
\]
Under the explicitly retained inputs, there are fixed $K<\infty$
and $h_*>0$ such that
\begin{equation}\label{eq:v87-main-envelope}
 \sup_{m\ge2,R,n,k}m^2\|V_h^{\rm reg}\|_\infty
                              \le Kh^{31/1408}\qquad(0<h<h_*).
\end{equation}
This includes all positive incidences and all values of $I_m$ on
the specified initial family. Both thresholds depend on the roof
width, not on the collision count, and the estimate holds uniformly
before any collision limit.
\end{theorem}
The proof is Theorems~\ref{thm:v87-summed-trace}--\ref{thm:v87-moving-envelope}.
It uses the actual moment $I_m^{33/32}$, obtained from the retained
mixed optical source. The elementary trace inequality is proved in
full, including its integration over complete physical segments.

The complementary initial mass is not discarded. Define
\[
 R_h^\varepsilon(t)=\int W^\varepsilon
            \one_{\{z<a(h)\ \mathrm{or}\ \ell<\delta(h)\}}\,d\mu_t.
\]
Theorem~\ref{thm:v87-endpoint} identifies the original pointwise error
with the outer limit of the ordered microscopic height of this source.
The result does not assert that this last charge is zero, nor that
previous sufficient optical charges have vanished. The original
target, arithmetic reference and scalar-first downstream order remain
unchanged.

\section{Explicit analytic and geometric inputs}
We collect the indispensable inherited assertions in the exact form
used by the new proofs. Their source theorems and operator dependencies
are specified here and in Appendix~\ref{app:v83-kernels}. The new
potential and first-exit estimates do not assume a surface spectral
insertion theorem or an unweighted pointwise local bound.

\begin{inputstatement}[Physical flow and source comparison]
\label{input:v83-geometric}
On the regular part $\Omega$ of each exact event, endpoint angles
$q=(\alpha_0,\alpha_m)$ give $F=L$, $g=\nabla_qF=R(-p_0,p_m)$,
$K=D_q^2F$ and $\vartheta=\arg g$. Put
\[
 \lambda_\circ=9\sqrt{19}/200,\qquad \Lambda_\circ=47/6,\qquad
 \omega=|\partial_{p_0}\alpha_m|^{-1},\qquad
 \chi=|\det D_q(F,\vartheta)|.
\]
The bounds $\lambda_\circ I\le K\le\Lambda_\circ I$ and
$\lambda_\circ\le\chi\le\Lambda_\circ$ hold uniformly.
The explicit field
\[
 X_q=\frac{K^{-1}g}{g^{\mathsf T}K^{-1}g}
\]
is stopped at true physical and section-decision boundaries, not at
chart seams. Its flow satisfies $XF=1$, $X\vartheta=0$ and
\begin{equation}\label{eq:v83-radial-input}
 2\lambda_\circ s\le |g(\Phi_sx)|^2-|g(x)|^2
                                      \le2\Lambda_\circ s.
\end{equation}
Its endpoint-curve length is at most $C\sqrt s$. For
$I_m=1+\sum_{i=1}^{m-1}q_0^{d_i}/c_i$ one has
$|X\omega|\le C\omega I_m/|g|$ on regular survivors. At a genuine
grazing exit some internal incidence tends to zero and $\omega$ tends
to zero. Forward interior critical termination cannot occur.

Every starting point with forward lifetime at most $r$ belongs to a
layer of one of the following forms, at an actual collision or flight:
\begin{equation}\label{eq:v83-layer-input}
 c_j\le C_Eq_0^{d_j}\sqrt r,\quad
 \Delta_j\le C_Eq_0^{d'_j}\sqrt r,\quad
 \operatorname{dist}(T_R^jx,\partial Y_R^*)\le C_Eq_0^{d_j}\sqrt r.
\end{equation}
The physical coarea covariance, proved in the appendix in the convention
used here, has Jacobian $J_h=(\omega(\Phi_hx)/\omega(x))\chi(x)/\chi(\Phi_hx)$.
Writing $E_h^\varepsilon$ for all weighted starting exit mass and
\[
 O_h^\varepsilon(t)=\int_{S_h(t)} W^\varepsilon(x)
            [1-\omega(\Phi_hx)/\omega(x)]_+\,d\mu_t(x),
\]
the complete source satisfies, for each fixed $r$ and almost every $t$,
\begin{equation}\label{eq:v83-shift-input}
 b^\varepsilon(t)\le2A_rb^\varepsilon(t)+\bar E_r^\varepsilon(t)
       +\bar O_r^\varepsilon(t)
       +C_\Theta(1+\varepsilon^{-1})\sqrt r\,A_rp(t),
 \qquad A_rf(t)=\frac1{2r}\int_{t-r}^{t+r}f(s)\,ds.
\end{equation}
\end{inputstatement}
These are the endpoint-field, all-contact motion, uniform cubic-action
and separated-volume results in retained modules 174, 177, 180--182.
The guard term in \eqref{eq:v83-shift-input} is the complete product
guard, not a protective subword source. The derivative and continuation
assertions are confined to regular physical trajectories.

\begin{inputstatement}[Complete local source estimates]
\label{input:v83-local}
For every fixed auxiliary band $B_*$, constants $C_{B_*}<\infty$ and
$\rho_{B_*}<1$ give the following bounds uniformly in radius, exact
labels, interval $J$ and mark $0\le j\le m$. Let $S_{R,s}$ be the
union of strips of width $Cs$ about the finitely many coordinate lines
supporting $\partial Y_R^*$. For $s$ small,
\begin{align}
 m^2\nu\{K_m^{\rm c}=k,A_m=n,L_m\in J,
                       T_R^jx\in S_{R,s}\}
 &\le Cs^{1/16}(|J|+B_*^{-1})
   +C_{B_*}(1+B_*|J|)m^2\rho_{B_*}^{m/2},\label{eq:v83-strip-input}\\
 m^2\int_Jb^\delta(t)\,dt
 &\le C\delta^{1/16}(|J|+B_*^{-1})
   +C_{B_*}(1+B_*|J|)m^3\rho_{B_*}^{m/2}.\label{eq:v83-guard-input}
\end{align}
The local extension to a fixed remaining compact range of $\delta$
is included. At auxiliary parameter one $W^1=1$, so $b^1=p$.

Put $z=|g|$, $\alpha=129/128$ and, for $0<\eta<1/2$, define the
actual nonnegative source measures
\begin{align}
 v^{\eta,b}(t)\,dt&=L_\#(I_m z^{-\eta}\one_{\{0<z<b\}}
                                       \one_\Pi\nu_R^*),\notag\\
 v^\eta(t)\,dt&=L_\#(I_m z^{-\eta}\one_\Pi\nu_R^*).
                                                   \label{eq:v83-fractional-input}
\end{align}
They are finite and absolutely continuous at every fixed count, and
for $0<b<b_1$,
\begin{equation}\label{eq:v83-fractional-finite-input}
 m^2\int_Jv^{\eta,b}(t)\,dt
 \le C_\eta b^{\alpha-\eta}(|J|+B_*^{-1})
  +C_{B_*,\eta}(1+B_*|J|)m^2b^{1/2-\eta}\rho_{B_*}^{m/4}.
\end{equation}
Adding the bounded-gradient part gives, for every fixed $L>0$,
\begin{equation}\label{eq:v83-fixed-width-input}
 \limsup_m\sup_{R,n,k}\sup_{|J|=L}m^2\int_Jv^\eta\le C_\eta L,
 \qquad
 \limsup_m\sup_{R,n,k}\sup_{|J|=L}m^2\int_Jv^{\eta,b}
                                      \le C_\eta b^{\alpha-\eta}L.
\end{equation}
Every auxiliary band is fixed before the collision limit and enlarged
only afterwards; leading constants are independent of that band and
of the previously fixed interval length.
\end{inputstatement}
Equation~\eqref{eq:v83-strip-input} is Supplement
Theorem~\ref{supp-thm:v48-marked-strip-local}. The physical source
bound \eqref{eq:v83-guard-input} is Supplement
Theorem~\ref{supp-thm:v49-physical-remainder-local}, in particular its
finite-window estimate. The fractional bounds are the positive
layer-cake conclusions of retained module 187, stated in full above.
The joint spectral and common-space assumptions used there are recorded
in the appendix. No estimate for an untruncated first inverse gradient
or for a singular roof convolution is part of this input.

\begin{inputstatement}[The scalar positive remainder]
\label{input:v83-scalar}
For large fixed $B$, with $\varepsilon(B)=A_0B^{-1/12}$,
\begin{equation}\label{eq:v76-remainder}
 m^2p-\mathcal L=m^2b^{\varepsilon(B)}+e_B,\qquad
 \limsup_m\sup_{R,n,k}\esssup_{\mathcal C(M)}|e_B|
                                      \le C_MB^{-1/192}.
\end{equation}
The source $b^\varepsilon$ is nonnegative and has a compactly supported
$BV$ density at fixed count. No collision-uniform variation bound is
assumed. Supplement Theorem~\ref{supp-thm:v51-positive-error} gives
this remainder with the same complete finite-peak reference.
\end{inputstatement}

\subsection{Maximal segments and the two existing charges}
The retained Borel parametrization, Supplement
Theorem~\ref{supp-v85archive:thm:v85-segment-map}, is a bijection
\[
 \Psi:\coprod_j\{(\theta,u):\theta\in E_j,\ a_j(\theta)<u<b_j(\theta)\}
                      \longrightarrow\Omega,
\]
with Borel inverse. Its exact density is
$A_j(u,\theta)=\omega(\Psi_j(\theta,u))/(4\pi c_*\chi(\Psi_j(\theta,u)))$.
Every segment remains in one regular exact event until a true boundary;
Borel chart seams are not exits. Nongrazing outgoing faces retain their
incident physical component and carry
$d\zeta_+=A_j(b_j-,\theta)\,d\theta$. The full angular variation is
$|\partial_u A_j|\le C A_j I_m/z$.
These are measure identities with the null-set convention already
specified, not bounds with an unknown number of words.

The true first-exit assignment is
\[
 \mathcal F_h(t)=\int_{\Gamma_{\rm ng}^+}
  \one_{\{0<b-t\le h,\ b-t<\tau^-(y)\}}\,d\zeta_+(y),\qquad
 \overline{\mathcal F}_r=r^{-1}\int_0^r\mathcal F_h\,dh.
\]
No lower gradient cutoff is imposed. Put $l=b-a$ and define
\[
 \mathcal Z_\delta(t)=\int_{\Gamma_{\rm ng}^+}
       \one_{\{0<l\le\delta\}}\one_{\{a<t<b\}}\,d\zeta_+,
 \qquad q_\delta(t)\,dt
  =F_\#(\one_{\Omega_+^{\rm ng}}\one_{\{l\le\delta\}}\bar\nu).
\]
The second object is the original collision volume restricted to the
same actual short segments, not a volume collar for a boundary.
For source families set
\[
 \mathcal H_M(f)=\limsup_m\sup_{R,n,k}m^2
            \esssup_{t\in\mathcal C_{m,R,n,k}(M)}f(t),
 \quad \mathfrak B_M=\lim_{\delta\downarrow0}\mathcal H_M(\mathcal Z_\delta),
 \quad \mathfrak Q_M=\lim_{\delta\downarrow0}\mathcal H_M(q_\delta).
\]
These limits are monotone extended nonnegative numbers.

For $0<\eta<1/2$, with $\kappa=(1-\eta)/2$, retain
\[
 \mathsf P_{\eta,h}v^\eta(t)=\int_0^h s^{-\kappa}v^\eta(t+s)\,ds,
 \qquad\mathfrak A_{\eta,M}
       =\lim_{h\downarrow0}\mathcal H_M(\mathsf P_{\eta,h}v^\eta).
\]
Supplement Theorems~\ref{supp-v85archive:thm:v85-cutoff-surface}
and \ref{supp-v85archive:thm:v85-boundary-charge} give
\begin{equation}\label{eq:v86-surface-input}
 \mathfrak B_M\le\mathcal H_M(\overline{\mathcal F}_r)
                         \le\mathfrak B_M+C_\eta r^{1/32}.
\end{equation}
Their band sum is finite before the collision limsup at each positive
residence cutoff. Its leading constant is cutoff independent; the
complete zero-residence tail is still inside $\mathcal H_M$ before
its cutoff tends to zero. The active offset correction to historical
module 189 remains Supplement equation
\eqref{supp-v83archive:eq:v83-dyadic-tail}; the historical inclusion is
explicitly marked superseded. None of these inputs asserts charge zero.

\input{core/199_complete_mixed_optical_moments}
\input{core/202_complete_segment_trace_estimate}
\input{core/203_shrinking_segment_optical_envelope}
\input{core/204_original_source_concentration_localization}

\section{The original endpoint and its remaining physical concentration}
The microscopic estimate \eqref{eq:v87-uniform-envelope} is uniform
at finite collision count. It does not restrict optical sensitivity,
and it keeps the complete source family defined by the two shrinking
initial thresholds. Its finite-count spectral remainder is retained
in \eqref{eq:v87-finite-drift}; the positive exponent makes the width
limit legitimate without changing the auxiliary band.

The full pointwise problem has not been replaced by that subfamily.
The exceptional initial mass is retained with the original weight
$W^\varepsilon$, and the complete positive-cutoff outgoing source
is paid by its independently proved residence theorem. The exact
remaining assertion is
\[
 \lim_{B\to\infty}\mathfrak R_{\varepsilon(B),M}=0
 \quad\text{for every }M.
\]
It concerns original source concentration near initial criticality or
zero total regular residence, including every true endpoint type.
The trace theorem does not estimate this last charge. It is not marked
zero in the present revision.

Once the scalar endpoint is proved, Supplement
Theorem~\ref{supp-thm:v73-scalar-path-equality} gives the same-roof
path numerator. Normalization remains restricted to positive arithmetic
reference classes. The actual-return clock and forward likelihood
retain their separate inherited hypotheses; zero classes are never
divided by. The complete preceding article and every inherited proof
are retained in the supplement, rather than deleted from the record.

\appendix
\input{appendices/v87_common_kernels_and_dependencies}

\endgroup
\part{Finite-count mixed moments and reachable optical concentration}
\begingroup
\let\ArchiveLabel\label
\let\ArchiveRef\ref
\renewcommand{\label}[1]{\ArchiveLabel{v86archive:#1}}
\renewcommand{\ref}[1]{%
  \IfBeginWith{#1}{supp-}{%
    \StrBehind{#1}{supp-}[\ArchiveTarget]%
    \ArchiveRef{\ArchiveTarget}%
  }{\ArchiveRef{v86archive:#1}}%
}
\makeatletter
\let\ltx@label\label
\makeatother
\let\ArchivePageRef\pageref
\renewcommand{\pageref}[1]{%
  \IfBeginWith{#1}{supp-}{%
    \StrBehind{#1}{supp-}[\ArchivePageTarget]%
    \ArchivePageRef{\ArchivePageTarget}%
  }{\ArchivePageRef{v86archive:#1}}%
}
\paragraph{Retained revision-86 abstract.}

We study the unchanged exact-label pointwise local limit problem for
actual returns in a triangular Lorentz gas. A physically capped joint
endpoint--incidence estimate gives complete mixed optical moments strictly
above the first power, with a finite-count remainder integrable at every
incidence and gradient scale. We use this additional moment to prove a
uniform essential-height envelope for the entire microscopic sector whose
optical sensitivity exceeds a prescribed inverse power of the roof offset.
For explicit exponents its finite-count bound is $K h^{1/8}$ and its
ordered height is $C h^{9/8}$. The infinite dyadic tail is summed before
the collision limit. We also retain the exact backward-lifetime selection
in the optical drift, so that only target states actually reachable from
the starting roof contribute. This gives a smaller sufficient optical
charge and removes the estimated high-sensitivity sector from it without
changing the physical source. The original pointwise theorem is reduced
to reachable moderate-sensitivity concentration and the original
zero-residence volume concentration; their vanishing remains to be proved.



\section{The unchanged physical problem}
Put $R_-=9/20$ and $R_+=47/100$, and let $R\in I=[R_-,R_+]$. Circular scatterers have centers in
$\Lambda=\Z(1,0)+\Z(1/2,\sqrt3/2)$. Write $T_R$ for the regular
collision map and $(\alpha,p)\in(\R/2\pi\Z)\times(-1,1)$ for its
coordinates; $p$ is the outgoing tangential velocity and
$\dd\nu=\dd\alpha\,\dd p/(4\pi)$. Physical flights satisfy
$\tau_-\le\tau_R\le\tau_+$, where $\tau_-=3/50$ and $\tau_+$
is the fixed finite-horizon bound.

The section is specified here to fix every event in the argument.
Put $h_0=1/200$. Let $Y_R$ be the union of the four open coordinate
squares of half-side $h_0$ centered at
\[
 (0,0),\quad(\pi/6,-1/2),\quad
 (-\theta_R,\sin\theta_R),\quad
 (\pi/3-\theta_R,\sin\theta_R).
\]
Here $\theta_R=2\arctan q_R$, and $q_R$ is the unique root in
$(49/100,53/100)$ of
\[
 8Rq(1+q^2)-4q(1-q^2)+\sqrt3(1-6q^2+q^4)=0.
\]
Set
\[
 U=\{|\alpha-\pi/6|<3/50,\ |p|<3/20\},\qquad
 Y_R^*=Y_R\cup U,\qquad c_*=\frac{91}{10000\pi}.
\]
The rectangles are disjoint and $\nu(Y_R^*)=c_*$. We always use
$\nu_R^*=c_*^{-1}\nu|_{Y_R^*}$, not a renormalization of a smaller
source.

At $m$ flights, let $L=L_{m,R}$ be their total length,
$K^{\rm c}_{m,R}$ their lattice displacement, and
$A_{m,R}=\sum_{j=0}^{m-1}\one_{Y_R^*}\circ T_R^j$. The event is
\begin{equation}\label{eq:v76-event}
 \Pi_{n,k,m,R}=Y_R^*\cap\{K^{\rm c}_{m,R}=k,\ A_{m,R}=n\}
                         \cap T_R^{-m}Y_R^*.
\end{equation}
The initial visit is counted; the terminal visit is tested separately.
On this event the actual $n$th return has collision count $m$ and
roof $L$. Denote its density by $p_{n,R}(k,m,t)$.

Write $c_j=\sqrt{1-p_j^2}$ and let $\Delta_j$ be the least clearance
of flight $j$ from any third disk, excluding its two incident disks.
Keep the original smooth function $0\le\Theta\le1$, zero on
$(-\infty,1]$ and one on $[2,\infty)$. With $q_0=47/53$,
$d_j=\min(j,m-j)$ and $d'_j=\min(j,m-1-j)$, set
\begin{align}
 H^\varepsilon&=
 \prod_{j=0}^{m}\Theta\!\left(\frac{c_j}{\varepsilon q_0^{d_j/4}}\right)
 \prod_{j=0}^{m-1}\Theta\!\left(\frac{\Delta_j}{\varepsilon q_0^{d'_j/4}}\right),
 \qquad W^\varepsilon=1-H^\varepsilon,\label{eq:v76-guard}\\
 b^\varepsilon(t)\dd t
   &=L_\#(W^\varepsilon\one_{\Pi_{n,k,m,R}}\nu_R^*).
 \label{eq:v76-source}
\end{align}
In particular first incidence and first clearance outside any selected
critical neighborhoods are still present. No section-distance factor
has been inserted into this guard.

\subsection{The arithmetic reference and the target}
For $y=(k_1,k_2,t,n)$ put
\[
 \mu_R=(0,0,\bar\tau_R,c_*),\quad
 Z=(y-m\mu_R)/\sqrt m,\quad
 \mathcal C_{m,R,n,k}(M)=\{t:|Z|\le M\}.
\]
We retain the finite spectral-peak reference, including peaks close to
but not on the unit circle. To specify its meaning, the collision
record observable is $(\kappa_1,\kappa_2,\tau,\one_{Y_R^*})$.
Let $\lambda_j(R,\vartheta)$ and $\mathsf P_j(R,\vartheta)$ be its
continued simple twisted eigenvalue and projection in each of the
fixed peak charts, and let $\chi_j$ be their frequency partition.
At the real maximum $a_j$ of $\log|\lambda_j|$ put
\[
 z_j=\lambda_j(R,a_j),\quad
 \mu_j=i^{-1}\nabla\log\lambda_j(R,a_j),\quad
 H_j=-\nabla^2\log\lambda_j(R,a_j),
\]
with $\operatorname{Re}H_j\ge aI_4$ for a common $a>0$. If
$\eta_R=\one_{Y_R^*}$ and $\ell$ denotes mass evaluation, put
$B_j=\chi_j(R,a_j)\ell(M_{\eta_R}\mathsf P_j(R,a_j)(\eta_R\nu))$.
Define
\begin{align}
 G_H(z)&=(2\pi)^{-4}\int_{\R^4}e^{-iv\cdot z-v^{\mathsf T}Hv/2}\dd v,
 \notag\\
 \mathcal L_{m,R}(y)&=\operatorname{Re}\left[
 c_*^{-1}\sum_j e^{-ia_j\cdot y}z_j^m B_j
      G_{H_j}\left((y-m\mu_j)/\sqrt m\right)\right].
 \label{eq:v76-reference}
\end{align}
Only the charts active at $R$ are summed. This is the same transition
reference as in the retained spectral theorem. The existence and
uniform estimates for these scalar data belong to that analytic
input, not to the transport argument below. No frequency branch or
arithmetic zero class is deleted.

The governing assertion remains
\begin{equation}\label{eq:v76-target}
 \sup_{R,\,1\le n\le m,\,k\in\Z^2}
 \esssup_{t\in\mathcal C_{m,R,n,k}(M)}
 |m^2p_{n,R}(k,m,t)-\mathcal L_{m,R}(k_1,k_2,t,n)|\longrightarrow0.
\end{equation}
Let $\mathcal E_M$ denote this error's collision limsup. An empty
central target contributes zero. The order of the uniformity is
$\sup_R\esssup_t$, so exceptional density sets may depend on $R$.

\subsection{A finite-count microscopic envelope}
The optical sensitivity is the original additive function
$I_m=1+\sum_{i=1}^{m-1}q_0^{d_i}/c_i$. Put $z=|g|$ on the
regular exact event, using the endpoint gradient specified below.
At the first use of a fiber expression we fix the common Borel kernels
$\mu_t$ of that event. For each fixed $m,R,n,k$, measure identities
hold outside a null set in $t$, and product identities outside a joint
null set in $(t,s)$. Countable selectors and dyadic scales use these
same kernels. We do not require one null set for all radii.

\begin{theorem}[Complete microscopic high-sensitivity control]
\label{thm:v86-main}
Under the explicitly retained spectral inputs below, define
\[
 D_h(t)=\int_0^h s^{-3/8}\int I_m(y)z(y)^{-1/4}
                     \one_{\{I_m(y)>s^{-16}\}}\,d\mu_{t+s}(y)\,ds.
\]
For each fixed auxiliary band $B_*$ there is a finite $K_{B_*}$ such
that, for every $0<h<h_0$,
\begin{align}
 \sup_{m\ge2,R,n,k}m^2\|D_h\|_\infty&\le K_{B_*}h^{1/8},
                                          \label{eq:v86-main-uniform}\\
 \limsup_m\sup_{R,n,k}m^2\|D_h\|_\infty&\le Ch^{9/8}.
                                          \label{eq:v86-main-ordered}
\end{align}
The first bound is uniform before the collision limit. The second
follows after that limit and enlargement of the fixed auxiliary band.
Both include all actual words, exact labels and all microscopic
scales down to zero in the stated sector.
\end{theorem}
Theorems~\ref{thm:v86-moments} and \ref{thm:v86-envelope} prove this
statement. The extra moment is $I_m^{33/32}z^{-1/4}$, not an assumed
bounded multiplier on an anisotropic space. The exponent $1/8$ makes
the complete dyadic tail summable together with its finite-count error.

The theorem estimates a positive part of the original optical source;
it does not replace that source by a protected one. Its complement
remains explicit. Theorem~\ref{thm:v86-endpoint} also retains the
backward-lifetime selector in every optical comparison, replacing the
unrestricted fractional majorant only by a proved reachable-fiber
bound. The original pointwise target \eqref{eq:v76-target}, and the
scalar-first order of its downstream consequences, are unchanged.

\section{Explicit analytic and geometric inputs}
We collect the indispensable inherited assertions in the exact form
used by the new proofs. Their source theorems and operator dependencies
are specified here and in Appendix~\ref{app:v83-kernels}. The new
potential and first-exit estimates do not assume a surface spectral
insertion theorem or an unweighted pointwise local bound.

\begin{inputstatement}[Physical flow and source comparison]
\label{input:v83-geometric}
On the regular part $\Omega$ of each exact event, endpoint angles
$q=(\alpha_0,\alpha_m)$ give $F=L$, $g=\nabla_qF=R(-p_0,p_m)$,
$K=D_q^2F$ and $\vartheta=\arg g$. Put
\[
 \lambda_\circ=9\sqrt{19}/200,\qquad \Lambda_\circ=47/6,\qquad
 \omega=|\partial_{p_0}\alpha_m|^{-1},\qquad
 \chi=|\det D_q(F,\vartheta)|.
\]
The bounds $\lambda_\circ I\le K\le\Lambda_\circ I$ and
$\lambda_\circ\le\chi\le\Lambda_\circ$ hold uniformly.
The explicit field
\[
 X_q=\frac{K^{-1}g}{g^{\mathsf T}K^{-1}g}
\]
is stopped at true physical and section-decision boundaries, not at
chart seams. Its flow satisfies $XF=1$, $X\vartheta=0$ and
\begin{equation}\label{eq:v83-radial-input}
 2\lambda_\circ s\le |g(\Phi_sx)|^2-|g(x)|^2
                                      \le2\Lambda_\circ s.
\end{equation}
Its endpoint-curve length is at most $C\sqrt s$. For
$I_m=1+\sum_{i=1}^{m-1}q_0^{d_i}/c_i$ one has
$|X\omega|\le C\omega I_m/|g|$ on regular survivors. At a genuine
grazing exit some internal incidence tends to zero and $\omega$ tends
to zero. Forward interior critical termination cannot occur.

Every starting point with forward lifetime at most $r$ belongs to a
layer of one of the following forms, at an actual collision or flight:
\begin{equation}\label{eq:v83-layer-input}
 c_j\le C_Eq_0^{d_j}\sqrt r,\quad
 \Delta_j\le C_Eq_0^{d'_j}\sqrt r,\quad
 \operatorname{dist}(T_R^jx,\partial Y_R^*)\le C_Eq_0^{d_j}\sqrt r.
\end{equation}
The physical coarea covariance, proved in the appendix in the convention
used here, has Jacobian $J_h=(\omega(\Phi_hx)/\omega(x))\chi(x)/\chi(\Phi_hx)$.
Writing $E_h^\varepsilon$ for all weighted starting exit mass and
\[
 O_h^\varepsilon(t)=\int_{S_h(t)} W^\varepsilon(x)
            [1-\omega(\Phi_hx)/\omega(x)]_+\,d\mu_t(x),
\]
the complete source satisfies, for each fixed $r$ and almost every $t$,
\begin{equation}\label{eq:v83-shift-input}
 b^\varepsilon(t)\le2A_rb^\varepsilon(t)+\bar E_r^\varepsilon(t)
       +\bar O_r^\varepsilon(t)
       +C_\Theta(1+\varepsilon^{-1})\sqrt r\,A_rp(t),
 \qquad A_rf(t)=\frac1{2r}\int_{t-r}^{t+r}f(s)\,ds.
\end{equation}
\end{inputstatement}
These are the endpoint-field, all-contact motion, uniform cubic-action
and separated-volume results in retained modules 174, 177, 180--182.
The guard term in \eqref{eq:v83-shift-input} is the complete product
guard, not a protective subword source. The derivative and continuation
assertions are confined to regular physical trajectories.

\begin{inputstatement}[Complete local source estimates]
\label{input:v83-local}
For every fixed auxiliary band $B_*$, constants $C_{B_*}<\infty$ and
$\rho_{B_*}<1$ give the following bounds uniformly in radius, exact
labels, interval $J$ and mark $0\le j\le m$. Let $S_{R,s}$ be the
union of strips of width $Cs$ about the finitely many coordinate lines
supporting $\partial Y_R^*$. For $s$ small,
\begin{align}
 m^2\nu\{K_m^{\rm c}=k,A_m=n,L_m\in J,
                       T_R^jx\in S_{R,s}\}
 &\le Cs^{1/16}(|J|+B_*^{-1})
   +C_{B_*}(1+B_*|J|)m^2\rho_{B_*}^{m/2},\label{eq:v83-strip-input}\\
 m^2\int_Jb^\delta(t)\,dt
 &\le C\delta^{1/16}(|J|+B_*^{-1})
   +C_{B_*}(1+B_*|J|)m^3\rho_{B_*}^{m/2}.\label{eq:v83-guard-input}
\end{align}
The local extension to a fixed remaining compact range of $\delta$
is included. At auxiliary parameter one $W^1=1$, so $b^1=p$.

Put $z=|g|$, $\alpha=129/128$ and, for $0<\eta<1/2$, define the
actual nonnegative source measures
\begin{align}
 v^{\eta,b}(t)\,dt&=L_\#(I_m z^{-\eta}\one_{\{0<z<b\}}
                                       \one_\Pi\nu_R^*),\notag\\
 v^\eta(t)\,dt&=L_\#(I_m z^{-\eta}\one_\Pi\nu_R^*).
                                                   \label{eq:v83-fractional-input}
\end{align}
They are finite and absolutely continuous at every fixed count, and
for $0<b<b_1$,
\begin{equation}\label{eq:v83-fractional-finite-input}
 m^2\int_Jv^{\eta,b}(t)\,dt
 \le C_\eta b^{\alpha-\eta}(|J|+B_*^{-1})
  +C_{B_*,\eta}(1+B_*|J|)m^2b^{1/2-\eta}\rho_{B_*}^{m/4}.
\end{equation}
Adding the bounded-gradient part gives, for every fixed $L>0$,
\begin{equation}\label{eq:v83-fixed-width-input}
 \limsup_m\sup_{R,n,k}\sup_{|J|=L}m^2\int_Jv^\eta\le C_\eta L,
 \qquad
 \limsup_m\sup_{R,n,k}\sup_{|J|=L}m^2\int_Jv^{\eta,b}
                                      \le C_\eta b^{\alpha-\eta}L.
\end{equation}
Every auxiliary band is fixed before the collision limit and enlarged
only afterwards; leading constants are independent of that band and
of the previously fixed interval length.
\end{inputstatement}
Equation~\eqref{eq:v83-strip-input} is Supplement
Theorem~\ref{supp-thm:v48-marked-strip-local}. The physical source
bound \eqref{eq:v83-guard-input} is Supplement
Theorem~\ref{supp-thm:v49-physical-remainder-local}, in particular its
finite-window estimate. The fractional bounds are the positive
layer-cake conclusions of retained module 187, stated in full above.
The joint spectral and common-space assumptions used there are recorded
in the appendix. No estimate for an untruncated first inverse gradient
or for a singular roof convolution is part of this input.

\begin{inputstatement}[The scalar positive remainder]
\label{input:v83-scalar}
For large fixed $B$, with $\varepsilon(B)=A_0B^{-1/12}$,
\begin{equation}\label{eq:v76-remainder}
 m^2p-\mathcal L=m^2b^{\varepsilon(B)}+e_B,\qquad
 \limsup_m\sup_{R,n,k}\esssup_{\mathcal C(M)}|e_B|
                                      \le C_MB^{-1/192}.
\end{equation}
The source $b^\varepsilon$ is nonnegative and has a compactly supported
$BV$ density at fixed count. No collision-uniform variation bound is
assumed. Supplement Theorem~\ref{supp-thm:v51-positive-error} gives
this remainder with the same complete finite-peak reference.
\end{inputstatement}

\subsection{Maximal segments and the two existing charges}
The retained Borel parametrization, Supplement
Theorem~\ref{supp-v85archive:thm:v85-segment-map}, is a bijection
\[
 \Psi:\coprod_j\{(\theta,u):\theta\in E_j,\ a_j(\theta)<u<b_j(\theta)\}
                      \longrightarrow\Omega,
\]
with Borel inverse. Its exact density is
$A_j(u,\theta)=\omega(\Psi_j(\theta,u))/(4\pi c_*\chi(\Psi_j(\theta,u)))$.
Every segment remains in one regular exact event until a true boundary;
Borel chart seams are not exits. Nongrazing outgoing faces retain their
incident physical component and carry
$d\zeta_+=A_j(b_j-,\theta)\,d\theta$. The full angular variation is
$|\partial_u A_j|\le C A_j I_m/z$.
These are measure identities with the null-set convention already
specified, not bounds with an unknown number of words.

The true first-exit assignment is
\[
 \mathcal F_h(t)=\int_{\Gamma_{\rm ng}^+}
  \one_{\{0<b-t\le h,\ b-t<\tau^-(y)\}}\,d\zeta_+(y),\qquad
 \overline{\mathcal F}_r=r^{-1}\int_0^r\mathcal F_h\,dh.
\]
No lower gradient cutoff is imposed. Put $l=b-a$ and define
\[
 \mathcal Z_\delta(t)=\int_{\Gamma_{\rm ng}^+}
       \one_{\{0<l\le\delta\}}\one_{\{a<t<b\}}\,d\zeta_+,
 \qquad q_\delta(t)\,dt
  =F_\#(\one_{\Omega_+^{\rm ng}}\one_{\{l\le\delta\}}\bar\nu).
\]
The second object is the original collision volume restricted to the
same actual short segments, not a volume collar for a boundary.
For source families set
\[
 \mathcal H_M(f)=\limsup_m\sup_{R,n,k}m^2
            \esssup_{t\in\mathcal C_{m,R,n,k}(M)}f(t),
 \quad \mathfrak B_M=\lim_{\delta\downarrow0}\mathcal H_M(\mathcal Z_\delta),
 \quad \mathfrak Q_M=\lim_{\delta\downarrow0}\mathcal H_M(q_\delta).
\]
These limits are monotone extended nonnegative numbers.

For $0<\eta<1/2$, with $\kappa=(1-\eta)/2$, retain
\[
 \mathsf P_{\eta,h}v^\eta(t)=\int_0^h s^{-\kappa}v^\eta(t+s)\,ds,
 \qquad\mathfrak A_{\eta,M}
       =\lim_{h\downarrow0}\mathcal H_M(\mathsf P_{\eta,h}v^\eta).
\]
Supplement Theorems~\ref{supp-v85archive:thm:v85-cutoff-surface}
and \ref{supp-v85archive:thm:v85-boundary-charge} give
\begin{equation}\label{eq:v86-surface-input}
 \mathfrak B_M\le\mathcal H_M(\overline{\mathcal F}_r)
                         \le\mathfrak B_M+C_\eta r^{1/32}.
\end{equation}
Their band sum is finite before the collision limsup at each positive
residence cutoff. Its leading constant is cutoff independent; the
complete zero-residence tail is still inside $\mathcal H_M$ before
its cutoff tends to zero. The active offset correction to historical
module 189 remains Supplement equation
\eqref{supp-v83archive:eq:v83-dyadic-tail}; the historical inclusion is
explicitly marked superseded. None of these inputs asserts charge zero.

\input{core/199_complete_mixed_optical_moments}
\input{core/200_microscopic_large_sensitivity_envelope}
\input{core/201_causal_optical_charge_and_endpoint}

\section{The scalar objective and the remaining concentration}
The finite-count envelope \eqref{eq:v86-uniform-dini} controls every
microscopic scale of its selected large-sensitivity sector at once.
It is not a reinterpretation of a fixed-width limit. Its proof keeps
the physical cap through both singular integrations and through the
complete dyadic sum. The complementary moderate-sensitivity sector,
however, has no analogous uniform envelope in this revision.

The two sufficient pairs in the original endpoint are now
$(\mathfrak C_M^{\rm mod},\mathfrak B_M)$ and
$(\mathfrak C_M^{\rm mod},\mathfrak Q_M)$, as displayed in
\eqref{eq:v86-endpoint-B}--\eqref{eq:v86-endpoint-Q}.
The reachable optical charge is at most the preceding unrestricted
fractional charge, but no converse is assumed. Both remaining
vanishing requirements are kept explicit. The organizing theorem
is not marked proved by the estimated sector alone.

After scalar closure, Supplement
Theorem~\ref{supp-thm:v73-scalar-path-equality} gives the same-roof
path numerator. Ratios remain restricted to positive arithmetic
reference classes; the return clock and forward likelihood retain
their separate inherited hypotheses. No zero class is normalized.
The full preceding article and every inherited proof are preserved
in the supplement, while source receipts and finite diagnostics remain
outside the journal narrative.

\appendix
\input{appendices/v86_common_kernels_and_dependencies}

\endgroup
\part{Complete residence bands and zero-residence source}
\begingroup
\let\ArchiveLabel\label
\let\ArchiveRef\ref
\renewcommand{\label}[1]{\ArchiveLabel{v85archive:#1}}
\renewcommand{\ref}[1]{%
  \IfBeginWith{#1}{supp-}{%
    \StrBehind{#1}{supp-}[\ArchiveTarget]%
    \ArchiveRef{\ArchiveTarget}%
  }{\ArchiveRef{v85archive:#1}}%
}
\makeatletter
\let\ltx@label\label
\makeatother
\let\ArchivePageRef\pageref
\renewcommand{\pageref}[1]{%
  \IfBeginWith{#1}{supp-}{%
    \StrBehind{#1}{supp-}[\ArchivePageTarget]%
    \ArchivePageRef{\ArchivePageTarget}%
  }{\ArchivePageRef{v85archive:#1}}%
}
\paragraph{Retained revision-85 abstract.}

We retain the exact-label pointwise local limit problem for actual returns
in a triangular Lorentz gas. A Borel transversal parametrizes complete
maximal regular roof segments and identifies their angular disintegration
with the original collision measure, without a multiplicity factor.
Residence persistence then controls the essential height of every fixed
positive residence band, including critical and nongrazing boundary pairs.
Summing any finite collection of such bands yields a full surface-height
bound of order $r^{1/32}$ at every positive residence cutoff, with a leading
constant independent of that cutoff. The short critical part has the
sharper bound $r^{129/256}(1+|\log r|)$ in the same ordered norm.
The complete zero-residence tail remains inside the collision-normalized
height before the cutoff is removed. We identify it with the limiting
surface height and compare its actual boundary traces to the original
short-segment volume source by a paid fractional optical potential.
The original pointwise objective, exact labels and arithmetic classes
are unchanged. Its final budget retains two unproved microscopic
vanishing requirements: fractional optical concentration and
zero-residence concentration.



\section{The unchanged physical problem}
Put $R_-=9/20$ and $R_+=47/100$, and let $R\in I=[R_-,R_+]$. Circular scatterers have centers in
$\Lambda=\Z(1,0)+\Z(1/2,\sqrt3/2)$. Write $T_R$ for the regular
collision map and $(\alpha,p)\in(\R/2\pi\Z)\times(-1,1)$ for its
coordinates; $p$ is the outgoing tangential velocity and
$\dd\nu=\dd\alpha\,\dd p/(4\pi)$. Physical flights satisfy
$\tau_-\le\tau_R\le\tau_+$, where $\tau_-=3/50$ and $\tau_+$
is the fixed finite-horizon bound.

The section is specified here to fix every event in the argument.
Put $h_0=1/200$. Let $Y_R$ be the union of the four open coordinate
squares of half-side $h_0$ centered at
\[
 (0,0),\quad(\pi/6,-1/2),\quad
 (-\theta_R,\sin\theta_R),\quad
 (\pi/3-\theta_R,\sin\theta_R).
\]
Here $\theta_R=2\arctan q_R$, and $q_R$ is the unique root in
$(49/100,53/100)$ of
\[
 8Rq(1+q^2)-4q(1-q^2)+\sqrt3(1-6q^2+q^4)=0.
\]
Set
\[
 U=\{|\alpha-\pi/6|<3/50,\ |p|<3/20\},\qquad
 Y_R^*=Y_R\cup U,\qquad c_*=\frac{91}{10000\pi}.
\]
The rectangles are disjoint and $\nu(Y_R^*)=c_*$. We always use
$\nu_R^*=c_*^{-1}\nu|_{Y_R^*}$, not a renormalization of a smaller
source.

At $m$ flights, let $L=L_{m,R}$ be their total length,
$K^{\rm c}_{m,R}$ their lattice displacement, and
$A_{m,R}=\sum_{j=0}^{m-1}\one_{Y_R^*}\circ T_R^j$. The event is
\begin{equation}\label{eq:v76-event}
 \Pi_{n,k,m,R}=Y_R^*\cap\{K^{\rm c}_{m,R}=k,\ A_{m,R}=n\}
                         \cap T_R^{-m}Y_R^*.
\end{equation}
The initial visit is counted; the terminal visit is tested separately.
On this event the actual $n$th return has collision count $m$ and
roof $L$. Denote its density by $p_{n,R}(k,m,t)$.

Write $c_j=\sqrt{1-p_j^2}$ and let $\Delta_j$ be the least clearance
of flight $j$ from any third disk, excluding its two incident disks.
Keep the original smooth function $0\le\Theta\le1$, zero on
$(-\infty,1]$ and one on $[2,\infty)$. With $q_0=47/53$,
$d_j=\min(j,m-j)$ and $d'_j=\min(j,m-1-j)$, set
\begin{align}
 H^\varepsilon&=
 \prod_{j=0}^{m}\Theta\!\left(\frac{c_j}{\varepsilon q_0^{d_j/4}}\right)
 \prod_{j=0}^{m-1}\Theta\!\left(\frac{\Delta_j}{\varepsilon q_0^{d'_j/4}}\right),
 \qquad W^\varepsilon=1-H^\varepsilon,\label{eq:v76-guard}\\
 b^\varepsilon(t)\dd t
   &=L_\#(W^\varepsilon\one_{\Pi_{n,k,m,R}}\nu_R^*).
 \label{eq:v76-source}
\end{align}
In particular first incidence and first clearance outside any selected
critical neighborhoods are still present. No section-distance factor
has been inserted into this guard.

\subsection{The arithmetic reference and the target}
For $y=(k_1,k_2,t,n)$ put
\[
 \mu_R=(0,0,\bar\tau_R,c_*),\quad
 Z=(y-m\mu_R)/\sqrt m,\quad
 \mathcal C_{m,R,n,k}(M)=\{t:|Z|\le M\}.
\]
We retain the finite spectral-peak reference, including peaks close to
but not on the unit circle. To specify its meaning, the collision
record observable is $(\kappa_1,\kappa_2,\tau,\one_{Y_R^*})$.
Let $\lambda_j(R,\vartheta)$ and $\mathsf P_j(R,\vartheta)$ be its
continued simple twisted eigenvalue and projection in each of the
fixed peak charts, and let $\chi_j$ be their frequency partition.
At the real maximum $a_j$ of $\log|\lambda_j|$ put
\[
 z_j=\lambda_j(R,a_j),\quad
 \mu_j=i^{-1}\nabla\log\lambda_j(R,a_j),\quad
 H_j=-\nabla^2\log\lambda_j(R,a_j),
\]
with $\operatorname{Re}H_j\ge aI_4$ for a common $a>0$. If
$\eta_R=\one_{Y_R^*}$ and $\ell$ denotes mass evaluation, put
$B_j=\chi_j(R,a_j)\ell(M_{\eta_R}\mathsf P_j(R,a_j)(\eta_R\nu))$.
Define
\begin{align}
 G_H(z)&=(2\pi)^{-4}\int_{\R^4}e^{-iv\cdot z-v^{\mathsf T}Hv/2}\dd v,
 \notag\\
 \mathcal L_{m,R}(y)&=\operatorname{Re}\left[
 c_*^{-1}\sum_j e^{-ia_j\cdot y}z_j^m B_j
      G_{H_j}\left((y-m\mu_j)/\sqrt m\right)\right].
 \label{eq:v76-reference}
\end{align}
Only the charts active at $R$ are summed. This is the same transition
reference as in the retained spectral theorem. The existence and
uniform estimates for these scalar data belong to that analytic
input, not to the transport argument below. No frequency branch or
arithmetic zero class is deleted.

The governing assertion remains
\begin{equation}\label{eq:v76-target}
 \sup_{R,\,1\le n\le m,\,k\in\Z^2}
 \esssup_{t\in\mathcal C_{m,R,n,k}(M)}
 |m^2p_{n,R}(k,m,t)-\mathcal L_{m,R}(k_1,k_2,t,n)|\longrightarrow0.
\end{equation}
Let $\mathcal E_M$ denote this error's collision limsup. An empty
central target contributes zero. The order of the uniformity is
$\sup_R\esssup_t$, so exceptional density sets may depend on $R$.

\subsection{Complete surface estimates on all positive residence scales}
An actual nongrazing outgoing face carries its backward regular lifetime
$l=b-a$ in the same exact event. Let $\overline{\mathcal F}_r$ denote
its full first-exit assignment, averaged over roof increments in $(0,r)$.
A superscript $\ge\delta$ retains only faces with $l\ge\delta$.
Among its short faces, $l\le2r$, let $\mathcal K_r$ retain
$|g(b-)|^2<8\Lambda_\circ r$.

\begin{theorem}[All positive residence scales]
\label{thm:v85-main}
Under the explicit retained inputs below, for every sufficiently small
fixed $r>0$ and every fixed $\delta>0$,
\begin{align}
 \limsup_m\sup_{R,n,k}m^2
       \|\overline{\mathcal F}_r^{\ge\delta}\|_\infty
                 &\le C_\eta r^{1/32},\label{eq:v85-main-surface}\\
 \limsup_m\sup_{R,n,k}m^2\|\mathcal K_r^{\ge\delta}\|_\infty
                 &\le C_\eta r^{129/256}(1+|\log r|).
                                               \label{eq:v85-main-critical}
\end{align}
The constants are independent of the positive residence cutoff
$\delta$. No incoming-end type, physical word or exact label is removed
within these selected families. The same first bound holds for their
short noncritical paired incoming assignment.
\end{theorem}
Theorem~\ref{thm:v85-cutoff-surface} proves these assertions from
persistence on each complete residence band. The underlying height
bound is not inferred from local mass alone. Its measurable summation
is supplied by the explicit Borel bijection in
Theorem~\ref{thm:v85-segment-map}. Proposition~\ref{prop:v85-cutoff}
also reproduces the finite-count computation of the independent
first-inverse-gradient cutoff constant used for paired transfer.

The uniform leading constants do not authorize removal of $\delta$
before the collision limit. Theorem~\ref{thm:v85-boundary-charge}
identifies the precise remaining zero-residence charge and proves that
it is the small-width limit of the full surface height. Its comparison
with an actual collision-source restriction is paid by optical
variation, not by replacing a surface with a volume collar. Together
with the canonical fractional charge, it gives the two-term
microscopic budget in Theorem~\ref{thm:v85-endpoint}. Zero-charge
estimates remain required for the original pointwise theorem.

\section{Explicit analytic and geometric inputs}
We collect the indispensable inherited assertions in the exact form
used by the new proofs. Their source theorems and operator dependencies
are specified here and in Appendix~\ref{app:v83-kernels}. The new
potential and first-exit estimates do not assume a surface spectral
insertion theorem or an unweighted pointwise local bound.

\begin{inputstatement}[Physical flow and source comparison]
\label{input:v83-geometric}
On the regular part $\Omega$ of each exact event, endpoint angles
$q=(\alpha_0,\alpha_m)$ give $F=L$, $g=\nabla_qF=R(-p_0,p_m)$,
$K=D_q^2F$ and $\vartheta=\arg g$. Put
\[
 \lambda_\circ=9\sqrt{19}/200,\qquad \Lambda_\circ=47/6,\qquad
 \omega=|\partial_{p_0}\alpha_m|^{-1},\qquad
 \chi=|\det D_q(F,\vartheta)|.
\]
The bounds $\lambda_\circ I\le K\le\Lambda_\circ I$ and
$\lambda_\circ\le\chi\le\Lambda_\circ$ hold uniformly.
The explicit field
\[
 X_q=\frac{K^{-1}g}{g^{\mathsf T}K^{-1}g}
\]
is stopped at true physical and section-decision boundaries, not at
chart seams. Its flow satisfies $XF=1$, $X\vartheta=0$ and
\begin{equation}\label{eq:v83-radial-input}
 2\lambda_\circ s\le |g(\Phi_sx)|^2-|g(x)|^2
                                      \le2\Lambda_\circ s.
\end{equation}
Its endpoint-curve length is at most $C\sqrt s$. For
$I_m=1+\sum_{i=1}^{m-1}q_0^{d_i}/c_i$ one has
$|X\omega|\le C\omega I_m/|g|$ on regular survivors. At a genuine
grazing exit some internal incidence tends to zero and $\omega$ tends
to zero. Forward interior critical termination cannot occur.

Every starting point with forward lifetime at most $r$ belongs to a
layer of one of the following forms, at an actual collision or flight:
\begin{equation}\label{eq:v83-layer-input}
 c_j\le C_Eq_0^{d_j}\sqrt r,\quad
 \Delta_j\le C_Eq_0^{d'_j}\sqrt r,\quad
 \operatorname{dist}(T_R^jx,\partial Y_R^*)\le C_Eq_0^{d_j}\sqrt r.
\end{equation}
The physical coarea covariance, proved in the appendix in the convention
used here, has Jacobian $J_h=(\omega(\Phi_hx)/\omega(x))\chi(x)/\chi(\Phi_hx)$.
Writing $E_h^\varepsilon$ for all weighted starting exit mass and
\[
 O_h^\varepsilon(t)=\int_{S_h(t)} W^\varepsilon(x)
            [1-\omega(\Phi_hx)/\omega(x)]_+\,d\mu_t(x),
\]
the complete source satisfies, for each fixed $r$ and almost every $t$,
\begin{equation}\label{eq:v83-shift-input}
 b^\varepsilon(t)\le2A_rb^\varepsilon(t)+\bar E_r^\varepsilon(t)
       +\bar O_r^\varepsilon(t)
       +C_\Theta(1+\varepsilon^{-1})\sqrt r\,A_rp(t),
 \qquad A_rf(t)=\frac1{2r}\int_{t-r}^{t+r}f(s)\,ds.
\end{equation}
\end{inputstatement}
These are the endpoint-field, all-contact motion, uniform cubic-action
and separated-volume results in retained modules 174, 177, 180--182.
The guard term in \eqref{eq:v83-shift-input} is the complete product
guard, not a protective subword source. The derivative and continuation
assertions are confined to regular physical trajectories.

\begin{inputstatement}[Complete local source estimates]
\label{input:v83-local}
For every fixed auxiliary band $B_*$, constants $C_{B_*}<\infty$ and
$\rho_{B_*}<1$ give the following bounds uniformly in radius, exact
labels, interval $J$ and mark $0\le j\le m$. Let $S_{R,s}$ be the
union of strips of width $Cs$ about the finitely many coordinate lines
supporting $\partial Y_R^*$. For $s$ small,
\begin{align}
 m^2\nu\{K_m^{\rm c}=k,A_m=n,L_m\in J,
                       T_R^jx\in S_{R,s}\}
 &\le Cs^{1/16}(|J|+B_*^{-1})
   +C_{B_*}(1+B_*|J|)m^2\rho_{B_*}^{m/2},\label{eq:v83-strip-input}\\
 m^2\int_Jb^\delta(t)\,dt
 &\le C\delta^{1/16}(|J|+B_*^{-1})
   +C_{B_*}(1+B_*|J|)m^3\rho_{B_*}^{m/2}.\label{eq:v83-guard-input}
\end{align}
The local extension to a fixed remaining compact range of $\delta$
is included. At auxiliary parameter one $W^1=1$, so $b^1=p$.

Put $z=|g|$, $\alpha=129/128$ and, for $0<\eta<1/2$, define the
actual nonnegative source measures
\begin{align}
 v^{\eta,b}(t)\,dt&=L_\#(I_m z^{-\eta}\one_{\{0<z<b\}}
                                       \one_\Pi\nu_R^*),\notag\\
 v^\eta(t)\,dt&=L_\#(I_m z^{-\eta}\one_\Pi\nu_R^*).
                                                   \label{eq:v83-fractional-input}
\end{align}
They are finite and absolutely continuous at every fixed count, and
for $0<b<b_1$,
\begin{equation}\label{eq:v83-fractional-finite-input}
 m^2\int_Jv^{\eta,b}(t)\,dt
 \le C_\eta b^{\alpha-\eta}(|J|+B_*^{-1})
  +C_{B_*,\eta}(1+B_*|J|)m^2b^{1/2-\eta}\rho_{B_*}^{m/4}.
\end{equation}
Adding the bounded-gradient part gives, for every fixed $L>0$,
\begin{equation}\label{eq:v83-fixed-width-input}
 \limsup_m\sup_{R,n,k}\sup_{|J|=L}m^2\int_Jv^\eta\le C_\eta L,
 \qquad
 \limsup_m\sup_{R,n,k}\sup_{|J|=L}m^2\int_Jv^{\eta,b}
                                      \le C_\eta b^{\alpha-\eta}L.
\end{equation}
Every auxiliary band is fixed before the collision limit and enlarged
only afterwards; leading constants are independent of that band and
of the previously fixed interval length.
\end{inputstatement}
Equation~\eqref{eq:v83-strip-input} is Supplement
Theorem~\ref{supp-thm:v48-marked-strip-local}. The physical source
bound \eqref{eq:v83-guard-input} is Supplement
Theorem~\ref{supp-thm:v49-physical-remainder-local}, in particular its
finite-window estimate. The fractional bounds are the positive
layer-cake conclusions of retained module 187, stated in full above.
The joint spectral and common-space assumptions used there are recorded
in the appendix. No estimate for an untruncated first inverse gradient
or for a singular roof convolution is part of this input.

\begin{inputstatement}[The scalar positive remainder]
\label{input:v83-scalar}
For large fixed $B$, with $\varepsilon(B)=A_0B^{-1/12}$,
\begin{equation}\label{eq:v76-remainder}
 m^2p-\mathcal L=m^2b^{\varepsilon(B)}+e_B,\qquad
 \limsup_m\sup_{R,n,k}\esssup_{\mathcal C(M)}|e_B|
                                      \le C_MB^{-1/192}.
\end{equation}
The source $b^\varepsilon$ is nonnegative and has a compactly supported
$BV$ density at fixed count. No collision-uniform variation bound is
assumed. Supplement Theorem~\ref{supp-thm:v51-positive-error} gives
this remainder with the same complete finite-peak reference.
\end{inputstatement}

\input{core/190_untruncated_optical_potentials}
\input{core/191_complete_first_exit_surface_mass}
\input{core/192_ordered_fractional_tail_and_endpoint}

\input{core/196_measurable_segments_and_cutoff_audit}
\input{core/197_complete_residence_band_heights}
\input{core/198_zero_residence_charge_and_original_source}


\section{The original scalar objective and downstream normalization}
The governing assertion remains \eqref{eq:v76-target}. The full
surface is now estimated in essential height at every previously fixed
positive residence cutoff, with one leading constant for all such
cutoffs. The critical short surface has the stronger height estimate
\eqref{eq:v85-main-critical} under that same ordering. These are
complete positive-scale theorems, not a sum over selected words.

The complete zero-residence tail, and the complete fractional optical
tail, remain in the original normalized essential supremum before the
collision limit. Neither vanishing assertion follows by letting a
cutoff depend on collision count. Their exact role is
\eqref{eq:v85-endpoint}; the surface-source comparison
\eqref{eq:v85-trace-volume} does not suppress a boundary flux or change
the guard. The manuscript does not mark the organizing theorem proved
until these microscopic estimates are established.

The active dyadic formula is the offset expression
\eqref{eq:v83-dyadic-tail}; its negative-shift domain is explicit in
\eqref{eq:v85-charge-domain}. The historical module 189 remains
qualified at its inclusion in the supplement, with original bytes
preserved. The supplement also contains a compact input-to-theorem,
equation and proof-page map (Supplement
page~\pageref{supp-v85:dependency-guide}).

After scalar closure, the positive path-remainder equality in
Supplement Theorem~\ref{supp-thm:v73-scalar-path-equality} supplies
the same-roof numerator. Ratios are formed only on positive arithmetic
reference classes. The actual-return clock and forward likelihood
retain their separately stated hypotheses. No surface estimate is
used to reverse this scalar-first logical order.

\appendix
\input{appendices/v83_common_kernels_and_dependencies}

\endgroup
\part{Residence persistence and paired physical boundary transfer}
\begingroup
\let\ArchiveLabel\label
\let\ArchiveRef\ref
\renewcommand{\label}[1]{\ArchiveLabel{v84archive:#1}}
\renewcommand{\ref}[1]{%
  \IfBeginWith{#1}{supp-}{%
    \StrBehind{#1}{supp-}[\ArchiveTarget]%
    \ArchiveRef{\ArchiveTarget}%
  }{\ArchiveRef{v84archive:#1}}%
}
\makeatletter
\let\ltx@label\label
\makeatother
\paragraph{Retained revision-84 abstract.}

We study the original exact-label pointwise local limit problem for
actual returns in a triangular Lorentz gas. The backward lifetime of
a one-sided outgoing face gives a residence kernel for its complete
first-exit assignment. This kernel converts the established local
surface-mass theorem into an essential-height bound of order
$r^{1/32}$ for all regular segments of length greater than $2r$.
For shorter segments away from a critical outgoing gradient, the
physical angular density transfers between the two true ends with
an error of order $r$ in the same collision-normalized essential norm.
Grazing incoming ends contribute no unpaid boundary trace. The
remaining short critical surface has a sharper local-mass bound.
We also identify a canonical microscopic concentration charge and
prove dyadic outer-scale rigidity of the fractional tail. The exact
labels, original section, guard, arithmetic reference and first-exit
selection remain unchanged. The final budget isolates paired short
nongrazing boundaries, short critical segments and a microscopic
charge; their remaining essential-height estimates are stated
separately from the new proved bounds.



\section{The unchanged physical problem}
Put $R_-=9/20$ and $R_+=47/100$, and let $R\in I=[R_-,R_+]$. Circular scatterers have centers in
$\Lambda=\Z(1,0)+\Z(1/2,\sqrt3/2)$. Write $T_R$ for the regular
collision map and $(\alpha,p)\in(\R/2\pi\Z)\times(-1,1)$ for its
coordinates; $p$ is the outgoing tangential velocity and
$\dd\nu=\dd\alpha\,\dd p/(4\pi)$. Physical flights satisfy
$\tau_-\le\tau_R\le\tau_+$, where $\tau_-=3/50$ and $\tau_+$
is the fixed finite-horizon bound.

The section is specified here to fix every event in the argument.
Put $h_0=1/200$. Let $Y_R$ be the union of the four open coordinate
squares of half-side $h_0$ centered at
\[
 (0,0),\quad(\pi/6,-1/2),\quad
 (-\theta_R,\sin\theta_R),\quad
 (\pi/3-\theta_R,\sin\theta_R).
\]
Here $\theta_R=2\arctan q_R$, and $q_R$ is the unique root in
$(49/100,53/100)$ of
\[
 8Rq(1+q^2)-4q(1-q^2)+\sqrt3(1-6q^2+q^4)=0.
\]
Set
\[
 U=\{|\alpha-\pi/6|<3/50,\ |p|<3/20\},\qquad
 Y_R^*=Y_R\cup U,\qquad c_*=\frac{91}{10000\pi}.
\]
The rectangles are disjoint and $\nu(Y_R^*)=c_*$. We always use
$\nu_R^*=c_*^{-1}\nu|_{Y_R^*}$, not a renormalization of a smaller
source.

At $m$ flights, let $L=L_{m,R}$ be their total length,
$K^{\rm c}_{m,R}$ their lattice displacement, and
$A_{m,R}=\sum_{j=0}^{m-1}\one_{Y_R^*}\circ T_R^j$. The event is
\begin{equation}\label{eq:v76-event}
 \Pi_{n,k,m,R}=Y_R^*\cap\{K^{\rm c}_{m,R}=k,\ A_{m,R}=n\}
                         \cap T_R^{-m}Y_R^*.
\end{equation}
The initial visit is counted; the terminal visit is tested separately.
On this event the actual $n$th return has collision count $m$ and
roof $L$. Denote its density by $p_{n,R}(k,m,t)$.

Write $c_j=\sqrt{1-p_j^2}$ and let $\Delta_j$ be the least clearance
of flight $j$ from any third disk, excluding its two incident disks.
Keep the original smooth function $0\le\Theta\le1$, zero on
$(-\infty,1]$ and one on $[2,\infty)$. With $q_0=47/53$,
$d_j=\min(j,m-j)$ and $d'_j=\min(j,m-1-j)$, set
\begin{align}
 H^\varepsilon&=
 \prod_{j=0}^{m}\Theta\!\left(\frac{c_j}{\varepsilon q_0^{d_j/4}}\right)
 \prod_{j=0}^{m-1}\Theta\!\left(\frac{\Delta_j}{\varepsilon q_0^{d'_j/4}}\right),
 \qquad W^\varepsilon=1-H^\varepsilon,\label{eq:v76-guard}\\
 b^\varepsilon(t)\dd t
   &=L_\#(W^\varepsilon\one_{\Pi_{n,k,m,R}}\nu_R^*).
 \label{eq:v76-source}
\end{align}
In particular first incidence and first clearance outside any selected
critical neighborhoods are still present. No section-distance factor
has been inserted into this guard.

\subsection{The arithmetic reference and the target}
For $y=(k_1,k_2,t,n)$ put
\[
 \mu_R=(0,0,\bar\tau_R,c_*),\quad
 Z=(y-m\mu_R)/\sqrt m,\quad
 \mathcal C_{m,R,n,k}(M)=\{t:|Z|\le M\}.
\]
We retain the finite spectral-peak reference, including peaks close to
but not on the unit circle. To specify its meaning, the collision
record observable is $(\kappa_1,\kappa_2,\tau,\one_{Y_R^*})$.
Let $\lambda_j(R,\vartheta)$ and $\mathsf P_j(R,\vartheta)$ be its
continued simple twisted eigenvalue and projection in each of the
fixed peak charts, and let $\chi_j$ be their frequency partition.
At the real maximum $a_j$ of $\log|\lambda_j|$ put
\[
 z_j=\lambda_j(R,a_j),\quad
 \mu_j=i^{-1}\nabla\log\lambda_j(R,a_j),\quad
 H_j=-\nabla^2\log\lambda_j(R,a_j),
\]
with $\operatorname{Re}H_j\ge aI_4$ for a common $a>0$. If
$\eta_R=\one_{Y_R^*}$ and $\ell$ denotes mass evaluation, put
$B_j=\chi_j(R,a_j)\ell(M_{\eta_R}\mathsf P_j(R,a_j)(\eta_R\nu))$.
Define
\begin{align}
 G_H(z)&=(2\pi)^{-4}\int_{\R^4}e^{-iv\cdot z-v^{\mathsf T}Hv/2}\dd v,
 \notag\\
 \mathcal L_{m,R}(y)&=\operatorname{Re}\left[
 c_*^{-1}\sum_j e^{-ia_j\cdot y}z_j^m B_j
      G_{H_j}\left((y-m\mu_j)/\sqrt m\right)\right].
 \label{eq:v76-reference}
\end{align}
Only the charts active at $R$ are summed. This is the same transition
reference as in the retained spectral theorem. The existence and
uniform estimates for these scalar data belong to that analytic
input, not to the transport argument below. No frequency branch or
arithmetic zero class is deleted.

The governing assertion remains
\begin{equation}\label{eq:v76-target}
 \sup_{R,\,1\le n\le m,\,k\in\Z^2}
 \esssup_{t\in\mathcal C_{m,R,n,k}(M)}
 |m^2p_{n,R}(k,m,t)-\mathcal L_{m,R}(k_1,k_2,t,n)|\longrightarrow0.
\end{equation}
Let $\mathcal E_M$ denote this error's collision limsup. An empty
central target contributes zero. The order of the uniformity is
$\sup_R\esssup_t$, so exceptional density sets may depend on $R$.

\subsection{The complete surface-height estimates proved here}
For each one-sided nongrazing outgoing face, let $l$ be its backward
lifetime in the original regular exact event. Write
$\overline{\mathcal F}_r$ for the complete first-exit surface assignment
averaged over roof increments in $(0,r)$. The source is not replaced
by a boundary collar. Its restriction to $l>2r$ is denoted by
$\mathcal L_r$. Among the remaining segments, let
$\mathcal S_r^{\rm nc}$ have outgoing $|g|^2\ge8\Lambda_\circ r$,
and let $\mathcal J_r$ carry their incoming physical traces, paired
with the same outgoing faces and the same exit-window kernel.

\begin{theorem}[Residence and two-ended optical transfer]
\label{thm:v84-main}
Under the explicit inherited inputs below, for every sufficiently
small previously fixed $r>0$,
\begin{align}
 \limsup_m\sup_{R,n,k}m^2\|\mathcal L_r\|_\infty
                          &\le Cr^{1/32},\label{eq:v84-main-long}\\
 \limsup_m\sup_{R,n,k}m^2
       \|\mathcal S_r^{\rm nc}-\mathcal J_r\|_\infty
                          &\le Cr.\label{eq:v84-main-pair}
\end{align}
The complete short noncritical part with grazing incoming end also
has height at most $Cr$. All comparisons retain the original exact
label and every incident physical branch; no word-count multiplier
is used.
\end{theorem}
The first conclusion is Theorem~\ref{thm:v84-long-height}. It uses a
persistence property of the actual surface kernel, not an inference
from small mass alone. The second is
Theorem~\ref{thm:v84-paired-transfer}; it estimates the change of
physical angular density between the two ends, including both angular
Jacobian factors. The positive lower gradient used inside that proof
is a function of the fixed roof width only.

The preceding untruncated fractional transport and local surface-mass
results are retained with their proofs. The new argument uses them to
control the complementary long-residence source in essential height.
The still unresolved sources are short paired nongrazing boundaries,
short critical segments, and a canonical microscopic concentration
charge. Theorem~\ref{thm:v84-endpoint} gives their exact role in the
unchanged pointwise assertion \eqref{eq:v76-target}.

\section{Explicit analytic and geometric inputs}
We collect the indispensable inherited assertions in the exact form
used by the new proofs. Their source theorems and operator dependencies
are specified here and in Appendix~\ref{app:v83-kernels}. The new
potential and first-exit estimates do not assume a surface spectral
insertion theorem or an unweighted pointwise local bound.

\begin{inputstatement}[Physical flow and source comparison]
\label{input:v83-geometric}
On the regular part $\Omega$ of each exact event, endpoint angles
$q=(\alpha_0,\alpha_m)$ give $F=L$, $g=\nabla_qF=R(-p_0,p_m)$,
$K=D_q^2F$ and $\vartheta=\arg g$. Put
\[
 \lambda_\circ=9\sqrt{19}/200,\qquad \Lambda_\circ=47/6,\qquad
 \omega=|\partial_{p_0}\alpha_m|^{-1},\qquad
 \chi=|\det D_q(F,\vartheta)|.
\]
The bounds $\lambda_\circ I\le K\le\Lambda_\circ I$ and
$\lambda_\circ\le\chi\le\Lambda_\circ$ hold uniformly.
The explicit field
\[
 X_q=\frac{K^{-1}g}{g^{\mathsf T}K^{-1}g}
\]
is stopped at true physical and section-decision boundaries, not at
chart seams. Its flow satisfies $XF=1$, $X\vartheta=0$ and
\begin{equation}\label{eq:v83-radial-input}
 2\lambda_\circ s\le |g(\Phi_sx)|^2-|g(x)|^2
                                      \le2\Lambda_\circ s.
\end{equation}
Its endpoint-curve length is at most $C\sqrt s$. For
$I_m=1+\sum_{i=1}^{m-1}q_0^{d_i}/c_i$ one has
$|X\omega|\le C\omega I_m/|g|$ on regular survivors. At a genuine
grazing exit some internal incidence tends to zero and $\omega$ tends
to zero. Forward interior critical termination cannot occur.

Every starting point with forward lifetime at most $r$ belongs to a
layer of one of the following forms, at an actual collision or flight:
\begin{equation}\label{eq:v83-layer-input}
 c_j\le C_Eq_0^{d_j}\sqrt r,\quad
 \Delta_j\le C_Eq_0^{d'_j}\sqrt r,\quad
 \operatorname{dist}(T_R^jx,\partial Y_R^*)\le C_Eq_0^{d_j}\sqrt r.
\end{equation}
The physical coarea covariance, proved in the appendix in the convention
used here, has Jacobian $J_h=(\omega(\Phi_hx)/\omega(x))\chi(x)/\chi(\Phi_hx)$.
Writing $E_h^\varepsilon$ for all weighted starting exit mass and
\[
 O_h^\varepsilon(t)=\int_{S_h(t)} W^\varepsilon(x)
            [1-\omega(\Phi_hx)/\omega(x)]_+\,d\mu_t(x),
\]
the complete source satisfies, for each fixed $r$ and almost every $t$,
\begin{equation}\label{eq:v83-shift-input}
 b^\varepsilon(t)\le2A_rb^\varepsilon(t)+\bar E_r^\varepsilon(t)
       +\bar O_r^\varepsilon(t)
       +C_\Theta(1+\varepsilon^{-1})\sqrt r\,A_rp(t),
 \qquad A_rf(t)=\frac1{2r}\int_{t-r}^{t+r}f(s)\,ds.
\end{equation}
\end{inputstatement}
These are the endpoint-field, all-contact motion, uniform cubic-action
and separated-volume results in retained modules 174, 177, 180--182.
The guard term in \eqref{eq:v83-shift-input} is the complete product
guard, not a protective subword source. The derivative and continuation
assertions are confined to regular physical trajectories.

\begin{inputstatement}[Complete local source estimates]
\label{input:v83-local}
For every fixed auxiliary band $B_*$, constants $C_{B_*}<\infty$ and
$\rho_{B_*}<1$ give the following bounds uniformly in radius, exact
labels, interval $J$ and mark $0\le j\le m$. Let $S_{R,s}$ be the
union of strips of width $Cs$ about the finitely many coordinate lines
supporting $\partial Y_R^*$. For $s$ small,
\begin{align}
 m^2\nu\{K_m^{\rm c}=k,A_m=n,L_m\in J,
                       T_R^jx\in S_{R,s}\}
 &\le Cs^{1/16}(|J|+B_*^{-1})
   +C_{B_*}(1+B_*|J|)m^2\rho_{B_*}^{m/2},\label{eq:v83-strip-input}\\
 m^2\int_Jb^\delta(t)\,dt
 &\le C\delta^{1/16}(|J|+B_*^{-1})
   +C_{B_*}(1+B_*|J|)m^3\rho_{B_*}^{m/2}.\label{eq:v83-guard-input}
\end{align}
The local extension to a fixed remaining compact range of $\delta$
is included. At auxiliary parameter one $W^1=1$, so $b^1=p$.

Put $z=|g|$, $\alpha=129/128$ and, for $0<\eta<1/2$, define the
actual nonnegative source measures
\begin{align}
 v^{\eta,b}(t)\,dt&=L_\#(I_m z^{-\eta}\one_{\{0<z<b\}}
                                       \one_\Pi\nu_R^*),\notag\\
 v^\eta(t)\,dt&=L_\#(I_m z^{-\eta}\one_\Pi\nu_R^*).
                                                   \label{eq:v83-fractional-input}
\end{align}
They are finite and absolutely continuous at every fixed count, and
for $0<b<b_1$,
\begin{equation}\label{eq:v83-fractional-finite-input}
 m^2\int_Jv^{\eta,b}(t)\,dt
 \le C_\eta b^{\alpha-\eta}(|J|+B_*^{-1})
  +C_{B_*,\eta}(1+B_*|J|)m^2b^{1/2-\eta}\rho_{B_*}^{m/4}.
\end{equation}
Adding the bounded-gradient part gives, for every fixed $L>0$,
\begin{equation}\label{eq:v83-fixed-width-input}
 \limsup_m\sup_{R,n,k}\sup_{|J|=L}m^2\int_Jv^\eta\le C_\eta L,
 \qquad
 \limsup_m\sup_{R,n,k}\sup_{|J|=L}m^2\int_Jv^{\eta,b}
                                      \le C_\eta b^{\alpha-\eta}L.
\end{equation}
Every auxiliary band is fixed before the collision limit and enlarged
only afterwards; leading constants are independent of that band and
of the previously fixed interval length.
\end{inputstatement}
Equation~\eqref{eq:v83-strip-input} is Supplement
Theorem~\ref{supp-thm:v48-marked-strip-local}. The physical source
bound \eqref{eq:v83-guard-input} is Supplement
Theorem~\ref{supp-thm:v49-physical-remainder-local}, in particular its
finite-window estimate. The fractional bounds are the positive
layer-cake conclusions of retained module 187, stated in full above.
The joint spectral and common-space assumptions used there are recorded
in the appendix. No estimate for an untruncated first inverse gradient
or for a singular roof convolution is part of this input.

\begin{inputstatement}[The scalar positive remainder]
\label{input:v83-scalar}
For large fixed $B$, with $\varepsilon(B)=A_0B^{-1/12}$,
\begin{equation}\label{eq:v76-remainder}
 m^2p-\mathcal L=m^2b^{\varepsilon(B)}+e_B,\qquad
 \limsup_m\sup_{R,n,k}\esssup_{\mathcal C(M)}|e_B|
                                      \le C_MB^{-1/192}.
\end{equation}
The source $b^\varepsilon$ is nonnegative and has a compactly supported
$BV$ density at fixed count. No collision-uniform variation bound is
assumed. Supplement Theorem~\ref{supp-thm:v51-positive-error} gives
this remainder with the same complete finite-peak reference.
\end{inputstatement}

\input{core/190_untruncated_optical_potentials}
\input{core/191_complete_first_exit_surface_mass}
\input{core/192_ordered_fractional_tail_and_endpoint}

\input{core/193_backward_residence_surface_height}
\input{core/194_paired_boundary_optical_transfer}
\input{core/195_canonical_concentration_and_endpoint}


\section{The original scalar objective and downstream normalization}
The governing assertion remains \eqref{eq:v76-target}. The long-residence
surface and the short noncritical optical transfer now have proved
essential-height bounds. Their complements are not removed: the
paired nongrazing trace and short critical surface remain inside the
same exact-label norm. The latter has the quantitative local-mass
bound \eqref{eq:v84-critical-mass}, which is not its height bound.
The canonical microscopic charge in \eqref{eq:v84-canonical-charge}
requires a separate vanishing estimate; dyadic outer-scale rigidity
shows precisely why reducing the reconstruction width cannot replace
that estimate.

The only active dyadic convention in this article is the offset formula
\eqref{eq:v83-dyadic-tail}. The retained module 189 is marked as
superseded at its point of inclusion in the supplement, and its original
bytes remain in the historical source. Its ambiguous integration
notation is not an active theorem used here.

After scalar closure, the positive path-remainder equality in
Supplement Theorem~\ref{supp-thm:v73-scalar-path-equality} supplies
the same-roof path numerator. Ratios are formed only on positive
arithmetic reference classes; the actual-return clock and forward
likelihood retain their separate inherited hypotheses. The unchanged
scalar-first order is not reversed by any surface comparison.

\appendix
\input{appendices/v83_common_kernels_and_dependencies}

\endgroup
\part{Untruncated optical potentials and complete first-exit surface mass}
\begingroup
\let\ArchiveLabel\label
\let\ArchiveRef\ref
\renewcommand{\label}[1]{\ArchiveLabel{v83archive:#1}}
\renewcommand{\ref}[1]{%
  \IfBeginWith{#1}{supp-}{%
    \StrBehind{#1}{supp-}[\ArchiveTarget]%
    \ArchiveRef{\ArchiveTarget}%
  }{\ArchiveRef{v83archive:#1}}%
}
\makeatletter
\let\ltx@label\label
\makeatother
\paragraph{Retained revision-83 abstract.}

We study the pointwise local limit problem for the full displacement,
collision-count, occupation and roof record of actual returns in a
triangular Lorentz gas. Actual fractional inverse-gradient insertions
yield complete optical transport estimates without an initial critical
cutoff. The optical loss and genuine grazing starting mass obey local
weak-$L^q$ bounds for every fixed $2<q<4$, with an additional power
on the critical initial source. We estimate the actual nongrazing
first-exit surface assignment by paying both signs of optical variation.
Together with all-depth physical margin localization, this gives a
collision-normalized $O(r^{33/32})$ bound for its mass on each previously
fixed roof interval of length $r$. The exact labels, original guard and
one-sided first-exit selection are retained. A separate ordered-tail
argument identifies the microscopic estimate still required for the
original essential-height law; local norm and surface mass bounds are
not substituted for that pointwise estimate.



\section{The unchanged physical problem}
Put $R_-=9/20$ and $R_+=47/100$, and let $R\in I=[R_-,R_+]$. Circular scatterers have centers in
$\Lambda=\Z(1,0)+\Z(1/2,\sqrt3/2)$. Write $T_R$ for the regular
collision map and $(\alpha,p)\in(\R/2\pi\Z)\times(-1,1)$ for its
coordinates; $p$ is the outgoing tangential velocity and
$\dd\nu=\dd\alpha\,\dd p/(4\pi)$. Physical flights satisfy
$\tau_-\le\tau_R\le\tau_+$, where $\tau_-=3/50$ and $\tau_+$
is the fixed finite-horizon bound.

The section is specified here to fix every event in the argument.
Put $h_0=1/200$. Let $Y_R$ be the union of the four open coordinate
squares of half-side $h_0$ centered at
\[
 (0,0),\quad(\pi/6,-1/2),\quad
 (-\theta_R,\sin\theta_R),\quad
 (\pi/3-\theta_R,\sin\theta_R).
\]
Here $\theta_R=2\arctan q_R$, and $q_R$ is the unique root in
$(49/100,53/100)$ of
\[
 8Rq(1+q^2)-4q(1-q^2)+\sqrt3(1-6q^2+q^4)=0.
\]
Set
\[
 U=\{|\alpha-\pi/6|<3/50,\ |p|<3/20\},\qquad
 Y_R^*=Y_R\cup U,\qquad c_*=\frac{91}{10000\pi}.
\]
The rectangles are disjoint and $\nu(Y_R^*)=c_*$. We always use
$\nu_R^*=c_*^{-1}\nu|_{Y_R^*}$, not a renormalization of a smaller
source.

At $m$ flights, let $L=L_{m,R}$ be their total length,
$K^{\rm c}_{m,R}$ their lattice displacement, and
$A_{m,R}=\sum_{j=0}^{m-1}\one_{Y_R^*}\circ T_R^j$. The event is
\begin{equation}\label{eq:v76-event}
 \Pi_{n,k,m,R}=Y_R^*\cap\{K^{\rm c}_{m,R}=k,\ A_{m,R}=n\}
                         \cap T_R^{-m}Y_R^*.
\end{equation}
The initial visit is counted; the terminal visit is tested separately.
On this event the actual $n$th return has collision count $m$ and
roof $L$. Denote its density by $p_{n,R}(k,m,t)$.

Write $c_j=\sqrt{1-p_j^2}$ and let $\Delta_j$ be the least clearance
of flight $j$ from any third disk, excluding its two incident disks.
Keep the original smooth function $0\le\Theta\le1$, zero on
$(-\infty,1]$ and one on $[2,\infty)$. With $q_0=47/53$,
$d_j=\min(j,m-j)$ and $d'_j=\min(j,m-1-j)$, set
\begin{align}
 H^\varepsilon&=
 \prod_{j=0}^{m}\Theta\!\left(\frac{c_j}{\varepsilon q_0^{d_j/4}}\right)
 \prod_{j=0}^{m-1}\Theta\!\left(\frac{\Delta_j}{\varepsilon q_0^{d'_j/4}}\right),
 \qquad W^\varepsilon=1-H^\varepsilon,\label{eq:v76-guard}\\
 b^\varepsilon(t)\dd t
   &=L_\#(W^\varepsilon\one_{\Pi_{n,k,m,R}}\nu_R^*).
 \label{eq:v76-source}
\end{align}
In particular first incidence and first clearance outside any selected
critical neighborhoods are still present. No section-distance factor
has been inserted into this guard.

\subsection{The arithmetic reference and the target}
For $y=(k_1,k_2,t,n)$ put
\[
 \mu_R=(0,0,\bar\tau_R,c_*),\quad
 Z=(y-m\mu_R)/\sqrt m,\quad
 \mathcal C_{m,R,n,k}(M)=\{t:|Z|\le M\}.
\]
We retain the finite spectral-peak reference, including peaks close to
but not on the unit circle. To specify its meaning, the collision
record observable is $(\kappa_1,\kappa_2,\tau,\one_{Y_R^*})$.
Let $\lambda_j(R,\vartheta)$ and $\mathsf P_j(R,\vartheta)$ be its
continued simple twisted eigenvalue and projection in each of the
fixed peak charts, and let $\chi_j$ be their frequency partition.
At the real maximum $a_j$ of $\log|\lambda_j|$ put
\[
 z_j=\lambda_j(R,a_j),\quad
 \mu_j=i^{-1}\nabla\log\lambda_j(R,a_j),\quad
 H_j=-\nabla^2\log\lambda_j(R,a_j),
\]
with $\operatorname{Re}H_j\ge aI_4$ for a common $a>0$. If
$\eta_R=\one_{Y_R^*}$ and $\ell$ denotes mass evaluation, put
$B_j=\chi_j(R,a_j)\ell(M_{\eta_R}\mathsf P_j(R,a_j)(\eta_R\nu))$.
Define
\begin{align}
 G_H(z)&=(2\pi)^{-4}\int_{\R^4}e^{-iv\cdot z-v^{\mathsf T}Hv/2}\dd v,
 \notag\\
 \mathcal L_{m,R}(y)&=\operatorname{Re}\left[
 c_*^{-1}\sum_j e^{-ia_j\cdot y}z_j^m B_j
      G_{H_j}\left((y-m\mu_j)/\sqrt m\right)\right].
 \label{eq:v76-reference}
\end{align}
Only the charts active at $R$ are summed. This is the same transition
reference as in the retained spectral theorem. The existence and
uniform estimates for these scalar data belong to that analytic
input, not to the transport argument below. No frequency branch or
arithmetic zero class is deleted.

The governing assertion remains
\begin{equation}\label{eq:v76-target}
 \sup_{R,\,1\le n\le m,\,k\in\Z^2}
 \esssup_{t\in\mathcal C_{m,R,n,k}(M)}
 |m^2p_{n,R}(k,m,t)-\mathcal L_{m,R}(k_1,k_2,t,n)|\longrightarrow0.
\end{equation}
Let $\mathcal E_M$ denote this error's collision limsup. An empty
central target contributes zero. The order of the uniformity is
$\sup_R\esssup_t$, so exceptional density sets may depend on $R$.

\subsection{The complete bounds proved here}
The two positive transport mechanisms will be treated at the level of
the full original source. Let $O_h^\varepsilon$ denote positive optical
loss on survivors, $E_h^{\varepsilon,\mathrm{gr}}$ the starting mass of
genuine grazing first exits, and $\mathcal F_h$ the nongrazing outgoing
first-exit flux assigned to the starting roof. Their precise definitions
are given below and in \eqref{eq:v83-untruncated-flux}. A bar denotes
$h$-averaging over $(0,r)$. For $q>1$ use the normalization
$\mathcal N_q$ in \eqref{eq:v83-lorentz-norm}.

\begin{theorem}[Complete optical and first-exit estimates]
\label{thm:v83-main}
Under the explicit inherited inputs below, for every fixed $2<q<4$
and every previously fixed $r,L>0$,
\begin{align}
 \limsup_m\sup_{R,n,k}\sup_{|J|=L}
 \mathcal N_q\bigl(m^2(\bar O_r^\varepsilon+
             \bar E_r^{\varepsilon,\mathrm{gr}});J\bigr)
                                     &\le C_q(L+r),\label{eq:v83-main-optical}\\
 \limsup_m\sup_{R,n,k}\sup_{|J|=r}
           m^2\int_J\overline{\mathcal F}_r(t)\,dt
                                     &\le C_q r^{33/32}
                               \quad(0<r<r_0).\label{eq:v83-main-surface}
\end{align}
Both sources are untruncated in the initial endpoint gradient and in
all positive internal incidences. They retain the original exact label,
initial and terminal tests, and true one-sided first-exit selection.
\end{theorem}
The first assertion is proved in Theorem~\ref{thm:v83-complete-weak-height}
with $\eta=1-2/q$, and includes sharper bounds on a critical initial
neighborhood. The second is Theorem~\ref{thm:v83-surface-local-mass}.
Their proofs are complete relative to the stated inputs. They do not
assert \eqref{eq:v76-target}, whose remaining essential-height estimate
is formulated in Theorem~\ref{thm:v83-endpoint}. In particular the
window $r$ is never chosen as a function of the collision count.

\section{Explicit analytic and geometric inputs}
We collect the indispensable inherited assertions in the exact form
used by the new proofs. Their source theorems and operator dependencies
are specified here and in Appendix~\ref{app:v83-kernels}. The new
potential and first-exit estimates do not assume a surface spectral
insertion theorem or an unweighted pointwise local bound.

\begin{inputstatement}[Physical flow and source comparison]
\label{input:v83-geometric}
On the regular part $\Omega$ of each exact event, endpoint angles
$q=(\alpha_0,\alpha_m)$ give $F=L$, $g=\nabla_qF=R(-p_0,p_m)$,
$K=D_q^2F$ and $\vartheta=\arg g$. Put
\[
 \lambda_\circ=9\sqrt{19}/200,\qquad \Lambda_\circ=47/6,\qquad
 \omega=|\partial_{p_0}\alpha_m|^{-1},\qquad
 \chi=|\det D_q(F,\vartheta)|.
\]
The bounds $\lambda_\circ I\le K\le\Lambda_\circ I$ and
$\lambda_\circ\le\chi\le\Lambda_\circ$ hold uniformly.
The explicit field
\[
 X_q=\frac{K^{-1}g}{g^{\mathsf T}K^{-1}g}
\]
is stopped at true physical and section-decision boundaries, not at
chart seams. Its flow satisfies $XF=1$, $X\vartheta=0$ and
\begin{equation}\label{eq:v83-radial-input}
 2\lambda_\circ s\le |g(\Phi_sx)|^2-|g(x)|^2
                                      \le2\Lambda_\circ s.
\end{equation}
Its endpoint-curve length is at most $C\sqrt s$. For
$I_m=1+\sum_{i=1}^{m-1}q_0^{d_i}/c_i$ one has
$|X\omega|\le C\omega I_m/|g|$ on regular survivors. At a genuine
grazing exit some internal incidence tends to zero and $\omega$ tends
to zero. Forward interior critical termination cannot occur.

Every starting point with forward lifetime at most $r$ belongs to a
layer of one of the following forms, at an actual collision or flight:
\begin{equation}\label{eq:v83-layer-input}
 c_j\le C_Eq_0^{d_j}\sqrt r,\quad
 \Delta_j\le C_Eq_0^{d'_j}\sqrt r,\quad
 \operatorname{dist}(T_R^jx,\partial Y_R^*)\le C_Eq_0^{d_j}\sqrt r.
\end{equation}
The physical coarea covariance, proved in the appendix in the convention
used here, has Jacobian $J_h=(\omega(\Phi_hx)/\omega(x))\chi(x)/\chi(\Phi_hx)$.
Writing $E_h^\varepsilon$ for all weighted starting exit mass and
\[
 O_h^\varepsilon(t)=\int_{S_h(t)} W^\varepsilon(x)
            [1-\omega(\Phi_hx)/\omega(x)]_+\,d\mu_t(x),
\]
the complete source satisfies, for each fixed $r$ and almost every $t$,
\begin{equation}\label{eq:v83-shift-input}
 b^\varepsilon(t)\le2A_rb^\varepsilon(t)+\bar E_r^\varepsilon(t)
       +\bar O_r^\varepsilon(t)
       +C_\Theta(1+\varepsilon^{-1})\sqrt r\,A_rp(t),
 \qquad A_rf(t)=\frac1{2r}\int_{t-r}^{t+r}f(s)\,ds.
\end{equation}
\end{inputstatement}
These are the endpoint-field, all-contact motion, uniform cubic-action
and separated-volume results in retained modules 174, 177, 180--182.
The guard term in \eqref{eq:v83-shift-input} is the complete product
guard, not a protective subword source. The derivative and continuation
assertions are confined to regular physical trajectories.

\begin{inputstatement}[Complete local source estimates]
\label{input:v83-local}
For every fixed auxiliary band $B_*$, constants $C_{B_*}<\infty$ and
$\rho_{B_*}<1$ give the following bounds uniformly in radius, exact
labels, interval $J$ and mark $0\le j\le m$. Let $S_{R,s}$ be the
union of strips of width $Cs$ about the finitely many coordinate lines
supporting $\partial Y_R^*$. For $s$ small,
\begin{align}
 m^2\nu\{K_m^{\rm c}=k,A_m=n,L_m\in J,
                       T_R^jx\in S_{R,s}\}
 &\le Cs^{1/16}(|J|+B_*^{-1})
   +C_{B_*}(1+B_*|J|)m^2\rho_{B_*}^{m/2},\label{eq:v83-strip-input}\\
 m^2\int_Jb^\delta(t)\,dt
 &\le C\delta^{1/16}(|J|+B_*^{-1})
   +C_{B_*}(1+B_*|J|)m^3\rho_{B_*}^{m/2}.\label{eq:v83-guard-input}
\end{align}
The local extension to a fixed remaining compact range of $\delta$
is included. At auxiliary parameter one $W^1=1$, so $b^1=p$.

Put $z=|g|$, $\alpha=129/128$ and, for $0<\eta<1/2$, define the
actual nonnegative source measures
\begin{align}
 v^{\eta,b}(t)\,dt&=L_\#(I_m z^{-\eta}\one_{\{0<z<b\}}
                                       \one_\Pi\nu_R^*),\notag\\
 v^\eta(t)\,dt&=L_\#(I_m z^{-\eta}\one_\Pi\nu_R^*).
                                                   \label{eq:v83-fractional-input}
\end{align}
They are finite and absolutely continuous at every fixed count, and
for $0<b<b_1$,
\begin{equation}\label{eq:v83-fractional-finite-input}
 m^2\int_Jv^{\eta,b}(t)\,dt
 \le C_\eta b^{\alpha-\eta}(|J|+B_*^{-1})
  +C_{B_*,\eta}(1+B_*|J|)m^2b^{1/2-\eta}\rho_{B_*}^{m/4}.
\end{equation}
Adding the bounded-gradient part gives, for every fixed $L>0$,
\begin{equation}\label{eq:v83-fixed-width-input}
 \limsup_m\sup_{R,n,k}\sup_{|J|=L}m^2\int_Jv^\eta\le C_\eta L,
 \qquad
 \limsup_m\sup_{R,n,k}\sup_{|J|=L}m^2\int_Jv^{\eta,b}
                                      \le C_\eta b^{\alpha-\eta}L.
\end{equation}
Every auxiliary band is fixed before the collision limit and enlarged
only afterwards; leading constants are independent of that band and
of the previously fixed interval length.
\end{inputstatement}
Equation~\eqref{eq:v83-strip-input} is Supplement
Theorem~\ref{supp-thm:v48-marked-strip-local}. The physical source
bound \eqref{eq:v83-guard-input} is Supplement
Theorem~\ref{supp-thm:v49-physical-remainder-local}, in particular its
finite-window estimate. The fractional bounds are the positive
layer-cake conclusions of retained module 187, stated in full above.
The joint spectral and common-space assumptions used there are recorded
in the appendix. No estimate for an untruncated first inverse gradient
or for a singular roof convolution is part of this input.

\begin{inputstatement}[The scalar positive remainder]
\label{input:v83-scalar}
For large fixed $B$, with $\varepsilon(B)=A_0B^{-1/12}$,
\begin{equation}\label{eq:v76-remainder}
 m^2p-\mathcal L=m^2b^{\varepsilon(B)}+e_B,\qquad
 \limsup_m\sup_{R,n,k}\esssup_{\mathcal C(M)}|e_B|
                                      \le C_MB^{-1/192}.
\end{equation}
The source $b^\varepsilon$ is nonnegative and has a compactly supported
$BV$ density at fixed count. No collision-uniform variation bound is
assumed. Supplement Theorem~\ref{supp-thm:v51-positive-error} gives
this remainder with the same complete finite-peak reference.
\end{inputstatement}

\input{core/190_untruncated_optical_potentials}
\input{core/191_complete_first_exit_surface_mass}
\input{core/192_ordered_fractional_tail_and_endpoint}

\section{The scalar objective and downstream normalization}
The organizing assertion remains \eqref{eq:v76-target}. The complete
optical norm and first-exit surface mass are now quantitatively bounded
without an initial critical cutoff. They are not essential-height
bounds, and the two height terms in \eqref{eq:v83-endpoint} remain
distinct proof obligations. In particular a first-exit surface is not
removed by an internal signed occupation difference, and a fixed-width
volume theorem is not used at a collision-dependent microscopic width.

After scalar closure, the positive path-remainder equality in
Supplement Theorem~\ref{supp-thm:v73-scalar-path-equality} supplies
the same-roof path numerator. Ratios are formed only on positive
arithmetic reference classes; the actual-return clock and forward
likelihood retain their stated inherited hypotheses. Nothing in the
new estimates divides by an arithmetic zero class or reverses this
scalar-first order.

\appendix
\input{appendices/v83_common_kernels_and_dependencies}

\endgroup
\part{Retained critical-source criterion}
\begingroup
\let\ArchiveLabel\label
\let\ArchiveRef\ref
\renewcommand{\label}[1]{\ArchiveLabel{v82archive:#1}}
\renewcommand{\ref}[1]{%
  \IfBeginWith{#1}{supp-}{%
    \StrBehind{#1}{supp-}[\ArchiveTarget]%
    \ArchiveRef{\ArchiveTarget}%
  }{\ArchiveRef{v82archive:#1}}%
}
\makeatletter
\let\ltx@label\label
\makeatother
\paragraph{Retained revision-82 abstract.}
\paragraph{Source clarification for the retained dyadic formula.}
The following historical article is reproduced without editing its source.
In its module 189, the intervals named $I_j(t)$ are absolute roof intervals,
whereas the integrand $v(t+s)$ uses an offset. The active revision resolves
this notation by integrating $s$ from $\ell_j/2$ to $\ell_j$, or equivalently
$v(u)$ over $(t+\ell_j/2,t+\ell_j]$. Its full microscopic tail is formed
before the collision limsup. See the active article's ordered-tail theorem;
no historical smallness condition is promoted by this clarification.

We retain the pointwise local limit problem for the full displacement,
collision-count, occupation and roof record of actual returns in a
triangular Lorentz gas. Joint localization of the two endpoint momenta
and one internal incidence yields a superlinear gain in both widths.
The proof distinguishes a small initial vector from a bounded
multiplier, and treats both chronological cases at every continued
arithmetic peak. A physical cap on the finite-count remainder permits
all inverse-incidence scales to be integrated. The resulting complete
optical source has a critical-width gain greater than one, actual
untruncated fractional inverse-gradient moments, and first
inverse-gradient bounds uniform in a fixed positive lower cutoff.
Applied on common exact-label coarea fibers, these estimates replace
the earlier optical and grazing-exit bound $Cr/a$ by $Cr$, and pay
the same uniform compression cost at nongrazing exits. The complete
critical initial defect and outgoing nongrazing flux remain the two
roof-height estimates needed for the original pointwise theorem;
no exchange of critical-cutoff and collision limits is inferred.



\section{The unchanged physical problem}
Put $R_-=9/20$ and $R_+=47/100$, and let $R\in I=[R_-,R_+]$. Circular scatterers have centers in
$\Lambda=\Z(1,0)+\Z(1/2,\sqrt3/2)$. Write $T_R$ for the regular
collision map and $(\alpha,p)\in(\R/2\pi\Z)\times(-1,1)$ for its
coordinates; $p$ is the outgoing tangential velocity and
$\dd\nu=\dd\alpha\,\dd p/(4\pi)$. Physical flights satisfy
$\tau_-\le\tau_R\le\tau_+$, where $\tau_-=3/50$ and $\tau_+$
is the fixed finite-horizon bound.

The section is specified here to fix every event in the argument.
Put $h_0=1/200$. Let $Y_R$ be the union of the four open coordinate
squares of half-side $h_0$ centered at
\[
 (0,0),\quad(\pi/6,-1/2),\quad
 (-\theta_R,\sin\theta_R),\quad
 (\pi/3-\theta_R,\sin\theta_R).
\]
Here $\theta_R=2\arctan q_R$, and $q_R$ is the unique root in
$(49/100,53/100)$ of
\[
 8Rq(1+q^2)-4q(1-q^2)+\sqrt3(1-6q^2+q^4)=0.
\]
Set
\[
 U=\{|\alpha-\pi/6|<3/50,\ |p|<3/20\},\qquad
 Y_R^*=Y_R\cup U,\qquad c_*=\frac{91}{10000\pi}.
\]
The rectangles are disjoint and $\nu(Y_R^*)=c_*$. We always use
$\nu_R^*=c_*^{-1}\nu|_{Y_R^*}$, not a renormalization of a smaller
source.

At $m$ flights, let $L=L_{m,R}$ be their total length,
$K^{\rm c}_{m,R}$ their lattice displacement, and
$A_{m,R}=\sum_{j=0}^{m-1}\one_{Y_R^*}\circ T_R^j$. The event is
\begin{equation}\label{eq:v76-event}
 \Pi_{n,k,m,R}=Y_R^*\cap\{K^{\rm c}_{m,R}=k,\ A_{m,R}=n\}
                         \cap T_R^{-m}Y_R^*.
\end{equation}
The initial visit is counted; the terminal visit is tested separately.
On this event the actual $n$th return has collision count $m$ and
roof $L$. Denote its density by $p_{n,R}(k,m,t)$.

Write $c_j=\sqrt{1-p_j^2}$ and let $\Delta_j$ be the least clearance
of flight $j$ from any third disk, excluding its two incident disks.
Keep the original smooth function $0\le\Theta\le1$, zero on
$(-\infty,1]$ and one on $[2,\infty)$. With $q_0=47/53$,
$d_j=\min(j,m-j)$ and $d'_j=\min(j,m-1-j)$, set
\begin{align}
 H^\varepsilon&=
 \prod_{j=0}^{m}\Theta\!\left(\frac{c_j}{\varepsilon q_0^{d_j/4}}\right)
 \prod_{j=0}^{m-1}\Theta\!\left(\frac{\Delta_j}{\varepsilon q_0^{d'_j/4}}\right),
 \qquad W^\varepsilon=1-H^\varepsilon,\label{eq:v76-guard}\\
 b^\varepsilon(t)\dd t
   &=L_\#(W^\varepsilon\one_{\Pi_{n,k,m,R}}\nu_R^*).
 \label{eq:v76-source}
\end{align}
In particular first incidence and first clearance outside any selected
critical neighborhoods are still present. No section-distance factor
has been inserted into this guard.

\subsection{The arithmetic reference and the target}
For $y=(k_1,k_2,t,n)$ put
\[
 \mu_R=(0,0,\bar\tau_R,c_*),\quad
 Z=(y-m\mu_R)/\sqrt m,\quad
 \mathcal C_{m,R,n,k}(M)=\{t:|Z|\le M\}.
\]
We retain the finite spectral-peak reference, including peaks close to
but not on the unit circle. To specify its meaning, the collision
record observable is $(\kappa_1,\kappa_2,\tau,\one_{Y_R^*})$.
Let $\lambda_j(R,\vartheta)$ and $\mathsf P_j(R,\vartheta)$ be its
continued simple twisted eigenvalue and projection in each of the
fixed peak charts, and let $\chi_j$ be their frequency partition.
At the real maximum $a_j$ of $\log|\lambda_j|$ put
\[
 z_j=\lambda_j(R,a_j),\quad
 \mu_j=i^{-1}\nabla\log\lambda_j(R,a_j),\quad
 H_j=-\nabla^2\log\lambda_j(R,a_j),
\]
with $\operatorname{Re}H_j\ge aI_4$ for a common $a>0$. If
$\eta_R=\one_{Y_R^*}$ and $\ell$ denotes mass evaluation, put
$B_j=\chi_j(R,a_j)\ell(M_{\eta_R}\mathsf P_j(R,a_j)(\eta_R\nu))$.
Define
\begin{align}
 G_H(z)&=(2\pi)^{-4}\int_{\R^4}e^{-iv\cdot z-v^{\mathsf T}Hv/2}\dd v,
 \notag\\
 \mathcal L_{m,R}(y)&=\operatorname{Re}\left[
 c_*^{-1}\sum_j e^{-ia_j\cdot y}z_j^m B_j
      G_{H_j}\left((y-m\mu_j)/\sqrt m\right)\right].
 \label{eq:v76-reference}
\end{align}
Only the charts active at $R$ are summed. This is the same transition
reference as in the retained spectral theorem. The existence and
uniform estimates for these scalar data belong to that analytic
input, not to the transport argument below. No frequency branch or
arithmetic zero class is deleted.

The governing assertion remains
\begin{equation}\label{eq:v76-target}
 \sup_{R,\,1\le n\le m,\,k\in\Z^2}
 \esssup_{t\in\mathcal C_{m,R,n,k}(M)}
 |m^2p_{n,R}(k,m,t)-\mathcal L_{m,R}(k_1,k_2,t,n)|\longrightarrow0.
\end{equation}
Let $\mathcal E_M$ denote this error's collision limsup. An empty
central target contributes zero. The order of the uniformity is
$\sup_R\esssup_t$, so exceptional density sets may depend on $R$.

\subsection{The analytic inputs in the form used here}
\begin{inputstatement}\label{input:v76-analytic}
The unchanged source \eqref{eq:v76-source} has a nonnegative,
compactly supported $BV$ density at every fixed count. For every
fixed $\varepsilon,\delta>0$, with
$A_\delta f(t)=(2\delta)^{-1}\int_{t-\delta}^{t+\delta}f(s)\dd s$,
\begin{equation}\label{eq:v76-average}
 \limsup_m\sup_{R,n,k}m^2\|A_\delta b^\varepsilon\|_\infty
                                      \le C\varepsilon^{1/16}.
\end{equation}
The constant is independent of the previously fixed width. For large
fixed $B$ and $\varepsilon(B)=A_0B^{-1/12}$,
\begin{equation}\label{eq:v76-remainder}
 m^2p-\mathcal L=m^2b^{\varepsilon(B)}+e_B,\qquad
 \limsup_m\sup_{R,n,k}\esssup_{\mathcal C(M)}|e_B|
                                     \le C_MB^{-1/192}.
\end{equation}
All auxiliary spectral bands in these statements are fixed before
the collision limit and enlarged only afterwards.
\end{inputstatement}
These are the complete-source tame theorem, the fixed-width
complete-source bound, and the positive scalar remainder theorem in
the technical supplement. Their proofs and historical hypotheses are
retained there. The fixed-count source uses threshold integration
\cite{CMInt}, definable monotonicity \cite{Coste}, and the
distributional characterization of $BV$ \cite{AFP}; it does not
assume definability of the original smooth guard. They are restated as inputs to distinguish them from
the new geometric and transport proofs. No long-time bound on a
fixed-count variation constant is part of this statement.



\subsection{Precise inherited inputs}\label{sec:v79-input-map}
The fixed-width estimate is Supplement Theorem~\ref{supp-thm:v49-physical-remainder-local},
in particular Supplement equation~\eqref{supp-eq:v49-local-limsup}.
It uses Supplement Theorem~\ref{supp-thm:v49-physical-marked-local}, on the
original exact displacement and occupation coefficient. For an interval
of length $2r$, division by $2r$ gives \eqref{eq:v76-average}. That
proof fixes the auxiliary band, takes the collision limit, and only
then enlarges the auxiliary band. Its leading constant is independent
of the previously fixed interval length. The extension from small
margins to a fixed remaining compact range is stated there explicitly.
At auxiliary parameter one, $\Theta(c_0)=0$ since $c_0\le1$, so $W^1=1$
and $b^1=p$. This supplies the unweighted fixed-width estimate without
asserting a raw pointwise bound. The positive remainder is the inherited
scalar input \eqref{eq:v76-remainder}, stated as Supplement
Theorem~\ref{supp-thm:v51-positive-error}; no surface-current estimate is
included in any of these volume-source inputs.

\subsection{The complete-source estimate proved here}
The physical endpoint variables are
$q=(\alpha_0,\alpha_m)$, $g=\nabla_qF=R(-p_0,p_m)$ and
$K=D_q^2F$. The retained canonical field satisfies $XF=1$,
$X\arg g=0$ and
$|g(\Phi_sx)|^2\ge |g(x)|^2+2\lambda_\circ s$ up to its true
lifetime. Its physical Jacobian factors into an inverse-twist ratio
and an angular Hessian ratio. The latter and the original guard
already have complete-source estimates.

The remaining optical drift is bounded by
$C\omega I_m/|g|$, where
\[
 I_m=1+\sum_{j=1}^{m-1}q_0^{d_j}/c_j.
\]
The preceding singular-insertion argument controlled $I_m$ but
replaced $|g|^{-1}$ by $a^{-1}$ on initial $|g|\ge a$.
The new argument keeps both components of $g$ in the spectral
pairing. With $\alpha=129/128$, it proves, for the complete original
source and each fixed roof width,
\[
 u^b(t)\,dt=L_\#(I_m\one_{\{|g|<b\}}\one_\Pi\nu_R^*),\qquad
 \limsup_m\sup_{R,n,k}m^2\|A_ru^b\|_\infty\le Cb^\alpha.
\]

Theorem~\ref{thm:v81-gradient-insertions} integrates this bound with
its finite-count remainder. It gives an actual untruncated mixed
insertion for powers $|g|^{-\eta}$ with $0<\eta<1/2$, and a bound
for $|g|^{-1}\one_{\{|g|\ge a\}}$ whose leading constant is
independent of the previously fixed $a>0$.
Theorem~\ref{thm:v81-optical-height} then proves a complete
$Cr$ optical and grazing starting-mass estimate on that initial
source. Nongrazing exits retain their exact outgoing flux; only
the paid pullback cost is improved. These statements keep the
original target \eqref{eq:v76-target}, rather than replacing it by
a local law for a different record.

\input{core/174_canonical_endpoint_momentum_flow}
\input{core/177_all_contact_motion_and_guard}
\input{core/180_uniform_cubic_reduced_action}
\input{core/181_separated_physical_volume_budget}
\input{core/182_optical_sensitivity_and_exit_flux}
\input{core/183_homogeneous_incidence_localization}
\input{core/184_singular_optical_transport}
\input{core/185_common_coarea_and_remaining_boundary}
\input{core/186_joint_endpoint_incidence_localization}
\input{core/187_mixed_critical_volume_insertions}
\input{core/188_uniform_critical_truncation_transport}

\section{The remaining scalar estimates and their consequences}
Equation~\eqref{eq:v81-endpoint-budget} is the revised error budget
for the original pointwise assertion. It retains the complete
critical initial defect and the nongrazing outgoing flux. The
uniformity of the new estimate removes the former restriction
$r/a\to0$ from the paid drift; it does not remove those two sources.
In particular the inverse-Jacobian cost has been estimated rather
than erased from the first-exit formula.

The quantitative distinction between a truncated and an untruncated
first inverse gradient is explicit in
\eqref{eq:v81-truncated-finite}. The actual fractional insertion
\eqref{eq:v81-fractional-finite} integrates both singular parameters
before the collision limit, whereas the first inverse power still
requires its positive cutoff to be fixed. The full pointwise theorem
is not asserted until the two remaining roof-height bounds are
proved in the stated order. All arithmetic peaks and zero classes
remain in the same reference.

After scalar closure, Supplement
Theorem~\ref{supp-thm:v73-scalar-path-equality} and the positive path
remainder give the same-roof path numerator. Normalization is only
on positive arithmetic reference classes. Collision-to-return clock
transfer and forward likelihood retain their separate inherited
hypotheses. The complete preceding article and every inherited proof
are retained in the supplement.

\begin{quote}\small
\textbf{Superseded dyadic notation.}
The next module is retained as a historical source, not as an active
reference for its dyadic integration formula. The active convention is
\[
 \ell_j=2^{-j}r,\qquad
 \int_{\ell_j/2}^{\ell_j}v(t+s)\,ds
   =\int_{t+\ell_j/2}^{t+\ell_j}v(u)\,du.
\]
The variable in the first integral is an offset. The complete microscopic
tail is formed before the collision limsup. This corrected convention,
used in the active article, supersedes the ambiguous historical display;
no decay conclusion is inferred from this editorial correction.
\end{quote}
\input{core/189_dyadic_critical_source_criterion}


\endgroup
\part{Joint endpoint localization and critical optical insertions}
\begingroup
\let\ArchiveLabel\label
\let\ArchiveRef\ref
\renewcommand{\label}[1]{\ArchiveLabel{v81archive:#1}}
\renewcommand{\ref}[1]{%
  \IfBeginWith{#1}{supp-}{%
    \StrBehind{#1}{supp-}[\ArchiveTarget]%
    \ArchiveRef{\ArchiveTarget}%
  }{\ArchiveRef{v81archive:#1}}%
}
\makeatletter
\let\ltx@label\label
\makeatother
\paragraph{Retained revision-81 abstract.}

We retain the pointwise local limit problem for the full displacement,
collision-count, occupation and roof record of actual returns in a
triangular Lorentz gas. Joint localization of the two endpoint momenta
and one internal incidence yields a superlinear gain in both widths.
The proof distinguishes a small initial vector from a bounded
multiplier, and treats both chronological cases at every continued
arithmetic peak. A physical cap on the finite-count remainder permits
all inverse-incidence scales to be integrated. The resulting complete
optical source has a critical-width gain greater than one, actual
untruncated fractional inverse-gradient moments, and first
inverse-gradient bounds uniform in a fixed positive lower cutoff.
Applied on common exact-label coarea fibers, these estimates replace
the earlier optical and grazing-exit bound $Cr/a$ by $Cr$, and pay
the same uniform compression cost at nongrazing exits. The complete
critical initial defect and outgoing nongrazing flux remain the two
roof-height estimates needed for the original pointwise theorem;
no exchange of critical-cutoff and collision limits is inferred.



\section{The unchanged physical problem}
Put $R_-=9/20$ and $R_+=47/100$, and let $R\in I=[R_-,R_+]$. Circular scatterers have centers in
$\Lambda=\Z(1,0)+\Z(1/2,\sqrt3/2)$. Write $T_R$ for the regular
collision map and $(\alpha,p)\in(\R/2\pi\Z)\times(-1,1)$ for its
coordinates; $p$ is the outgoing tangential velocity and
$\dd\nu=\dd\alpha\,\dd p/(4\pi)$. Physical flights satisfy
$\tau_-\le\tau_R\le\tau_+$, where $\tau_-=3/50$ and $\tau_+$
is the fixed finite-horizon bound.

The section is specified here to fix every event in the argument.
Put $h_0=1/200$. Let $Y_R$ be the union of the four open coordinate
squares of half-side $h_0$ centered at
\[
 (0,0),\quad(\pi/6,-1/2),\quad
 (-\theta_R,\sin\theta_R),\quad
 (\pi/3-\theta_R,\sin\theta_R).
\]
Here $\theta_R=2\arctan q_R$, and $q_R$ is the unique root in
$(49/100,53/100)$ of
\[
 8Rq(1+q^2)-4q(1-q^2)+\sqrt3(1-6q^2+q^4)=0.
\]
Set
\[
 U=\{|\alpha-\pi/6|<3/50,\ |p|<3/20\},\qquad
 Y_R^*=Y_R\cup U,\qquad c_*=\frac{91}{10000\pi}.
\]
The rectangles are disjoint and $\nu(Y_R^*)=c_*$. We always use
$\nu_R^*=c_*^{-1}\nu|_{Y_R^*}$, not a renormalization of a smaller
source.

At $m$ flights, let $L=L_{m,R}$ be their total length,
$K^{\rm c}_{m,R}$ their lattice displacement, and
$A_{m,R}=\sum_{j=0}^{m-1}\one_{Y_R^*}\circ T_R^j$. The event is
\begin{equation}\label{eq:v76-event}
 \Pi_{n,k,m,R}=Y_R^*\cap\{K^{\rm c}_{m,R}=k,\ A_{m,R}=n\}
                         \cap T_R^{-m}Y_R^*.
\end{equation}
The initial visit is counted; the terminal visit is tested separately.
On this event the actual $n$th return has collision count $m$ and
roof $L$. Denote its density by $p_{n,R}(k,m,t)$.

Write $c_j=\sqrt{1-p_j^2}$ and let $\Delta_j$ be the least clearance
of flight $j$ from any third disk, excluding its two incident disks.
Keep the original smooth function $0\le\Theta\le1$, zero on
$(-\infty,1]$ and one on $[2,\infty)$. With $q_0=47/53$,
$d_j=\min(j,m-j)$ and $d'_j=\min(j,m-1-j)$, set
\begin{align}
 H^\varepsilon&=
 \prod_{j=0}^{m}\Theta\!\left(\frac{c_j}{\varepsilon q_0^{d_j/4}}\right)
 \prod_{j=0}^{m-1}\Theta\!\left(\frac{\Delta_j}{\varepsilon q_0^{d'_j/4}}\right),
 \qquad W^\varepsilon=1-H^\varepsilon,\label{eq:v76-guard}\\
 b^\varepsilon(t)\dd t
   &=L_\#(W^\varepsilon\one_{\Pi_{n,k,m,R}}\nu_R^*).
 \label{eq:v76-source}
\end{align}
In particular first incidence and first clearance outside any selected
critical neighborhoods are still present. No section-distance factor
has been inserted into this guard.

\subsection{The arithmetic reference and the target}
For $y=(k_1,k_2,t,n)$ put
\[
 \mu_R=(0,0,\bar\tau_R,c_*),\quad
 Z=(y-m\mu_R)/\sqrt m,\quad
 \mathcal C_{m,R,n,k}(M)=\{t:|Z|\le M\}.
\]
We retain the finite spectral-peak reference, including peaks close to
but not on the unit circle. To specify its meaning, the collision
record observable is $(\kappa_1,\kappa_2,\tau,\one_{Y_R^*})$.
Let $\lambda_j(R,\vartheta)$ and $\mathsf P_j(R,\vartheta)$ be its
continued simple twisted eigenvalue and projection in each of the
fixed peak charts, and let $\chi_j$ be their frequency partition.
At the real maximum $a_j$ of $\log|\lambda_j|$ put
\[
 z_j=\lambda_j(R,a_j),\quad
 \mu_j=i^{-1}\nabla\log\lambda_j(R,a_j),\quad
 H_j=-\nabla^2\log\lambda_j(R,a_j),
\]
with $\operatorname{Re}H_j\ge aI_4$ for a common $a>0$. If
$\eta_R=\one_{Y_R^*}$ and $\ell$ denotes mass evaluation, put
$B_j=\chi_j(R,a_j)\ell(M_{\eta_R}\mathsf P_j(R,a_j)(\eta_R\nu))$.
Define
\begin{align}
 G_H(z)&=(2\pi)^{-4}\int_{\R^4}e^{-iv\cdot z-v^{\mathsf T}Hv/2}\dd v,
 \notag\\
 \mathcal L_{m,R}(y)&=\operatorname{Re}\left[
 c_*^{-1}\sum_j e^{-ia_j\cdot y}z_j^m B_j
      G_{H_j}\left((y-m\mu_j)/\sqrt m\right)\right].
 \label{eq:v76-reference}
\end{align}
Only the charts active at $R$ are summed. This is the same transition
reference as in the retained spectral theorem. The existence and
uniform estimates for these scalar data belong to that analytic
input, not to the transport argument below. No frequency branch or
arithmetic zero class is deleted.

The governing assertion remains
\begin{equation}\label{eq:v76-target}
 \sup_{R,\,1\le n\le m,\,k\in\Z^2}
 \esssup_{t\in\mathcal C_{m,R,n,k}(M)}
 |m^2p_{n,R}(k,m,t)-\mathcal L_{m,R}(k_1,k_2,t,n)|\longrightarrow0.
\end{equation}
Let $\mathcal E_M$ denote this error's collision limsup. An empty
central target contributes zero. The order of the uniformity is
$\sup_R\esssup_t$, so exceptional density sets may depend on $R$.

\subsection{The analytic inputs in the form used here}
\begin{inputstatement}\label{input:v76-analytic}
The unchanged source \eqref{eq:v76-source} has a nonnegative,
compactly supported $BV$ density at every fixed count. For every
fixed $\varepsilon,\delta>0$, with
$A_\delta f(t)=(2\delta)^{-1}\int_{t-\delta}^{t+\delta}f(s)\dd s$,
\begin{equation}\label{eq:v76-average}
 \limsup_m\sup_{R,n,k}m^2\|A_\delta b^\varepsilon\|_\infty
                                      \le C\varepsilon^{1/16}.
\end{equation}
The constant is independent of the previously fixed width. For large
fixed $B$ and $\varepsilon(B)=A_0B^{-1/12}$,
\begin{equation}\label{eq:v76-remainder}
 m^2p-\mathcal L=m^2b^{\varepsilon(B)}+e_B,\qquad
 \limsup_m\sup_{R,n,k}\esssup_{\mathcal C(M)}|e_B|
                                     \le C_MB^{-1/192}.
\end{equation}
All auxiliary spectral bands in these statements are fixed before
the collision limit and enlarged only afterwards.
\end{inputstatement}
These are the complete-source tame theorem, the fixed-width
complete-source bound, and the positive scalar remainder theorem in
the technical supplement. Their proofs and historical hypotheses are
retained there. The fixed-count source uses threshold integration
\cite{CMInt}, definable monotonicity \cite{Coste}, and the
distributional characterization of $BV$ \cite{AFP}; it does not
assume definability of the original smooth guard. They are restated as inputs to distinguish them from
the new geometric and transport proofs. No long-time bound on a
fixed-count variation constant is part of this statement.



\subsection{Precise inherited inputs}\label{sec:v79-input-map}
The fixed-width estimate is Supplement Theorem~\ref{supp-thm:v49-physical-remainder-local},
in particular Supplement equation~\eqref{supp-eq:v49-local-limsup}.
It uses Supplement Theorem~\ref{supp-thm:v49-physical-marked-local}, on the
original exact displacement and occupation coefficient. For an interval
of length $2r$, division by $2r$ gives \eqref{eq:v76-average}. That
proof fixes the auxiliary band, takes the collision limit, and only
then enlarges the auxiliary band. Its leading constant is independent
of the previously fixed interval length. The extension from small
margins to a fixed remaining compact range is stated there explicitly.
At auxiliary parameter one, $\Theta(c_0)=0$ since $c_0\le1$, so $W^1=1$
and $b^1=p$. This supplies the unweighted fixed-width estimate without
asserting a raw pointwise bound. The positive remainder is the inherited
scalar input \eqref{eq:v76-remainder}, stated as Supplement
Theorem~\ref{supp-thm:v51-positive-error}; no surface-current estimate is
included in any of these volume-source inputs.

\subsection{The complete-source estimate proved here}
The physical endpoint variables are
$q=(\alpha_0,\alpha_m)$, $g=\nabla_qF=R(-p_0,p_m)$ and
$K=D_q^2F$. The retained canonical field satisfies $XF=1$,
$X\arg g=0$ and
$|g(\Phi_sx)|^2\ge |g(x)|^2+2\lambda_\circ s$ up to its true
lifetime. Its physical Jacobian factors into an inverse-twist ratio
and an angular Hessian ratio. The latter and the original guard
already have complete-source estimates.

The remaining optical drift is bounded by
$C\omega I_m/|g|$, where
\[
 I_m=1+\sum_{j=1}^{m-1}q_0^{d_j}/c_j.
\]
The preceding singular-insertion argument controlled $I_m$ but
replaced $|g|^{-1}$ by $a^{-1}$ on initial $|g|\ge a$.
The new argument keeps both components of $g$ in the spectral
pairing. With $\alpha=129/128$, it proves, for the complete original
source and each fixed roof width,
\[
 u^b(t)\,dt=L_\#(I_m\one_{\{|g|<b\}}\one_\Pi\nu_R^*),\qquad
 \limsup_m\sup_{R,n,k}m^2\|A_ru^b\|_\infty\le Cb^\alpha.
\]

Theorem~\ref{thm:v81-gradient-insertions} integrates this bound with
its finite-count remainder. It gives an actual untruncated mixed
insertion for powers $|g|^{-\eta}$ with $0<\eta<1/2$, and a bound
for $|g|^{-1}\one_{\{|g|\ge a\}}$ whose leading constant is
independent of the previously fixed $a>0$.
Theorem~\ref{thm:v81-optical-height} then proves a complete
$Cr$ optical and grazing starting-mass estimate on that initial
source. Nongrazing exits retain their exact outgoing flux; only
the paid pullback cost is improved. These statements keep the
original target \eqref{eq:v76-target}, rather than replacing it by
a local law for a different record.

\input{core/174_canonical_endpoint_momentum_flow}
\input{core/177_all_contact_motion_and_guard}
\input{core/180_uniform_cubic_reduced_action}
\input{core/181_separated_physical_volume_budget}
\input{core/182_optical_sensitivity_and_exit_flux}
\input{core/183_homogeneous_incidence_localization}
\input{core/184_singular_optical_transport}
\input{core/185_common_coarea_and_remaining_boundary}
\input{core/186_joint_endpoint_incidence_localization}
\input{core/187_mixed_critical_volume_insertions}
\input{core/188_uniform_critical_truncation_transport}

\section{The remaining scalar estimates and their consequences}
Equation~\eqref{eq:v81-endpoint-budget} is the revised error budget
for the original pointwise assertion. It retains the complete
critical initial defect and the nongrazing outgoing flux. The
uniformity of the new estimate removes the former restriction
$r/a\to0$ from the paid drift; it does not remove those two sources.
In particular the inverse-Jacobian cost has been estimated rather
than erased from the first-exit formula.

The quantitative distinction between a truncated and an untruncated
first inverse gradient is explicit in
\eqref{eq:v81-truncated-finite}. The actual fractional insertion
\eqref{eq:v81-fractional-finite} integrates both singular parameters
before the collision limit, whereas the first inverse power still
requires its positive cutoff to be fixed. The full pointwise theorem
is not asserted until the two remaining roof-height bounds are
proved in the stated order. All arithmetic peaks and zero classes
remain in the same reference.

After scalar closure, Supplement
Theorem~\ref{supp-thm:v73-scalar-path-equality} and the positive path
remainder give the same-roof path numerator. Normalization is only
on positive arithmetic reference classes. Collision-to-return clock
transfer and forward likelihood retain their separate inherited
hypotheses. The complete preceding article and every inherited proof
are retained in the supplement.


\endgroup
\part{Singular insertions and common exact-label fibers}
\begingroup
\let\ArchiveLabel\label
\let\ArchiveRef\ref
\renewcommand{\label}[1]{\ArchiveLabel{v80archive:#1}}
\renewcommand{\ref}[1]{%
  \IfBeginWith{#1}{supp-}{%
    \StrBehind{#1}{supp-}[\ArchiveTarget]%
    \ArchiveRef{\ArchiveTarget}%
  }{\ArchiveRef{v80archive:#1}}%
}
\makeatletter
\let\ltx@label\label
\makeatother
\paragraph{Retained revision-80 abstract.}

We study the pointwise local limit problem for the complete displacement,
collision-count, occupation and roof record of actual returns in a
triangular Lorentz gas. Homogeneous grazing strips give a superlinear
weak estimate for real-frequency physical powers. Combining it with
a fixed refinement of the spectral peaks and a physical cap on the
remainder yields a marked-incidence bound which can be integrated
against inverse incidence. This proves a complete fixed-width local
estimate for unbounded insertions, without assuming that they are
bounded strong-space multipliers. Applied to the canonical optical
flow, it controls inverse-twist loss and the initial mass of genuine
grazing exits at every internal incidence scale, on sources with
noncritical initial endpoint momentum. The same drift estimate pays
the compression cost in the nongrazing first-exit pullback. Common
coarea kernels and a flow-box exhaustion keep all exact labels and
true boundaries visible. The original scalar endpoint is reduced to
the critical initial source and the remaining nongrazing boundary
flux; their collision-uniform height estimates are not inferred from
fixed-width volume concentration.



\section{The unchanged physical problem}
Put $R_-=9/20$ and $R_+=47/100$, and let $R\in I=[R_-,R_+]$. Circular scatterers have centers in
$\Lambda=\Z(1,0)+\Z(1/2,\sqrt3/2)$. Write $T_R$ for the regular
collision map and $(\alpha,p)\in(\R/2\pi\Z)\times(-1,1)$ for its
coordinates; $p$ is the outgoing tangential velocity and
$\dd\nu=\dd\alpha\,\dd p/(4\pi)$. Physical flights satisfy
$\tau_-\le\tau_R\le\tau_+$, where $\tau_-=3/50$ and $\tau_+$
is the fixed finite-horizon bound.

The section is specified here to fix every event in the argument.
Put $h_0=1/200$. Let $Y_R$ be the union of the four open coordinate
squares of half-side $h_0$ centered at
\[
 (0,0),\quad(\pi/6,-1/2),\quad
 (-\theta_R,\sin\theta_R),\quad
 (\pi/3-\theta_R,\sin\theta_R).
\]
Here $\theta_R=2\arctan q_R$, and $q_R$ is the unique root in
$(49/100,53/100)$ of
\[
 8Rq(1+q^2)-4q(1-q^2)+\sqrt3(1-6q^2+q^4)=0.
\]
Set
\[
 U=\{|\alpha-\pi/6|<3/50,\ |p|<3/20\},\qquad
 Y_R^*=Y_R\cup U,\qquad c_*=\frac{91}{10000\pi}.
\]
The rectangles are disjoint and $\nu(Y_R^*)=c_*$. We always use
$\nu_R^*=c_*^{-1}\nu|_{Y_R^*}$, not a renormalization of a smaller
source.

At $m$ flights, let $L=L_{m,R}$ be their total length,
$K^{\rm c}_{m,R}$ their lattice displacement, and
$A_{m,R}=\sum_{j=0}^{m-1}\one_{Y_R^*}\circ T_R^j$. The event is
\begin{equation}\label{eq:v76-event}
 \Pi_{n,k,m,R}=Y_R^*\cap\{K^{\rm c}_{m,R}=k,\ A_{m,R}=n\}
                         \cap T_R^{-m}Y_R^*.
\end{equation}
The initial visit is counted; the terminal visit is tested separately.
On this event the actual $n$th return has collision count $m$ and
roof $L$. Denote its density by $p_{n,R}(k,m,t)$.

Write $c_j=\sqrt{1-p_j^2}$ and let $\Delta_j$ be the least clearance
of flight $j$ from any third disk, excluding its two incident disks.
Keep the original smooth function $0\le\Theta\le1$, zero on
$(-\infty,1]$ and one on $[2,\infty)$. With $q_0=47/53$,
$d_j=\min(j,m-j)$ and $d'_j=\min(j,m-1-j)$, set
\begin{align}
 H^\varepsilon&=
 \prod_{j=0}^{m}\Theta\!\left(\frac{c_j}{\varepsilon q_0^{d_j/4}}\right)
 \prod_{j=0}^{m-1}\Theta\!\left(\frac{\Delta_j}{\varepsilon q_0^{d'_j/4}}\right),
 \qquad W^\varepsilon=1-H^\varepsilon,\label{eq:v76-guard}\\
 b^\varepsilon(t)\dd t
   &=L_\#(W^\varepsilon\one_{\Pi_{n,k,m,R}}\nu_R^*).
 \label{eq:v76-source}
\end{align}
In particular first incidence and first clearance outside any selected
critical neighborhoods are still present. No section-distance factor
has been inserted into this guard.

\subsection{The arithmetic reference and the target}
For $y=(k_1,k_2,t,n)$ put
\[
 \mu_R=(0,0,\bar\tau_R,c_*),\quad
 Z=(y-m\mu_R)/\sqrt m,\quad
 \mathcal C_{m,R,n,k}(M)=\{t:|Z|\le M\}.
\]
We retain the finite spectral-peak reference, including peaks close to
but not on the unit circle. To specify its meaning, the collision
record observable is $(\kappa_1,\kappa_2,\tau,\one_{Y_R^*})$.
Let $\lambda_j(R,\vartheta)$ and $\mathsf P_j(R,\vartheta)$ be its
continued simple twisted eigenvalue and projection in each of the
fixed peak charts, and let $\chi_j$ be their frequency partition.
At the real maximum $a_j$ of $\log|\lambda_j|$ put
\[
 z_j=\lambda_j(R,a_j),\quad
 \mu_j=i^{-1}\nabla\log\lambda_j(R,a_j),\quad
 H_j=-\nabla^2\log\lambda_j(R,a_j),
\]
with $\operatorname{Re}H_j\ge aI_4$ for a common $a>0$. If
$\eta_R=\one_{Y_R^*}$ and $\ell$ denotes mass evaluation, put
$B_j=\chi_j(R,a_j)\ell(M_{\eta_R}\mathsf P_j(R,a_j)(\eta_R\nu))$.
Define
\begin{align}
 G_H(z)&=(2\pi)^{-4}\int_{\R^4}e^{-iv\cdot z-v^{\mathsf T}Hv/2}\dd v,
 \notag\\
 \mathcal L_{m,R}(y)&=\operatorname{Re}\left[
 c_*^{-1}\sum_j e^{-ia_j\cdot y}z_j^m B_j
      G_{H_j}\left((y-m\mu_j)/\sqrt m\right)\right].
 \label{eq:v76-reference}
\end{align}
Only the charts active at $R$ are summed. This is the same transition
reference as in the retained spectral theorem. The existence and
uniform estimates for these scalar data belong to that analytic
input, not to the transport argument below. No frequency branch or
arithmetic zero class is deleted.

The governing assertion remains
\begin{equation}\label{eq:v76-target}
 \sup_{R,\,1\le n\le m,\,k\in\Z^2}
 \esssup_{t\in\mathcal C_{m,R,n,k}(M)}
 |m^2p_{n,R}(k,m,t)-\mathcal L_{m,R}(k_1,k_2,t,n)|\longrightarrow0.
\end{equation}
Let $\mathcal E_M$ denote this error's collision limsup. An empty
central target contributes zero. The order of the uniformity is
$\sup_R\esssup_t$, so exceptional density sets may depend on $R$.

\subsection{The analytic inputs in the form used here}
\begin{inputstatement}\label{input:v76-analytic}
The unchanged source \eqref{eq:v76-source} has a nonnegative,
compactly supported $BV$ density at every fixed count. For every
fixed $\varepsilon,\delta>0$, with
$A_\delta f(t)=(2\delta)^{-1}\int_{t-\delta}^{t+\delta}f(s)\dd s$,
\begin{equation}\label{eq:v76-average}
 \limsup_m\sup_{R,n,k}m^2\|A_\delta b^\varepsilon\|_\infty
                                      \le C\varepsilon^{1/16}.
\end{equation}
The constant is independent of the previously fixed width. For large
fixed $B$ and $\varepsilon(B)=A_0B^{-1/12}$,
\begin{equation}\label{eq:v76-remainder}
 m^2p-\mathcal L=m^2b^{\varepsilon(B)}+e_B,\qquad
 \limsup_m\sup_{R,n,k}\esssup_{\mathcal C(M)}|e_B|
                                     \le C_MB^{-1/192}.
\end{equation}
All auxiliary spectral bands in these statements are fixed before
the collision limit and enlarged only afterwards.
\end{inputstatement}
These are the complete-source tame theorem, the fixed-width
complete-source bound, and the positive scalar remainder theorem in
the technical supplement. Their proofs and historical hypotheses are
retained there. The fixed-count source uses threshold integration
\cite{CMInt}, definable monotonicity \cite{Coste}, and the
distributional characterization of $BV$ \cite{AFP}; it does not
assume definability of the original smooth guard. They are restated as inputs to distinguish them from
the new geometric and transport proofs. No long-time bound on a
fixed-count variation constant is part of this statement.



\subsection{Precise inherited inputs}\label{sec:v79-input-map}
The fixed-width estimate is Supplement Theorem~\ref{supp-thm:v49-physical-remainder-local},
in particular Supplement equation~\eqref{supp-eq:v49-local-limsup}.
It uses Supplement Theorem~\ref{supp-thm:v49-physical-marked-local}, on the
original exact displacement and occupation coefficient. For an interval
of length $2r$, division by $2r$ gives \eqref{eq:v76-average}. That
proof fixes the auxiliary band, takes the collision limit, and only
then enlarges the auxiliary band. Its leading constant is independent
of the previously fixed interval length. The extension from small
margins to a fixed remaining compact range is stated there explicitly.
At auxiliary parameter one, $\Theta(c_0)=0$ since $c_0\le1$, so $W^1=1$
and $b^1=p$. This supplies the unweighted fixed-width estimate without
asserting a raw pointwise bound. The positive remainder is the inherited
scalar input \eqref{eq:v76-remainder}, stated as Supplement
Theorem~\ref{supp-thm:v51-positive-error}; no surface-current estimate is
included in any of these volume-source inputs.

\subsection{Singular insertions and physical optical transport}
The first-variation coordinates are $q=(\alpha_0,\alpha_m)$,
$g=\nabla_qF=R(-p_0,p_m)$, $K=D_q^2F$ and
$\vartheta=\arg(-p_0,p_m)$. The canonical field satisfies
$XF=1$, $X\vartheta=0$, and is stopped at the true boundary of the
regular exact event. Put
$\omega=|\partial_{p_0}\alpha_m|^{-1}$ and
$\chi=\det D_q(F,\vartheta)$. The retained geometric proofs below
give uniform positive bounds for $K$ and $\chi$, a uniform cubic
bound for $F$, and the physical Jacobian factorization
\[
             J_h=\frac{\omega(\Phi_hx)}{\omega(x)}
                       \frac{\chi(x)}{\chi(\Phi_hx)}.
\]
The angular ratio and the original guard have already been controlled
on the complete source. The remaining optical sensitivity contains
$\sum_i q_0^{d_i}/c_i$, which is singular at small internal incidences.

The new estimate for this sensitivity is a local concentration theorem
for the actual unbounded volume insertion
\[
 I_m=1+\sum_{i=1}^{m-1}q_0^{d_i}c_i^{-1},\qquad
 u(t)\,dt=L_\#(I_m\one_{\Pi_{n,k,m,R}}\nu_R^*).
\]
Corollary~\ref{cor:v80-depth-source} proves, for each fixed $r>0$,
\[
              \limsup_m\sup_{R,n,k}m^2\|A_r u\|_\infty\le C.
\]
The proof uses a capped superlinear marked-incidence estimate, not a
pointwise bound for another bad source and not a bounded multiplier
assertion for $c^{-1}$. All grazing thresholds are summed before the
collision limit, while the auxiliary Fourier band remains fixed.

Let $E_h^\varepsilon$ be the original weighted true-exit mass, and
let $O_h^\varepsilon$ be the positive inverse-twist loss on survivors,
as defined in Section~\ref{sec:v79-volume}. For a fixed
$a>0$, restrict their initial source to $|g|\ge a$; use $\mathrm{gr}$
for a true exit with some internal incidence tending to zero.
Theorem~\ref{thm:v80-noncritical-optical} gives the full estimate
\[
 \limsup_m\sup_{R,n,k}m^2
       \|\bar O_r^{\varepsilon,a}
           +\bar E_r^{\varepsilon,\mathrm{gr},a}\|_\infty
                  \le Cr/a.
\]
This includes every positive internal incidence, without an incidence
corridor. The zero optical density at a grazing exit does not erase
its positive starting mass: that mass is paid by the negative optical
variation. The corresponding nongrazing boundary inequality is
Theorem~\ref{thm:v80-boundary-reduction}; it retains the outgoing
physical flux and explicitly pays the missing pullback by volume drift.

Corollary~\ref{cor:v80-endpoint-budget} inserts these estimates into
the original pointwise objective. The initial critical neighborhood
and the nongrazing boundary flux still require their own height
estimates. An integrable singular insertion at every fixed width is
not a uniform pointwise bound for a critical Abel convolution, nor
is it a theorem about a surface-current insertion.

\subsection{Contact estimates and the literal guard}
The constant in each entry below is uniform in the collision count.
Here $\ell_q$ is the length of a regular endpoint curve and the displayed
powers refer to the contact and flight depths defined above.
\begin{center}
\begin{tabular}{lll}
Quantity & Change bounded by & Original guard factor\\[2pt]
\hline
Incidence $c_j$ & $Cq_0^{d_j}\ell_q$ & Yes\\
Third-disk clearance $\Delta_j$ & $Cq_0^{d'_j}\ell_q$ & Yes\\
Section-edge distance & $Cq_0^{d_j}\ell_q$ & No\\
Flight length $\ell_j$ & $Cq_0^{d'_j}\ell_q$ & No
\end{tabular}
\end{center}
Section decisions remain exact event tests. They are not new smooth
guard factors. The canonical field crosses coordinate cuts by change
of chart; it stops at the prescribed physical and section-decision
boundaries. Exhaustion boundaries enter only the limiting flux identity
in Proposition~\ref{prop:v79-source-exit-flux}.

\input{core/174_canonical_endpoint_momentum_flow}
\input{core/177_all_contact_motion_and_guard}
\input{core/180_uniform_cubic_reduced_action}
\input{core/181_separated_physical_volume_budget}
\input{core/182_optical_sensitivity_and_exit_flux}

\input{core/183_homogeneous_incidence_localization}
\input{core/184_singular_optical_transport}
\input{core/185_common_coarea_and_remaining_boundary}

\section{The remaining scalar estimates and their consequences}
The organizing assertion remains \eqref{eq:v76-target}, with the
complete arithmetic transition, original normalization, exact labels
and original smooth guard. The new singular insertion theorem closes
the incidence-scale summation needed for optical drift away from the
initial critical set. The boundary inequality estimates the same
pullback cost for nongrazing exits, but does not bound their outgoing
flux height. The two residual sources are displayed in
\eqref{eq:v80-endpoint-budget}; neither is suppressed by changing the
event or by a signed occupation-current cancellation.

All widths, including the initial critical threshold, are fixed
before the collision limit. After that limit the auxiliary spectral
band is enlarged, and the outer reconstruction parameter is taken
last. The full scalar theorem is not asserted until the two remaining
height estimates have been proved in this order. The new capped
remainder, common-kernel argument and singular-volume estimate do not
license exchanging these limits.

After scalar closure, Supplement
Theorem~\ref{supp-thm:v73-scalar-path-equality} and the positive path
remainder give the same-roof path numerator. Normalization is made
only on positive arithmetic reference classes. Collision-to-return
clock transfer and forward likelihood retain their separate inherited
hypotheses; arithmetic zero classes are never divided by. The complete
preceding article and every inherited proof remain in the supplement.


\endgroup
\part{Uniform cubic geometry and separated optical transport}
\begingroup
\let\ArchiveLabel\label
\let\ArchiveRef\ref
\renewcommand{\label}[1]{\ArchiveLabel{v79archive:#1}}
\renewcommand{\ref}[1]{\ArchiveRef{v79archive:#1}}
\makeatletter
\let\ltx@label\label
\makeatother
\paragraph{Retained revision-79 abstract.}

We study the pointwise local limit problem for the full displacement,
collision-count, occupation and roof record of actual returns in a
triangular Lorentz gas. We prove a collision-uniform third-derivative
bound for the reduced roof at every regular physical word, without a
lower internal-incidence cutoff. A stationary-action identity eliminates
all second derivatives of the internal contact map. The physical
Jacobian of the canonical endpoint-momentum flow then splits into an
angular factor and an inverse-twist ratio. The complete angular-volume
error is bounded by a constant times the square root of the previously
fixed roof width at the required collision normalization. Positive
optical sensitivities give the same type of estimate on incidence
corridors, with the remaining thin part charged to the original source.
We also express first-exit mass through its initial-weight pullback to
outgoing boundaries. The resulting estimate separates true exits from
optical loss while retaining every exact label and arithmetic class.
The uniform roof-height bounds for these two remaining terms are stated
separately; they are not inferred from the controlled angular factor.



\section{The unchanged physical problem}
Put $R_-=9/20$ and $R_+=47/100$, and let $R\in I=[R_-,R_+]$. Circular scatterers have centers in
$\Lambda=\Z(1,0)+\Z(1/2,\sqrt3/2)$. Write $T_R$ for the regular
collision map and $(\alpha,p)\in(\R/2\pi\Z)\times(-1,1)$ for its
coordinates; $p$ is the outgoing tangential velocity and
$\dd\nu=\dd\alpha\,\dd p/(4\pi)$. Physical flights satisfy
$\tau_-\le\tau_R\le\tau_+$, where $\tau_-=3/50$ and $\tau_+$
is the fixed finite-horizon bound.

The section is specified here to fix every event in the argument.
Put $h_0=1/200$. Let $Y_R$ be the union of the four open coordinate
squares of half-side $h_0$ centered at
\[
 (0,0),\quad(\pi/6,-1/2),\quad
 (-\theta_R,\sin\theta_R),\quad
 (\pi/3-\theta_R,\sin\theta_R).
\]
Here $\theta_R=2\arctan q_R$, and $q_R$ is the unique root in
$(49/100,53/100)$ of
\[
 8Rq(1+q^2)-4q(1-q^2)+\sqrt3(1-6q^2+q^4)=0.
\]
Set
\[
 U=\{|\alpha-\pi/6|<3/50,\ |p|<3/20\},\qquad
 Y_R^*=Y_R\cup U,\qquad c_*=\frac{91}{10000\pi}.
\]
The rectangles are disjoint and $\nu(Y_R^*)=c_*$. We always use
$\nu_R^*=c_*^{-1}\nu|_{Y_R^*}$, not a renormalization of a smaller
source.

At $m$ flights, let $L=L_{m,R}$ be their total length,
$K^{\rm c}_{m,R}$ their lattice displacement, and
$A_{m,R}=\sum_{j=0}^{m-1}\one_{Y_R^*}\circ T_R^j$. The event is
\begin{equation}\label{eq:v76-event}
 \Pi_{n,k,m,R}=Y_R^*\cap\{K^{\rm c}_{m,R}=k,\ A_{m,R}=n\}
                         \cap T_R^{-m}Y_R^*.
\end{equation}
The initial visit is counted; the terminal visit is tested separately.
On this event the actual $n$th return has collision count $m$ and
roof $L$. Denote its density by $p_{n,R}(k,m,t)$.

Write $c_j=\sqrt{1-p_j^2}$ and let $\Delta_j$ be the least clearance
of flight $j$ from any third disk, excluding its two incident disks.
Keep the original smooth function $0\le\Theta\le1$, zero on
$(-\infty,1]$ and one on $[2,\infty)$. With $q_0=47/53$,
$d_j=\min(j,m-j)$ and $d'_j=\min(j,m-1-j)$, set
\begin{align}
 H^\varepsilon&=
 \prod_{j=0}^{m}\Theta\!\left(\frac{c_j}{\varepsilon q_0^{d_j/4}}\right)
 \prod_{j=0}^{m-1}\Theta\!\left(\frac{\Delta_j}{\varepsilon q_0^{d'_j/4}}\right),
 \qquad W^\varepsilon=1-H^\varepsilon,\label{eq:v76-guard}\\
 b^\varepsilon(t)\dd t
   &=L_\#(W^\varepsilon\one_{\Pi_{n,k,m,R}}\nu_R^*).
 \label{eq:v76-source}
\end{align}
In particular first incidence and first clearance outside any selected
critical neighborhoods are still present. No section-distance factor
has been inserted into this guard.

\subsection{The arithmetic reference and the target}
For $y=(k_1,k_2,t,n)$ put
\[
 \mu_R=(0,0,\bar\tau_R,c_*),\quad
 Z=(y-m\mu_R)/\sqrt m,\quad
 \mathcal C_{m,R,n,k}(M)=\{t:|Z|\le M\}.
\]
We retain the finite spectral-peak reference, including peaks close to
but not on the unit circle. To specify its meaning, the collision
record observable is $(\kappa_1,\kappa_2,\tau,\one_{Y_R^*})$.
Let $\lambda_j(R,\vartheta)$ and $\mathsf P_j(R,\vartheta)$ be its
continued simple twisted eigenvalue and projection in each of the
fixed peak charts, and let $\chi_j$ be their frequency partition.
At the real maximum $a_j$ of $\log|\lambda_j|$ put
\[
 z_j=\lambda_j(R,a_j),\quad
 \mu_j=i^{-1}\nabla\log\lambda_j(R,a_j),\quad
 H_j=-\nabla^2\log\lambda_j(R,a_j),
\]
with $\operatorname{Re}H_j\ge aI_4$ for a common $a>0$. If
$\eta_R=\one_{Y_R^*}$ and $\ell$ denotes mass evaluation, put
$B_j=\chi_j(R,a_j)\ell(M_{\eta_R}\mathsf P_j(R,a_j)(\eta_R\nu))$.
Define
\begin{align}
 G_H(z)&=(2\pi)^{-4}\int_{\R^4}e^{-iv\cdot z-v^{\mathsf T}Hv/2}\dd v,
 \notag\\
 \mathcal L_{m,R}(y)&=\operatorname{Re}\left[
 c_*^{-1}\sum_j e^{-ia_j\cdot y}z_j^m B_j
      G_{H_j}\left((y-m\mu_j)/\sqrt m\right)\right].
 \label{eq:v76-reference}
\end{align}
Only the charts active at $R$ are summed. This is the same transition
reference as in the retained spectral theorem. The existence and
uniform estimates for these scalar data belong to that analytic
input, not to the transport argument below. No frequency branch or
arithmetic zero class is deleted.

The governing assertion remains
\begin{equation}\label{eq:v76-target}
 \sup_{R,\,1\le n\le m,\,k\in\Z^2}
 \esssup_{t\in\mathcal C_{m,R,n,k}(M)}
 |m^2p_{n,R}(k,m,t)-\mathcal L_{m,R}(k_1,k_2,t,n)|\longrightarrow0.
\end{equation}
Let $\mathcal E_M$ denote this error's collision limsup. An empty
central target contributes zero. The order of the uniformity is
$\sup_R\esssup_t$, so exceptional density sets may depend on $R$.

\subsection{The analytic inputs in the form used here}
\begin{inputstatement}\label{input:v76-analytic}
The unchanged source \eqref{eq:v76-source} has a nonnegative,
compactly supported $BV$ density at every fixed count. For every
fixed $\varepsilon,\delta>0$, with
$A_\delta f(t)=(2\delta)^{-1}\int_{t-\delta}^{t+\delta}f(s)\dd s$,
\begin{equation}\label{eq:v76-average}
 \limsup_m\sup_{R,n,k}m^2\|A_\delta b^\varepsilon\|_\infty
                                      \le C\varepsilon^{1/16}.
\end{equation}
The constant is independent of the previously fixed width. For large
fixed $B$ and $\varepsilon(B)=A_0B^{-1/12}$,
\begin{equation}\label{eq:v76-remainder}
 m^2p-\mathcal L=m^2b^{\varepsilon(B)}+e_B,\qquad
 \limsup_m\sup_{R,n,k}\esssup_{\mathcal C(M)}|e_B|
                                     \le C_MB^{-1/192}.
\end{equation}
All auxiliary spectral bands in these statements are fixed before
the collision limit and enlarged only afterwards.
\end{inputstatement}
These are the complete-source tame theorem, the fixed-width
complete-source bound, and the positive scalar remainder theorem in
the technical supplement. Their proofs and historical hypotheses are
retained there. The fixed-count source uses threshold integration
\cite{CMInt}, definable monotonicity \cite{Coste}, and the
distributional characterization of $BV$ \cite{AFP}; it does not
assume definability of the original smooth guard. They are restated as inputs to distinguish them from
the new geometric and transport proofs. No long-time bound on a
fixed-count variation constant is part of this statement.



\subsection{Precise inherited inputs}\label{sec:v79-input-map}
The fixed-width estimate is Supplement Theorem~\ArchiveRef{thm:v49-physical-remainder-local},
in particular Supplement equation~\textup{(\ArchiveRef{eq:v49-local-limsup})}.
It uses Supplement Theorem~\ArchiveRef{thm:v49-physical-marked-local}, on the
original exact displacement and occupation coefficient. For an interval
of length $2r$, division by $2r$ gives \eqref{eq:v76-average}. That
proof fixes the auxiliary band, takes the collision limit, and only
then enlarges the auxiliary band. Its leading constant is independent
of the previously fixed interval length. The extension from small
margins to a fixed remaining compact range is stated there explicitly.
At auxiliary parameter one, $\Theta(c_0)=0$ since $c_0\le1$, so $W^1=1$
and $b^1=p$. This supplies the unweighted fixed-width estimate without
asserting a raw pointwise bound. The positive remainder is the inherited
scalar input \eqref{eq:v76-remainder}, stated as Supplement
Theorem~\ArchiveRef{thm:v51-positive-error}; no surface-current estimate is
included in any of these volume-source inputs.

\subsection{The principal geometric estimates}
Let $q=(\alpha_0,\alpha_m)$ be local endpoint angles, $F$ the reduced
roof, $K=D_q^2F$, and $\vartheta=\arg(-p_0,p_m)$. The canonical field
satisfies $XL=1$ and $X\vartheta=0$ and is stopped at the boundary of
the regular exact event. Set $\omega=|\partial_{p_0}\alpha_m|^{-1}$ and
$\chi=\det D_q(F,\vartheta)$. These quantities and their positive
uniform bounds are derived in Section~\ref{sec:v77-canonical}.
Theorem~\ref{thm:v79-uniform-cubic} proves
\[
                         \|D_q^3F\|\le T_*
\]
with one constant for all regular physical words. It is a consequence
of the local stationary equations and the all-contact first derivatives,
not of a uniform implicit-function radius near grazing.

To state the endpoint consequence precisely, let $d\mu_t$ be the
unweighted coarea fiber of the original normalized source, $S_h(t)$
the survivors, and $J_h$ the area Jacobian relative to that source.
The two remaining quantities at a fixed count are
\[
 E_h^\varepsilon(t)=\int_{\{L=t\}\setminus S_h(t)}W^\varepsilon\,d\mu_t,
 \qquad
 O_h^\varepsilon(t)=\int_{S_h(t)}W^\varepsilon
                 \left[1-\frac{\omega(\Phi_hx)}{\omega(x)}\right]_+d\mu_t.
\]
The physical Jacobian is
$J_h=(\omega\circ\Phi_h/\omega)(\chi/\chi\circ\Phi_h)$.
Theorem~\ref{thm:v79-volume-factorization} controls, on the complete
source, the difference between the physical volume loss and $O_h$.
Averaging at a width fixed before the collision limit gives an
$O(\sqrt r)$ bound after the $m^2$ normalization. If $\mathfrak X$ and
$\mathfrak O$ denote the separate ordered heights of the averages of
$E_h$ and $O_h$, respectively, Theorem~\ref{thm:v79-separated-endpoint}
gives
\[
 \mathcal E_M\le C_MB^{-1/192}
       +\mathfrak X_{\varepsilon(B),M}(r)
       +\mathfrak O_{\varepsilon(B),M}(r)
       +C_\Theta(1+\varepsilon(B)^{-1})\sqrt r.
\]
The order is: fix $B$, $\varepsilon(B)$, $r$ and the auxiliary band;
then take the collision limit; enlarge the auxiliary band afterwards;
and take $B\to\infty$ last. No positive exit mass has been removed by
an internal occupation-current cancellation. Such a cancellation is a
signed statement, with nonzero occupation-frequency phases still present.

\subsection{Contact estimates and the literal guard}
The constant in each entry below is uniform in the collision count.
Here $\ell_q$ is the length of a regular endpoint curve and the displayed
powers refer to the contact and flight depths defined above.
\begin{center}
\begin{tabular}{lll}
Quantity & Change bounded by & Original guard factor\\[2pt]
\hline
Incidence $c_j$ & $Cq_0^{d_j}\ell_q$ & Yes\\
Third-disk clearance $\Delta_j$ & $Cq_0^{d'_j}\ell_q$ & Yes\\
Section-edge distance & $Cq_0^{d_j}\ell_q$ & No\\
Flight length $\ell_j$ & $Cq_0^{d'_j}\ell_q$ & No
\end{tabular}
\end{center}
Section decisions remain exact event tests. They are not new smooth
guard factors. The canonical field crosses coordinate cuts by change
of chart; it stops at the prescribed physical and section-decision
boundaries. Exhaustion boundaries enter only the limiting flux identity
in Proposition~\ref{prop:v79-source-exit-flux}.

\input{core/174_canonical_endpoint_momentum_flow}
\input{core/177_all_contact_motion_and_guard}
\input{core/180_uniform_cubic_reduced_action}
\input{core/181_separated_physical_volume_budget}
\input{core/182_optical_sensitivity_and_exit_flux}

\section{Completion of the original pointwise theorem}
The original assertion is \eqref{eq:v76-target}, with its full
arithmetic transition and all exact labels. The new uniform cubic
estimate and the complete angular-volume estimate require no new
statistical hypothesis beyond the stated fixed-width input.
The incidence-corridor theorem also estimates a complete nonnegative
part of the optical loss; its thin complement remains charged to the
original source rather than omitted.

The two quantitative bounds in \eqref{eq:v79-two-hypotheses}, or the
separated qualitative criterion \eqref{eq:v79-separated-criterion},
are still required to conclude the organizing theorem. The first-exit
formula shows why ordinary surface flux without the inverse Jacobian
is not such a bound. The optical sensitivity formula shows which
inverse-incidence factors survive after the uniform cubic estimate.
No theorem for bounded volume multipliers is promoted to a theorem for
these surface or singular insertions. Thus the pointwise law and its
unconditional consequences are not asserted on the basis of those formulas.

After scalar closure, Supplement Theorem~\ArchiveRef{thm:v73-scalar-path-equality}
and the retained positive path remainder supply the
same-roof path numerator; normalization is made only on positive
arithmetic reference classes. Collision-to-return clock transfer and
forward likelihood retain their stated inherited hypotheses. Zero
classes are never divided by. The complete earlier arguments, including
the angular recoupling and occupation-current proofs, are retained in
the technical supplement. They do not replace the two remaining estimates.


\endgroup
\part{Complete guard distortion and exact occupation currents}
\begingroup
\let\ArchiveLabel\label
\let\ArchiveRef\ref
\renewcommand{\label}[1]{\ArchiveLabel{v78archive:#1}}
\renewcommand{\ref}[1]{\ArchiveRef{v78archive:#1}}
\makeatletter
\let\ltx@label\label
\makeatother
\paragraph{Retained revision-78 abstract.}

We study the pointwise local limit problem for the complete displacement,
collision-count, occupation and roof record of actual returns in a
triangular Lorentz gas. The physical endpoint-momentum field gives
collision-uniform control of every contact, with an exponential loss
of influence toward the middle of the word. This holds on all regular
surviving trajectories, without a lower incidence cutoff. It yields a
uniform square-root modulus for the original full guard. By transporting
its error to the target fiber and averaging before the collision limit,
we estimate the complete guard-distortion contribution at the required
normalization using only fixed-width concentration. The remaining
geometric contribution retains true exits and physical coarea-volume
loss. We also prove an exact discrete-divergence identity for internal
section-decision currents, with a peak-aware Fourier estimate that keeps
nontrivial arithmetic phases. These results advance the original
pointwise theorem without replacing its exact labels or its source;
the complete true-exit and geometric-distortion estimate is the
quantitative step still to be established.



\section{The unchanged physical problem}
Put $R_-=9/20$ and $R_+=47/100$, and let $R\in I=[R_-,R_+]$. Circular scatterers have centers in
$\Lambda=\Z(1,0)+\Z(1/2,\sqrt3/2)$. Write $T_R$ for the regular
collision map and $(\alpha,p)\in(\R/2\pi\Z)\times(-1,1)$ for its
coordinates; $p$ is the outgoing tangential velocity and
$\dd\nu=\dd\alpha\,\dd p/(4\pi)$. Physical flights satisfy
$\tau_-\le\tau_R\le\tau_+$, where $\tau_-=3/50$ and $\tau_+$
is the fixed finite-horizon bound.

The section is specified here to fix every event in the argument.
Put $h_0=1/200$. Let $Y_R$ be the union of the four open coordinate
squares of half-side $h_0$ centered at
\[
 (0,0),\quad(\pi/6,-1/2),\quad
 (-\theta_R,\sin\theta_R),\quad
 (\pi/3-\theta_R,\sin\theta_R).
\]
Here $\theta_R=2\arctan q_R$, and $q_R$ is the unique root in
$(49/100,53/100)$ of
\[
 8Rq(1+q^2)-4q(1-q^2)+\sqrt3(1-6q^2+q^4)=0.
\]
Set
\[
 U=\{|\alpha-\pi/6|<3/50,\ |p|<3/20\},\qquad
 Y_R^*=Y_R\cup U,\qquad c_*=\frac{91}{10000\pi}.
\]
The rectangles are disjoint and $\nu(Y_R^*)=c_*$. We always use
$\nu_R^*=c_*^{-1}\nu|_{Y_R^*}$, not a renormalization of a smaller
source.

At $m$ flights, let $L=L_{m,R}$ be their total length,
$K^{\rm c}_{m,R}$ their lattice displacement, and
$A_{m,R}=\sum_{j=0}^{m-1}\one_{Y_R^*}\circ T_R^j$. The event is
\begin{equation}\label{eq:v76-event}
 \Pi_{n,k,m,R}=Y_R^*\cap\{K^{\rm c}_{m,R}=k,\ A_{m,R}=n\}
                         \cap T_R^{-m}Y_R^*.
\end{equation}
The initial visit is counted; the terminal visit is tested separately.
On this event the actual $n$th return has collision count $m$ and
roof $L$. Denote its density by $p_{n,R}(k,m,t)$.

Write $c_j=\sqrt{1-p_j^2}$ and let $\Delta_j$ be the least clearance
of flight $j$ from any third disk, excluding its two incident disks.
Keep the original smooth function $0\le\Theta\le1$, zero on
$(-\infty,1]$ and one on $[2,\infty)$. With $q_0=47/53$,
$d_j=\min(j,m-j)$ and $d'_j=\min(j,m-1-j)$, set
\begin{align}
 H^\varepsilon&=
 \prod_{j=0}^{m}\Theta\!\left(\frac{c_j}{\varepsilon q_0^{d_j/4}}\right)
 \prod_{j=0}^{m-1}\Theta\!\left(\frac{\Delta_j}{\varepsilon q_0^{d'_j/4}}\right),
 \qquad W^\varepsilon=1-H^\varepsilon,\label{eq:v76-guard}\\
 b^\varepsilon(t)\dd t
   &=L_\#(W^\varepsilon\one_{\Pi_{n,k,m,R}}\nu_R^*).
 \label{eq:v76-source}
\end{align}
In particular first incidence and first clearance outside any selected
critical neighborhoods are still present. No section-distance factor
has been inserted into this guard.

\subsection{The arithmetic reference and the target}
For $y=(k_1,k_2,t,n)$ put
\[
 \mu_R=(0,0,\bar\tau_R,c_*),\quad
 Z=(y-m\mu_R)/\sqrt m,\quad
 \mathcal C_{m,R,n,k}(M)=\{t:|Z|\le M\}.
\]
We retain the finite spectral-peak reference, including peaks close to
but not on the unit circle. To specify its meaning, the collision
record observable is $(\kappa_1,\kappa_2,\tau,\one_{Y_R^*})$.
Let $\lambda_j(R,\vartheta)$ and $\mathsf P_j(R,\vartheta)$ be its
continued simple twisted eigenvalue and projection in each of the
fixed peak charts, and let $\chi_j$ be their frequency partition.
At the real maximum $a_j$ of $\log|\lambda_j|$ put
\[
 z_j=\lambda_j(R,a_j),\quad
 \mu_j=i^{-1}\nabla\log\lambda_j(R,a_j),\quad
 H_j=-\nabla^2\log\lambda_j(R,a_j),
\]
with $\operatorname{Re}H_j\ge aI_4$ for a common $a>0$. If
$\eta_R=\one_{Y_R^*}$ and $\ell$ denotes mass evaluation, put
$B_j=\chi_j(R,a_j)\ell(M_{\eta_R}\mathsf P_j(R,a_j)(\eta_R\nu))$.
Define
\begin{align}
 G_H(z)&=(2\pi)^{-4}\int_{\R^4}e^{-iv\cdot z-v^{\mathsf T}Hv/2}\dd v,
 \notag\\
 \mathcal L_{m,R}(y)&=\operatorname{Re}\left[
 c_*^{-1}\sum_j e^{-ia_j\cdot y}z_j^m B_j
      G_{H_j}\left((y-m\mu_j)/\sqrt m\right)\right].
 \label{eq:v76-reference}
\end{align}
Only the charts active at $R$ are summed. This is the same transition
reference as in the retained spectral theorem. The existence and
uniform estimates for these scalar data belong to that analytic
input, not to the transport argument below. No frequency branch or
arithmetic zero class is deleted.

The governing assertion remains
\begin{equation}\label{eq:v76-target}
 \sup_{R,\,1\le n\le m,\,k\in\Z^2}
 \esssup_{t\in\mathcal C_{m,R,n,k}(M)}
 |m^2p_{n,R}(k,m,t)-\mathcal L_{m,R}(k_1,k_2,t,n)|\longrightarrow0.
\end{equation}
Let $\mathcal E_M$ denote this error's collision limsup. An empty
central target contributes zero. The order of the uniformity is
$\sup_R\esssup_t$, so exceptional density sets may depend on $R$.

\subsection{The analytic inputs in the form used here}
\begin{inputstatement}\label{input:v76-analytic}
The unchanged source \eqref{eq:v76-source} has a nonnegative,
compactly supported $BV$ density at every fixed count. For every
fixed $\varepsilon,\delta>0$, with
$A_\delta f(t)=(2\delta)^{-1}\int_{t-\delta}^{t+\delta}f(s)\dd s$,
\begin{equation}\label{eq:v76-average}
 \limsup_m\sup_{R,n,k}m^2\|A_\delta b^\varepsilon\|_\infty
                                      \le C\varepsilon^{1/16}.
\end{equation}
The constant is independent of the previously fixed width. For large
fixed $B$ and $\varepsilon(B)=A_0B^{-1/12}$,
\begin{equation}\label{eq:v76-remainder}
 m^2p-\mathcal L=m^2b^{\varepsilon(B)}+e_B,\qquad
 \limsup_m\sup_{R,n,k}\esssup_{\mathcal C(M)}|e_B|
                                     \le C_MB^{-1/192}.
\end{equation}
All auxiliary spectral bands in these statements are fixed before
the collision limit and enlarged only afterwards.
\end{inputstatement}
These are the complete-source tame theorem, the fixed-width
complete-source bound, and the positive scalar remainder theorem in
the technical supplement. Their proofs and historical hypotheses are
retained there. The fixed-count source uses threshold integration
\cite{CMInt}, definable monotonicity \cite{Coste}, and the
distributional characterization of $BV$ \cite{AFP}; it does not
assume definability of the original smooth guard. They are restated as inputs to distinguish them from
the new geometric and transport proofs. No long-time bound on a
fixed-count variation constant is part of this statement.


\subsection{The physical construction and the new estimate}
The endpoint-momentum field and complete angular recoupling are
retained below with their proofs. The new argument then estimates
one full dynamical contribution, rather than replacing it by a
per-word bound. All-contact optical response gives a uniform modulus
for the unchanged guard, even with arbitrarily small internal
incidences. Its transported error is measured against a translated
unweighted target; averaging that target permits the fixed-width
analytic input to be used without a pointwise density assumption.

For the remaining physical geometry, internal section interfaces
admit an exact signed difference in the occupation index. That
identity is stated before taking absolute values, and its Fourier
consequence is formulated at every peak rather than only at the
origin. True singular and outer section boundaries remain explicit.

\input{core/174_canonical_endpoint_momentum_flow}
\input{core/175_complete_angular_recoupling}
\input{core/176_angular_current_and_ordered_endpoint}

\input{core/177_all_contact_motion_and_guard}
\input{core/178_averaged_guard_distortion}
\input{core/179_exact_occupation_boundary_current}

\section{The complete geometric estimate and the original theorem}
The pointwise assertion remains \eqref{eq:v76-target}, with the full
arithmetic transition, the same normalization and all zero classes.
Theorem~\ref{thm:v78-separation} proves the complete guard-distortion
bound \eqref{eq:v78-guard-closure}; the constant is uniform in collision
count and no selected physical word has replaced the source.
Proposition~\ref{prop:v78-exit-localization} locates every true forward
exit, but does not bound concentration on the starting roof.
Theorem~\ref{thm:v78-decision-current} computes the signed internal
section contribution without suppressing its exact-label coefficients.

To conclude the organizing theorem, one must still estimate the
complete geometric term in \eqref{eq:v78-ordered-geometry}, or the
complete angular defect in \eqref{eq:v77-ordered-defect}. The new
guard estimate does not control the optical Jacobian loss. An
integrated boundary-layer estimate is not a pointwise source estimate,
and a spectral estimate for volume multipliers is not automatically
an estimate for surface-current insertions. In particular
\eqref{eq:v78-peak-difference} retains nonzero arithmetic leading
terms. The proof ledger records these distinctions explicitly.

The organizing pointwise theorem is therefore not marked proved in
this revision. Its path and likelihood consequences retain the same
scalar-first logical order. The complete inherited mathematical
source and the preceding journal article are retained in the technical
supplement and provenance, rather than deleted or recast as a different
problem.


\endgroup
\part{Endpoint momenta and complete angular recoupling}
\begingroup
\let\ArchiveLabel\label
\let\ArchiveRef\ref
\renewcommand{\label}[1]{\ArchiveLabel{v77archive:#1}}
\renewcommand{\ref}[1]{\ArchiveRef{v77archive:#1}}
\makeatletter
\let\ltx@label\label
\makeatother
\paragraph{Retained revision-77 abstract.}

We retain the pointwise local limit problem for the full displacement,
collision count, occupation, and roof record of actual returns in a
triangular Lorentz gas. The exact initial and terminal section tests
give collision-uniform positive bounds for the Hessian in physical
endpoint angles. They determine an explicit roof flow preserving the
angle of the endpoint momenta. Its forward motion is uniformly
square-root controlled and never reaches an interior critical point;
no Morse collars or cap-dependent patching are needed. We disintegrate
the complete original source jointly over roof and momentum angle,
and construct a maximal angle-preserving subcoupling after summing
all actual words with the same exact label. This recoupling can recover
mass lost by an individual stopped orbit. A finite signed angular
current retains true boundary flux and the positive critical residues.
These results sharpen the complete-source persistence estimate for
the original pointwise law. The collision-normalized concentration
bound for the complete angular defect remains the quantitative
estimate required to close that law.



\section{The unchanged physical problem}
Put $R_-=9/20$ and $R_+=47/100$, and let $R\in I=[R_-,R_+]$. Circular scatterers have centers in
$\Lambda=\Z(1,0)+\Z(1/2,\sqrt3/2)$. Write $T_R$ for the regular
collision map and $(\alpha,p)\in(\R/2\pi\Z)\times(-1,1)$ for its
coordinates; $p$ is the outgoing tangential velocity and
$\dd\nu=\dd\alpha\,\dd p/(4\pi)$. Physical flights satisfy
$\tau_-\le\tau_R\le\tau_+$, where $\tau_-=3/50$ and $\tau_+$
is the fixed finite-horizon bound.

The section is specified here to fix every event in the argument.
Put $h_0=1/200$. Let $Y_R$ be the union of the four open coordinate
squares of half-side $h_0$ centered at
\[
 (0,0),\quad(\pi/6,-1/2),\quad
 (-\theta_R,\sin\theta_R),\quad
 (\pi/3-\theta_R,\sin\theta_R).
\]
Here $\theta_R=2\arctan q_R$, and $q_R$ is the unique root in
$(49/100,53/100)$ of
\[
 8Rq(1+q^2)-4q(1-q^2)+\sqrt3(1-6q^2+q^4)=0.
\]
Set
\[
 U=\{|\alpha-\pi/6|<3/50,\ |p|<3/20\},\qquad
 Y_R^*=Y_R\cup U,\qquad c_*=\frac{91}{10000\pi}.
\]
The rectangles are disjoint and $\nu(Y_R^*)=c_*$. We always use
$\nu_R^*=c_*^{-1}\nu|_{Y_R^*}$, not a renormalization of a smaller
source.

At $m$ flights, let $L=L_{m,R}$ be their total length,
$K^{\rm c}_{m,R}$ their lattice displacement, and
$A_{m,R}=\sum_{j=0}^{m-1}\one_{Y_R^*}\circ T_R^j$. The event is
\begin{equation}\label{eq:v76-event}
 \Pi_{n,k,m,R}=Y_R^*\cap\{K^{\rm c}_{m,R}=k,\ A_{m,R}=n\}
                         \cap T_R^{-m}Y_R^*.
\end{equation}
The initial visit is counted; the terminal visit is tested separately.
On this event the actual $n$th return has collision count $m$ and
roof $L$. Denote its density by $p_{n,R}(k,m,t)$.

Write $c_j=\sqrt{1-p_j^2}$ and let $\Delta_j$ be the least clearance
of flight $j$ from any third disk, excluding its two incident disks.
Keep the original smooth function $0\le\Theta\le1$, zero on
$(-\infty,1]$ and one on $[2,\infty)$. With $q_0=47/53$,
$d_j=\min(j,m-j)$ and $d'_j=\min(j,m-1-j)$, set
\begin{align}
 H^\varepsilon&=
 \prod_{j=0}^{m}\Theta\!\left(\frac{c_j}{\varepsilon q_0^{d_j/4}}\right)
 \prod_{j=0}^{m-1}\Theta\!\left(\frac{\Delta_j}{\varepsilon q_0^{d'_j/4}}\right),
 \qquad W^\varepsilon=1-H^\varepsilon,\label{eq:v76-guard}\\
 b^\varepsilon(t)\dd t
   &=L_\#(W^\varepsilon\one_{\Pi_{n,k,m,R}}\nu_R^*).
 \label{eq:v76-source}
\end{align}
In particular first incidence and first clearance outside any selected
critical neighborhoods are still present. No section-distance factor
has been inserted into this guard.

\subsection{The arithmetic reference and the target}
For $y=(k_1,k_2,t,n)$ put
\[
 \mu_R=(0,0,\bar\tau_R,c_*),\quad
 Z=(y-m\mu_R)/\sqrt m,\quad
 \mathcal C_{m,R,n,k}(M)=\{t:|Z|\le M\}.
\]
We retain the finite spectral-peak reference, including peaks close to
but not on the unit circle. To specify its meaning, the collision
record observable is $(\kappa_1,\kappa_2,\tau,\one_{Y_R^*})$.
Let $\lambda_j(R,\vartheta)$ and $\mathsf P_j(R,\vartheta)$ be its
continued simple twisted eigenvalue and projection in each of the
fixed peak charts, and let $\chi_j$ be their frequency partition.
At the real maximum $a_j$ of $\log|\lambda_j|$ put
\[
 z_j=\lambda_j(R,a_j),\quad
 \mu_j=i^{-1}\nabla\log\lambda_j(R,a_j),\quad
 H_j=-\nabla^2\log\lambda_j(R,a_j),
\]
with $\operatorname{Re}H_j\ge aI_4$ for a common $a>0$. If
$\eta_R=\one_{Y_R^*}$ and $\ell$ denotes mass evaluation, put
$B_j=\chi_j(R,a_j)\ell(M_{\eta_R}\mathsf P_j(R,a_j)(\eta_R\nu))$.
Define
\begin{align}
 G_H(z)&=(2\pi)^{-4}\int_{\R^4}e^{-iv\cdot z-v^{\mathsf T}Hv/2}\dd v,
 \notag\\
 \mathcal L_{m,R}(y)&=\operatorname{Re}\left[
 c_*^{-1}\sum_j e^{-ia_j\cdot y}z_j^m B_j
      G_{H_j}\left((y-m\mu_j)/\sqrt m\right)\right].
 \label{eq:v76-reference}
\end{align}
Only the charts active at $R$ are summed. This is the same transition
reference as in the retained spectral theorem. The existence and
uniform estimates for these scalar data belong to that analytic
input, not to the transport argument below. No frequency branch or
arithmetic zero class is deleted.

The governing assertion remains
\begin{equation}\label{eq:v76-target}
 \sup_{R,\,1\le n\le m,\,k\in\Z^2}
 \esssup_{t\in\mathcal C_{m,R,n,k}(M)}
 |m^2p_{n,R}(k,m,t)-\mathcal L_{m,R}(k_1,k_2,t,n)|\longrightarrow0.
\end{equation}
Let $\mathcal E_M$ denote this error's collision limsup. An empty
central target contributes zero. The order of the uniformity is
$\sup_R\esssup_t$, so exceptional density sets may depend on $R$.

\subsection{The analytic inputs in the form used here}
\begin{inputstatement}\label{input:v76-analytic}
The unchanged source \eqref{eq:v76-source} has a nonnegative,
compactly supported $BV$ density at every fixed count. For every
fixed $\varepsilon,\delta>0$, with
$A_\delta f(t)=(2\delta)^{-1}\int_{t-\delta}^{t+\delta}f(s)\dd s$,
\begin{equation}\label{eq:v76-average}
 \limsup_m\sup_{R,n,k}m^2\|A_\delta b^\varepsilon\|_\infty
                                      \le C\varepsilon^{1/16}.
\end{equation}
The constant is independent of the previously fixed width. For large
fixed $B$ and $\varepsilon(B)=A_0B^{-1/12}$,
\begin{equation}\label{eq:v76-remainder}
 m^2p-\mathcal L=m^2b^{\varepsilon(B)}+e_B,\qquad
 \limsup_m\sup_{R,n,k}\esssup_{\mathcal C(M)}|e_B|
                                     \le C_MB^{-1/192}.
\end{equation}
All auxiliary spectral bands in these statements are fixed before
the collision limit and enlarged only afterwards.
\end{inputstatement}
These are the complete-source tame theorem, the fixed-width
complete-source bound, and the positive scalar remainder theorem in
the technical supplement. Their proofs and historical hypotheses are
retained there. The fixed-count source uses threshold integration
\cite{CMInt}, definable monotonicity \cite{Coste}, and the
distributional characterization of $BV$ \cite{AFP}; it does not
assume definability of the original smooth guard. They are restated as inputs to distinguish them from
the new geometric and transport proofs. No long-time bound on a
fixed-count variation constant is part of this statement.


\subsection{The new physical construction}
The retained convex-support and optical arguments provide genuine
information on every regular physical word. They do not on their
own control the complete weighted sum on an exact roof. The present
construction uses a different observable from the auxiliary cap
normal: the angle of the two physical endpoint momenta. The fixed
section makes its coarea Jacobian uniformly positive after the
optical twist is factored out. The resulting field is explicit and
requires no collision-dependent critical neighborhoods.

The new subcoupling is constructed after summing all actual words of
one exact event. It may match mass between different true words,
while preserving their exact displacement, occupation, terminal test,
and endpoint-momentum angle. This is more permissive than continuation
of a single physical orbit and more restrictive than unconstrained
scalar matching. Its finite-current estimate is an unconditional
fixed-count statement. Its collision-normalized concentration remains
a distinct dynamical question.

\input{core/174_canonical_endpoint_momentum_flow}
\input{core/175_complete_angular_recoupling}
\input{core/176_angular_current_and_ordered_endpoint}

\section{The remaining complete-source estimate}
The principal assertion remains \eqref{eq:v76-target}, with its full
arithmetic transition and all zero classes. The new uniform bounds
concern the actual endpoint Hessian and the forward motion of an
explicit physical field. They remove arbitrary field choice and
Morse-radius patching from that construction. The exact angular
sum and its maximal recoupling also avoid charging a separate loss
when another actual word supplies matching mass at the same angle
and label. True boundaries are neither declared null in flux nor
removed from the original guard.

What is still needed is \eqref{eq:v77-remaining}, or a direct estimate
of comparable strength for the original complete source. Neither
$\int A^\varepsilon\le1$ nor the finite constants $N_m,J_m,C_m$
control concentration at one roof after multiplication by $m^2$.
The explicit drift formula also involves third derivatives and source
derivatives for which no collision-uniform summed bound has been
obtained here. Thus the organizing pointwise theorem and its
unconditional path and likelihood consequences are not claimed as
proved. All earlier arguments, including their exact hypotheses,
are retained in the technical supplement; the target is not replaced
by a smoothed law or a different record.


\endgroup
\part{Stopped physical transport and optical persistence}
\paragraph{Retained revision-76 abstract.}

We retain the pointwise local limit problem for the full displacement,
collision count, occupation, and roof record of actual returns in a
triangular Lorentz gas. We replace matching of coordinate-dependent
Gauss traces by a stopped transport of the original physical level
measures. The resulting identity charges only genuine exits and
relative source distortion; changes of endpoint caps make no
contribution. A positive transfer-matrix comparison bounds the inverse
twist of every regular word by an explicit hyperbolic-sine factor,
while retaining simultaneous incidence gains. At an interior critical
roof, a Morse-adapted transport has zero planar Jacobian distortion
and a square-root source defect, without charging the positive jump.
These results give a complete-source persistence inequality and a
self-contained implication to the original arithmetic pointwise law.
The collision-normalized sum of the true exit and distortion defects
is the additional estimate required for that implication; it is not
inferred from the single-word comparison or from fixed-count collars.



\section{The unchanged physical problem}
Put $R_-=9/20$ and $R_+=47/100$, and let $R\in I=[R_-,R_+]$. Circular scatterers have centers in
$\Lambda=\Z(1,0)+\Z(1/2,\sqrt3/2)$. Write $T_R$ for the regular
collision map and $(\alpha,p)\in(\R/2\pi\Z)\times(-1,1)$ for its
coordinates; $p$ is the outgoing tangential velocity and
$\dd\nu=\dd\alpha\,\dd p/(4\pi)$. Physical flights satisfy
$\tau_-\le\tau_R\le\tau_+$, where $\tau_-=3/50$ and $\tau_+$
is the fixed finite-horizon bound.

The section is specified here to fix every event in the argument.
Put $h_0=1/200$. Let $Y_R$ be the union of the four open coordinate
squares of half-side $h_0$ centered at
\[
 (0,0),\quad(\pi/6,-1/2),\quad
 (-\theta_R,\sin\theta_R),\quad
 (\pi/3-\theta_R,\sin\theta_R).
\]
Here $\theta_R=2\arctan q_R$, and $q_R$ is the unique root in
$(49/100,53/100)$ of
\[
 8Rq(1+q^2)-4q(1-q^2)+\sqrt3(1-6q^2+q^4)=0.
\]
Set
\[
 U=\{|\alpha-\pi/6|<3/50,\ |p|<3/20\},\qquad
 Y_R^*=Y_R\cup U,\qquad c_*=\frac{91}{10000\pi}.
\]
The rectangles are disjoint and $\nu(Y_R^*)=c_*$. We always use
$\nu_R^*=c_*^{-1}\nu|_{Y_R^*}$, not a renormalization of a smaller
source.

At $m$ flights, let $L=L_{m,R}$ be their total length,
$K^{\rm c}_{m,R}$ their lattice displacement, and
$A_{m,R}=\sum_{j=0}^{m-1}\one_{Y_R^*}\circ T_R^j$. The event is
\begin{equation}\label{eq:v76-event}
 \Pi_{n,k,m,R}=Y_R^*\cap\{K^{\rm c}_{m,R}=k,\ A_{m,R}=n\}
                         \cap T_R^{-m}Y_R^*.
\end{equation}
The initial visit is counted; the terminal visit is tested separately.
On this event the actual $n$th return has collision count $m$ and
roof $L$. Denote its density by $p_{n,R}(k,m,t)$.

Write $c_j=\sqrt{1-p_j^2}$ and let $\Delta_j$ be the least clearance
of flight $j$ from any third disk, excluding its two incident disks.
Keep the original smooth function $0\le\Theta\le1$, zero on
$(-\infty,1]$ and one on $[2,\infty)$. With $q_0=47/53$,
$d_j=\min(j,m-j)$ and $d'_j=\min(j,m-1-j)$, set
\begin{align}
 H^\varepsilon&=
 \prod_{j=0}^{m}\Theta\!\left(\frac{c_j}{\varepsilon q_0^{d_j/4}}\right)
 \prod_{j=0}^{m-1}\Theta\!\left(\frac{\Delta_j}{\varepsilon q_0^{d'_j/4}}\right),
 \qquad W^\varepsilon=1-H^\varepsilon,\label{eq:v76-guard}\\
 b^\varepsilon(t)\dd t
   &=L_\#(W^\varepsilon\one_{\Pi_{n,k,m,R}}\nu_R^*).
 \label{eq:v76-source}
\end{align}
In particular first incidence and first clearance outside any selected
critical neighborhoods are still present. No section-distance factor
has been inserted into this guard.

\subsection{The arithmetic reference and the target}
For $y=(k_1,k_2,t,n)$ put
\[
 \mu_R=(0,0,\bar\tau_R,c_*),\quad
 Z=(y-m\mu_R)/\sqrt m,\quad
 \mathcal C_{m,R,n,k}(M)=\{t:|Z|\le M\}.
\]
We retain the finite spectral-peak reference, including peaks close to
but not on the unit circle. To specify its meaning, the collision
record observable is $(\kappa_1,\kappa_2,\tau,\one_{Y_R^*})$.
Let $\lambda_j(R,\vartheta)$ and $\mathsf P_j(R,\vartheta)$ be its
continued simple twisted eigenvalue and projection in each of the
fixed peak charts, and let $\chi_j$ be their frequency partition.
At the real maximum $a_j$ of $\log|\lambda_j|$ put
\[
 z_j=\lambda_j(R,a_j),\quad
 \mu_j=i^{-1}\nabla\log\lambda_j(R,a_j),\quad
 H_j=-\nabla^2\log\lambda_j(R,a_j),
\]
with $\operatorname{Re}H_j\ge aI_4$ for a common $a>0$. If
$\eta_R=\one_{Y_R^*}$ and $\ell$ denotes mass evaluation, put
$B_j=\chi_j(R,a_j)\ell(M_{\eta_R}\mathsf P_j(R,a_j)(\eta_R\nu))$.
Define
\begin{align}
 G_H(z)&=(2\pi)^{-4}\int_{\R^4}e^{-iv\cdot z-v^{\mathsf T}Hv/2}\dd v,
 \notag\\
 \mathcal L_{m,R}(y)&=\operatorname{Re}\left[
 c_*^{-1}\sum_j e^{-ia_j\cdot y}z_j^m B_j
      G_{H_j}\left((y-m\mu_j)/\sqrt m\right)\right].
 \label{eq:v76-reference}
\end{align}
Only the charts active at $R$ are summed. This is the same transition
reference as in the retained spectral theorem. The existence and
uniform estimates for these scalar data belong to that analytic
input, not to the transport argument below. No frequency branch or
arithmetic zero class is deleted.

The governing assertion remains
\begin{equation}\label{eq:v76-target}
 \sup_{R,\,1\le n\le m,\,k\in\Z^2}
 \esssup_{t\in\mathcal C_{m,R,n,k}(M)}
 |m^2p_{n,R}(k,m,t)-\mathcal L_{m,R}(k_1,k_2,t,n)|\longrightarrow0.
\end{equation}
Let $\mathcal E_M$ denote this error's collision limsup. An empty
central target contributes zero. The order of the uniformity is
$\sup_R\esssup_t$, so exceptional density sets may depend on $R$.

\subsection{The analytic inputs in the form used here}
\begin{inputstatement}\label{input:v76-analytic}
The unchanged source \eqref{eq:v76-source} has a nonnegative,
compactly supported $BV$ density at every fixed count. For every
fixed $\varepsilon,\delta>0$, with
$A_\delta f(t)=(2\delta)^{-1}\int_{t-\delta}^{t+\delta}f(s)\dd s$,
\begin{equation}\label{eq:v76-average}
 \limsup_m\sup_{R,n,k}m^2\|A_\delta b^\varepsilon\|_\infty
                                      \le C\varepsilon^{1/16}.
\end{equation}
The constant is independent of the previously fixed width. For large
fixed $B$ and $\varepsilon(B)=A_0B^{-1/12}$,
\begin{equation}\label{eq:v76-remainder}
 m^2p-\mathcal L=m^2b^{\varepsilon(B)}+e_B,\qquad
 \limsup_m\sup_{R,n,k}\esssup_{\mathcal C(M)}|e_B|
                                     \le C_MB^{-1/192}.
\end{equation}
All auxiliary spectral bands in these statements are fixed before
the collision limit and enlarged only afterwards.
\end{inputstatement}
These are the complete-source tame theorem, the fixed-width
complete-source bound, and the positive scalar remainder theorem in
the technical supplement. Their proofs and historical hypotheses are
retained there. The fixed-count source uses threshold integration
\cite{CMInt}, definable monotonicity \cite{Coste}, and the
distributional characterization of $BV$ \cite{AFP}; it does not
assume definability of the original smooth guard. They are restated as inputs to distinguish them from
the new geometric and transport proofs. No long-time bound on a
fixed-count variation constant is part of this statement.

For context, the preceding geometric theorem supplies at most $256$
endpoint-cap rectangles for each physical word, with one strongly
convex supporting action per rectangle. On every regular physical
piece its Hessian lies between $(50/47)I_2$ and $(1360/9)I_2$.
All its level arcs together have coarea integral at most
$2\pi/(50/47)$. Its positive Gauss formula is exact, but comparing
the numerical normal angle at different roofs depends on the chosen
caps. The construction below instead takes place on the original
physical collision manifold. Its dependence on the chosen smooth
roof vector field is explicit; it has no dependence on a cap
partition used merely to calculate its measures.

\input{core/171_stopped_physical_roof_transport}
\input{core/172_optical_weights_and_morse_transport}
\input{core/173_complete_source_persistence}

\section{The mathematical boundary of the present argument}
The optical comparison controls every regular word and every selected
set of incidence defects. The stopped transport is a complete-source
identity, not a protected-source replacement. Its losses distinguish
physical exits from the changes of coordinates that contributed to
the earlier sufficient trace modulus. The local Morse calculation
also avoids assigning the height of a positive critical jump to a
forward transport defect.

Neither statement bounds the complete normalized exit and distortion
budget in Theorem~\ref{thm:v76-closure}. There may be many words,
small true physical domains, and small critical neighborhoods. Their
weighted concentration on an exact-label roof interval must still be
estimated. The fixed-count Morse radii are not used as widths in a
collision limit. Likewise, the inverse-twist comparison is not summed
by replacing the weighted physical word set with its cardinality.
Proving this remaining dynamical estimate, on the original source
and in the stated order of limits, continues to be the route to
\eqref{eq:v76-target}; the target has not been replaced by a different
record or by a smoothed law.


\part{Convex support and the complete physical Gauss trace}
\paragraph{Retained revision-75 abstract.}

We study exact displacement, collision count and flight time for actual
returns in a triangular periodic Lorentz gas. For every regular physical
word we construct, on a fixed finite cover of endpoint projection
coordinates, a strongly convex supporting action which agrees with the
original roof on its entire physical endpoint image. Both local curvature
bounds are independent of the collision count and of every positive
incidence and clearance. The convex comparison controls all physical
level pieces together: their coarea integral has a count-independent
bound and their Gauss map has at most one inverse per normal direction.
This removes the wordwise level-complexity loss from the complete
inverse-incidence height estimate. We obtain an exact positive Gauss
coarea formula retaining every original return label and physical source.
The remaining collision-uniform short-excursion estimate is expressed
on this explicit inverse-twist trace. The geometric estimates are proved
without assuming that weighted dynamical estimate.



\section{The original record and the geometric estimate}
\label{sec:v75-introduction}
Let $R\in I=[9/20,47/100]$, and let
$J_{n,R}=(K_{n,R},N_{n,R},T_{n,R})$ be the actual return record
under the original section probability
$\nu_R^*=c_*^{-1}\nu|_{Y_R^*}$, where $c_*=91/(10000\pi)$.
The circular scatterers have centers in the triangular lattice
$\Lambda=\Z(1,0)+\Z(1/2,\sqrt3/2)$. In collision coordinates
$(\alpha,p)$, $d\nu=d\alpha\,dp/(4\pi)$ and $p$ is the tangential
component of the outgoing unit velocity. The section, including its
moving centers, is specified without alteration in the supplement's
Section~\ref{sec:v74-physical-graph}.

At $m$ flights let $L_{m,R}$ be their total length, $K^{\rm c}_{m,R}$
the accumulated displacement, and
$A_{m,R}=\sum_{j=0}^{m-1}\one_{Y_R^*}\circ T_R^j$ the occupation.
The exact event used here is
\[
 \Pi_{n,k,m,R}=Y_R^*\cap\{K^{\rm c}_{m,R}=k,\ A_{m,R}=n\}
                            \cap T_R^{-m}Y_R^*.
\]
Thus the initial visit is counted, the terminal visit is tested separately,
and the actual $n$th return has collision count $m$ on this event.
The density of its roof under $\nu_R^*$ is $p_{n,R}(k,m,t)$.
Set
\[
 Z_{m,R}(k,t,n)=\frac{(k_1,k_2,t-m\bar\tau_R,n-mc_*)}{\sqrt m},
 \qquad \mathcal C_{m,R,n,k}(M)=\{t:|Z_{m,R}(k,t,n)|\le M\}.
\]
The governing pointwise assertion remains, for every finite $M$,
\begin{equation}\label{eq:v75-original-target}
 \sup_{R,\,1\le n\le m,\,k\in\Z^2}
 \esssup_{t\in\mathcal C_{m,R,n,k}(M)}
 |m^2p_{n,R}(k,m,t)-\mathcal L_{m,R}(k_1,k_2,t,n)|\longrightarrow0.
\end{equation}
The reference $\mathcal L_{m,R}$ is the complete arithmetic transition
reference in the supplement, not a Gaussian obtained by pooling labels.
Every arithmetic zero class is retained. Denote the collision limsup
of the left side by $\mathcal E_M$.

\subsection{The complete positive source and the imported statements}
Write $c_j=\sqrt{1-p_j^2}$ and let $\Delta_j$ be the actual
third-disk clearance of flight $j$. With $q_0=47/53$, set
$d_j=\min(j,m-j)$ and $d'_j=\min(j,m-1-j)$. The original guard is
\[
 H_{m,R}^\varepsilon=
 \prod_{j=0}^{m}\Theta\!\left(\frac{c_j}{\varepsilon q_0^{d_j/4}}\right)
 \prod_{j=0}^{m-1}\Theta\!\left(\frac{\Delta_j}{\varepsilon q_0^{d'_j/4}}\right),
\]
where the unchanged smooth cutoff satisfies $0\le\Theta\le1$,
$\Theta=0$ on $(-\infty,1]$ and $\Theta=1$ on $[2,\infty)$.
No definability of this smooth cutoff is assumed. Define
\begin{equation}\label{eq:v75-source}
 b^\varepsilon(t)\,dt
   =(L_{m,R})_\#\bigl((1-H_{m,R}^\varepsilon)
                         \one_{\Pi_{n,k,m,R}}\nu_R^*\bigr).
\end{equation}
Every incidence and clearance factor remains in this definition.
In particular the first-clearance source outside selected critical
disks and the first-incidence source are both included.

\begin{inputstatement}[Analytic interfaces on the unchanged source]
\label{input:v75-analytic}
For each fixed $m,R,\varepsilon,n,k$, the density in
\eqref{eq:v75-source} is nonnegative, compactly supported and $BV$.
For each fixed $\varepsilon,\delta>0$, with
$A_\delta f(t)=(2\delta)^{-1}\int_{t-\delta}^{t+\delta}f(s)\,ds$,
\begin{equation}\label{eq:v75-average-input}
 \limsup_m\sup_{R,n,k}m^2\|A_\delta b^\varepsilon\|_\infty
                                  \le C\varepsilon^{1/16}.
\end{equation}
For sufficiently large fixed $B$ and $\varepsilon(B)=A_0B^{-1/12}$,
\begin{equation}\label{eq:v75-positive-input}
 m^2p-\mathcal L=m^2b^{\varepsilon(B)}+e_B,\qquad
 \limsup_m\sup_{R,n,k}\esssup_{\mathcal C(M)}|e_B|
                                  \le C_MB^{-1/192}.
\end{equation}
The constant in \eqref{eq:v75-average-input} is independent of the
previously fixed width. Its auxiliary spectral band is fixed before
the collision limit and enlarged only afterwards. For nonnegative
compactly supported $BV$ densities,
\[
 [f-2A_\delta f]_+\le S_\delta f,\qquad
 S_\delta f(t)=\frac1\delta\int_0^\delta T_f(t,r)\,dr,
\]
where $T_f(t,r)$ counts the height levels whose essential superlevel
component through $t$ has length less than $r$.
\end{inputstatement}

The exact sources are the supplement's
Theorem~\ref{thm:v74-tame-family},
\eqref{eq:v72-finite-average-height},
\eqref{eq:v51-positive-error}, and
Theorem~\ref{thm:v74-level-bound} with its length-tail identity.
The positivity comparison yields
\[
 \mathcal E_M\le\mathcal S_M(B,\delta)+C_MB^{-1/192},\qquad
 \mathcal S_M(B,\delta)=\limsup_m\sup_{R,n,k}m^2
                           \esssup_{\mathcal C(M)}S_\delta b^{\varepsilon(B)}.
\]
The full two-sided comparison is
Theorem~\ref{thm:v74-endpoint-equivalence}. The new geometric
proofs below do not use any collision-uniform variation or tail bound.

For downstream path conclusions we additionally retain
Theorem~\ref{thm:v52-path-remainder}: the path-numerator error has
a positive remainder with exactly the scalar mass $m^2b^{\varepsilon(B)}$
and a bounded-Lipschitz dual error with the same vanishing outer-band
bound. The ordered path and scalar errors then agree by
Theorem~\ref{thm:v73-scalar-path-equality}. Conditional ratios are
formed only on central reference classes bounded below by a positive
constant; arithmetic zero classes are not normalized. Actual-return
path statements further retain the original clock-transfer hypotheses.
These analytic and path interfaces are source-pinned inputs, not
independent spectral or clock proofs supplied by the present argument.

\subsection{The new physical mechanism}
A positive Hessian on the physical endpoint image is not enough:
that image may have many components separated by first-hit and
exact-label boundaries. We instead construct one convex supporting
action on an entire endpoint rectangle. A fixed 16-direction cover
at each endpoint accommodates all positive incidences, even near
grazing. On the physical pieces its Hessian obeys
\[
            \frac{50}{47}I_2\le D^2F_w\le\frac{1360}{9}I_2.
\]
Theorem~\ref{thm:v75-convex-envelope} proves this statement.
Lemma~\ref{lem:v75-total-turning} then bounds all physical level
pieces of a word together by $2\pi/\lambda$, with at most one
inverse for each outward normal. Theorem~\ref{thm:v75-exact-gauss}
uses that inverse to express the complete exact-label density as a
positive Gauss coarea integral. This removes a collision-dependent
wordwise level-complexity loss, rather than concealing it in a
constant.

The remaining sum of inverse-twist weights over physical words is
not thereby controlled at the $m^{-2}$ scale. Section~\ref{sec:v75-physical-tail}
states precisely the paired modulus of that trace sufficient to
complete \eqref{eq:v75-original-target}, and distinguishes it from
the uniform per-word geometry actually proved here. The short
article and the complete supplement retain the same original
pointwise problem.

\input{core/168_convex_support_endpoint_caps}
\input{core/169_wordwise_gauss_coarea}
\input{core/170_physical_gauss_trace_tail}

\part{The retained complete source and short excursions}
\paragraph{Retained revision-74 abstract.}

We study exact displacement, collision count and flight time for actual
returns in a triangular periodic Lorentz gas. For the complete physical
boundary source, we give an explicit tame integration family, including
the moving section and every exact label. We then match incoming and
outgoing current on the same superlevel component. This yields a
pointwise short-excursion bound sharper than the minimum of the two
Jordan masses, a sharp integrated estimate, and an exact length-tail
identity. The ordered short-excursion budget has the same outer limit
as the original arithmetic pointwise error and, under the common
positive path-remainder theorem, its bounded-Lipschitz path-numerator
error. A collision-uniform length-tail estimate sufficient for the
full raw pointwise law is isolated with its order of limits. That
dynamical estimate is not inferred from fixed-count finite variation.



\section{The physical record and the pointwise theorem}
\label{sec:v74-introduction}
Let $R\in I=[9/20,47/100]$. Under the original section probability
$\nu_R^*=c_*^{-1}\nu|_{Y_R^*}$, with $c_*=91/(10000\pi)$, write
$J_{n,R}=(K_{n,R},N_{n,R},T_{n,R})$ for the actual return record.
The density at the exact displacement $k$ and collision count $m$
is $p_{n,R}(k,m,t)$. Occupation counts collisions $0,\ldots,m-1$;
terminal section membership is tested at collision $m$. Set
\[
 Z_{m,R}(k,t,n)=\frac{(k_1,k_2,t-m\bar\tau_R,n-mc_*)}{\sqrt m},
 \qquad \mathcal C_{m,R,n,k}(M)=\{t:|Z_{m,R}(k,t,n)|\le M\}.
\]
The organizing pointwise assertion is, for every finite $M$,
\begin{equation}\label{eq:v74-original-target}
 \sup_{R,\,1\le n\le m,\,k\in\Z^2}
 \esssup_{t\in\mathcal C_{m,R,n,k}(M)}
 |m^2p_{n,R}(k,m,t)-\mathcal L_{m,R}(k_1,k_2,t,n)|\longrightarrow0.
\end{equation}
Here $\mathcal L_{m,R}$ is the full arithmetic transition reference
of the technical supplement, including every zero class. Denote the
collision limsup of the left side by $\mathcal E_M$. No discrete
coordinate is pooled and no roof set is removed beyond the null sets
inherent in density versions.

For a fixed guard width $\varepsilon>0$, the complete positive
source is
\begin{equation}\label{eq:v74-full-source}
 b^\varepsilon(t)\,dt
  =L_\#\bigl((1-H^\varepsilon)\one_{\Pi_{n,k,m,R}}\nu_R^*\bigr).
\end{equation}
The guard contains all original incidence and clearance factors.
In the earlier source decomposition it is
$b=s_{71,1}+d_{71,1}+b^{\varepsilon,\mathrm{inc}}+e_{70}$,
where $e_{70}$ is first clearance outside the selected disks,
with first incidence separate. Every term remains present.
Define, before any averaging argument,
\[
 A_\delta f(t)=\frac1{2\delta}\int_{t-\delta}^{t+\delta}f(s)\,ds,
 \qquad
 \mathcal H(f)=\limsup_{m\to\infty}\sup_{R,n,k}m^2\|f\|_\infty.
\]
Widths are fixed before this collision limit.

\subsection{The precise inherited inputs}
The technical supplement contains all 163 preceding proof modules,
their appendices and bibliography. Its theorem numbering is distinct
from the present article. The following statement identifies exactly
which interfaces are used; the subsequent arguments do not re-prove
the full spectral, arithmetic or return-clock chain.

\begin{inputstatement}[Source-pinned analytic interfaces]
\label{input:v74-dependencies}
The complete fixed-count height theorem gives
$\|\rho_{m,R}\|_\infty\le CA^m$, uniformly in $R$, and hence bounds
every positive exact-label restriction by $L_m=CA^m/c_*$.
For each fixed $\varepsilon,\delta>0$, the complete-source
fixed-width theorem gives
\begin{equation}\label{eq:v74-average-input}
                    \mathcal H(A_\delta b^\varepsilon)
                                         \le C\varepsilon^{1/16}.
\end{equation}
For sufficiently large fixed $B$ and
$\varepsilon(B)=A_0B^{-1/12}$, the positive scalar remainder satisfies
\begin{equation}\label{eq:v74-positive-input}
 m^2p-\mathcal L=m^2b^{\varepsilon(B)}+e_B,
 \qquad \limsup_m\sup_{R,n,k}\esssup_{\mathcal C(M)}|e_B|
                                              \le C_MB^{-1/192}.
\end{equation}
The path consequence uses additionally the common positive
path-remainder theorem with this same scalar mass, and actual-return
path consequences additionally use the original clock-transfer
hypotheses on positive arithmetic classes.
\end{inputstatement}

The exact sources are respectively the supplement's
Theorem~\ref{thm:v44-complete-density-bound},
\eqref{eq:v72-finite-average-height},
\eqref{eq:v51-positive-error}, and
Theorem~\ref{thm:v52-path-remainder}. The common-kernel version
convention and the scalar/path comparison are recorded there in
Lemma~\ref{lem:v73-common-kernel} and
Theorem~\ref{thm:v73-scalar-path-equality}.
In \eqref{eq:v74-average-input}, the auxiliary band is fixed before
the collision limit and enlarged afterwards. The resulting constant
is independent of the previously fixed width. These are the precise
interfaces used in the main error comparison.

\subsection{The new estimate}
For a nonnegative compactly supported $BV$ density $f$, let $I_u(t)$
be the component of the essential superlevel set $\{f>u\}$ containing
$t$. Define its short-excursion budget by
\[
 S_\delta f(t)=\int_0^{f(t)}(1-|I_u(t)|/\delta)_+\,du.
\]
Section~\ref{sec:v74-level-matching} specifies the measurable versions
and proves
\begin{equation}\label{eq:v74-leading-estimate}
 [f-2A_\delta f]_+\le S_\delta f
 \le\min\{(Df)^+([t-\delta,t]),(Df)^-([t,t+\delta])\}.
\end{equation}
Unlike the last minimum, $S_\delta f$ matches the entrance and exit
on the same height layer and the same component through the roof.
It can vanish when both of those Jordan masses are positive.
A narrow pulse remains charged with its exact excess.
Moreover,
\[
 \int S_\delta f\le\frac\delta8\|Df\|_{\TV},
\]
with sharp constant. This integrated inequality is not used as a
pointwise density bound.

Theorem~\ref{thm:v74-tame-family} supplies the finite current of the
complete original source with an explicit moving-section graph.
Theorem~\ref{thm:v74-endpoint-equivalence} then shows that its
ordered short-excursion budget tends, as the outer band grows, to
exactly $\mathcal E_M$. Theorem~\ref{thm:v74-tail-criterion} proves
a sufficient quantitative length-tail criterion, with every limit
in the required order. The Lorentz-specific uniform decay in that
criterion is still to be established; the original pointwise
assertion \eqref{eq:v74-original-target} is therefore not claimed
as a consequence of the present fixed-count and deterministic
results alone. The present article, its technical supplement and
the frozen provenance are separate parts of one source-preserving
revision, not different mathematical topics.

\input{core/164_explicit_tame_physical_graph}
\input{core/165_level_matched_short_excursions}
\input{core/166_short_excursion_endpoint_budget}
\input{core/167_same_roof_normalization}

\part{Finite physical currents and the pointwise error}
\input{core/160_common_source_versions_and_interfaces}
\input{core/161_full_source_bv_by_threshold_mixtures}
\input{core/162_paired_flux_capacity_bound}
\input{core/163_scalar_path_error_equivalence}

\input{build/v72-body}
\clearpage
\input{build/v72-frontmatter}
\input{build/v102-archive-references}
\end{document}
